% GENERATED by tools/buildDocs.py from tests/tNgf*.m (methods, tags, comments, recorded
% bounds) and tests/last_run.txt (measured values).
% DO NOT EDIT: auditPackage.py compares this file by content and refuses a stale one.
% Regenerate with: python tools/buildDocs.py

\noindent This catalogue is generated from the test classes by \texttt{tools/buildDocs.py}: 514 test methods in 23 classes, 526 recorded bounds, toolbox version 0.19.0. For each method it gives the category its \texttt{TestTags} block assigns, the unit \texttt{auditTestCategories} credits it to (the longest catalogued name the method name begins with), the comment under its function line, and every bound it records through \texttt{verifyBoundAndRecord}: label, bound, units and the reason stated for the bound's size. The measured value beside a bound is the one \texttt{tests/last\_run.txt} holds, from the run: \emph{NonGravForces 0.18.0  run 2026-11-03 18:58 UTC  MATLAB 26.1.0.3346908 (R2026a) Update 5 tests 511  passed 511  failed 0  incomplete 0  duration 1048.6 s  GATE TRUE}; a dash means that run did not record the label (a newer test, or a label built at run time). A method whose label is computed at run time is shown as \emph{(computed label)}. The seven categories and the rules behind them are those of the section on testing.

\subsection{\texttt{tNgfApp}}
\noindent Tests of the application layer (LP-7): app.forceBudget, app.densityInversion.

\begin{description}[leftmargin=1em,style=nextline,itemsep=2pt,font=\normalfont\ttfamily\small]
\item[forceBudgetContract] {\small\emph{contract \textperiodcentered\ \texttt{app.\allowbreak{}forceBudget}}}
\item[forceBudgetReference] {\small\emph{reference \textperiodcentered\ \texttt{app.\allowbreak{}forceBudget}}} --- ORACLE IS Model.budget, because this unit is a reduction over the same evaluation and nothing else. One chunk covering the whole arc must therefore be bitwise identical: sqrt(sum(x.\textasciicircum{}2)/K) is what mean(x.\textasciicircum{}2) computes, in the same order.
\begin{itemize}[leftmargin=1.2em,itemsep=0pt,topsep=1pt]\small
\item forceBudget vs Model.budget, one chunk: measured 0.0000e+00, bound 0 m/s\textasciicircum{}2 --- \emph{the same reduction over the same evaluation: bitwise or the unit computes physics of its own}
\end{itemize}
\item[forceBudgetPrecision] {\small\emph{precision \textperiodcentered\ \texttt{app.\allowbreak{}forceBudget}}} --- Chunking changes the association of the sum of squares and nothing else, so the RMS may move at the rounding level and the maxima may not move at all.
\begin{itemize}[leftmargin=1.2em,itemsep=0pt,topsep=1pt]\small
\item forceBudget RMS, 40 epochs in chunks of 7 vs one chunk: measured 1.4708e-16, bound 1e-14 --- \emph{K squares summed in one pass or in six: association only}
\item forceBudget maxima under chunking: measured 0.0000e+00, bound 0 m/s\textasciicircum{}2 --- \emph{a maximum of maxima is exact whatever the grouping}
\item forceBudget share vs rms\_norm ratio: measured 0.0000e+00, bound 0 --- \emph{one division, quoted back}
\end{itemize}
\item[forceBudgetTransientChunks] {\small\emph{precision \textperiodcentered\ \texttt{app.\allowbreak{}forceBudget}}} --- Transient walls carry their temperatures from chunk to chunk. Until 2026-10-31 the next chunk STARTED at the last temperatures, so the interval across every boundary was held instead of integrated; since then a later chunk is evaluated from one epoch early and drops it. Chunked and unchunked then integrate the same intervals from the same states, and differ by the rounding of the overlap epoch evaluated in another batch.
\begin{itemize}[leftmargin=1.2em,itemsep=0pt,topsep=1pt]\small
\item forceBudget, transient walls, chunks of 7 vs one chunk: measured 0.0000e+00, bound 1e-12 --- \emph{the boundary intervals are integrated as in one call; a held step would show at \textasciitilde{}1e-4 relative in the thermal row}
\item Model.series, transient walls, chunks of 7 vs 8640: measured 0.0000e+00, bound 1e-12 --- \emph{as above, for series' own chunking}
\end{itemize}
\item[forceBudgetVectorization] {\small\emph{vectorization \textperiodcentered\ \texttt{app.\allowbreak{}forceBudget}}} --- Two models in one call must give exactly what two calls give, one each: stacking is bookkeeping and may not touch a number.
\begin{itemize}[leftmargin=1.2em,itemsep=0pt,topsep=1pt]\small
\item forceBudget, two models stacked vs one each: measured 0.0000e+00, bound 0 m/s\textasciicircum{}2 --- \emph{stacking is bookkeeping; a difference here would mean the models see each other}
\end{itemize}
\item[forceBudgetRobustness] {\small\emph{robustness \textperiodcentered\ \texttt{app.\allowbreak{}forceBudget}}}
\begin{itemize}[leftmargin=1.2em,itemsep=0pt,topsep=1pt]\small
\item forceBudget, K = 1: rms equals max: measured 0.0000e+00, bound 0 m/s\textasciicircum{}2 --- \emph{one sample: the root mean square of one number is its absolute value}
\item forceBudget, explicit q vs the attitude provider: measured 0.0000e+00, bound 0 m/s\textasciicircum{}2 --- \emph{the provider is asked for the same epochs, so the quaternions are the same numbers}
\end{itemize}
\item[forceBudgetMetamorphic] {\small\emph{metamorphic \textperiodcentered\ \texttt{app.\allowbreak{}forceBudget}}}
\begin{itemize}[leftmargin=1.2em,itemsep=0pt,topsep=1pt]\small
\item forceBudget, gcrs vs body maximum magnitude ratio: measured 1.0010e+00, bound 1.6 --- \emph{per-axis maxima can fall on different epochs after a rotation, so the norm of the maxima may grow; it cannot grow without bound}
\item forceBudget: cdScale 2 doubles the drag row: measured 0.0000e+00, bound 1e-15 --- \emph{cdScale multiplies the per-panel gas force and nothing else}
\item forceBudget: cdScale leaves the radiation row alone: measured 0.0000e+00, bound 1e-15 --- \emph{the effects are separate reductions over separate per-panel arrays}
\end{itemize}
\item[forceBudgetSpeed] {\small\emph{speed \textperiodcentered\ \texttt{app.\allowbreak{}forceBudget}}} --- Chunking exists to bound the memory, not to be free. What must be true is that it does not cost a MULTIPLE of the unchunked run: the evaluation dominates and the accumulation is four sums of three numbers per chunk.
\begin{itemize}[leftmargin=1.2em,itemsep=0pt,topsep=1pt]\small
\item forceBudget, 4 chunks / 1 chunk over 20 epochs: measured 2.9840e+00, bound 4.0 --- \emph{four Model.accel calls instead of one on the same 20 epochs; the per-call overhead is what this measures, and a ratio near 4 would mean the evaluation does not vectorise over epochs at all}
\end{itemize}
\item[densityInversionContract] {\small\emph{contract \textperiodcentered\ \texttt{app.\allowbreak{}densityInversion}}}
\item[densityInversionReference] {\small\emph{reference \textperiodcentered\ \texttt{app.\allowbreak{}densityInversion}}} --- ROUND TRIP IS THE ORACLE. A model's own acceleration, fed back in, must return the model's own density: the inversion undoes exactly what the forward model did, and anything else would mean the two disagree about what the gas force is.
\begin{itemize}[leftmargin=1.2em,itemsep=0pt,topsep=1pt]\small
\item density round trip: the model's own acceleration returns its own density: measured 4.5427e-16, bound 1e-14 --- \emph{the gas force is exactly linear in rho, so the inversion is a division and the round trip is exact to rounding}
\item density round trip: the recovered scale is one: measured 4.4409e-16, bound 1e-14 --- \emph{same statement, read off the scale column}
\item \emph{(computed label)}: measured ---, bound 1e-13 --- \emph{the measurement's gas part was scaled by a known factor and nothing else was touched}
\item \emph{(computed label)}: measured ---, bound 1e-13 --- \emph{scaling along the drag direction cannot create a component across it}
\end{itemize}
\item[densityInversionPrecision] {\small\emph{precision \textperiodcentered\ \texttt{app.\allowbreak{}densityInversion}}} --- Chunking must not move the answer: every epoch is inverted on its own, so a different grouping is a different loop over the same independent divisions.
\begin{itemize}[leftmargin=1.2em,itemsep=0pt,topsep=1pt]\small
\item densityInversion, chunks of 7 vs one chunk: measured 0.0000e+00, bound 0 kg/m\textasciicircum{}3 --- \emph{each epoch is an independent division; unlike an RMS there is nothing to associate differently}
\item rho equals scale times the model density: measured 0.0000e+00, bound 0 kg/m\textasciicircum{}3 --- \emph{one multiplication, reported twice}
\end{itemize}
\item[densityInversionVectorization] {\small\emph{vectorization \textperiodcentered\ \texttt{app.\allowbreak{}densityInversion}}} --- One epoch at a time must equal the whole arc, epoch by epoch.
\begin{itemize}[leftmargin=1.2em,itemsep=0pt,topsep=1pt]\small
\item densityInversion, one epoch at a time vs the whole arc: measured 0.0000e+00, bound 1e-13 --- \emph{the inversion has no cross-epoch term; what is left is Model.accel associating differently at K = 1 than at K = 10}
\end{itemize}
\item[densityInversionRobustness] {\small\emph{robustness \textperiodcentered\ \texttt{app.\allowbreak{}densityInversion}}}
\begin{itemize}[leftmargin=1.2em,itemsep=0pt,topsep=1pt]\small
\item a measurement against the modelled drag returns a negative scale: measured 2.2204e-16, bound 1e-13 --- \emph{the sign is data, not an error: clipping it at zero would hide the case it reports}
\end{itemize}
\item[densityInversionMetamorphic] {\small\emph{metamorphic \textperiodcentered\ \texttt{app.\allowbreak{}densityInversion}}}
\begin{itemize}[leftmargin=1.2em,itemsep=0pt,topsep=1pt]\small
\item doubling the measured gas part doubles the density: measured 1.1375e-16, bound 1e-13 --- \emph{linearity in rho, seen from the other side}
\item a radiation error along the drag direction moves rho by exactly its size: measured 2.2204e-16, bound 1e-12 --- \emph{nothing separates the two along that direction, which is why residualPerp is reported beside every value}
\item an error across the drag direction leaves rho untouched: measured 1.1375e-16, bound 1e-12 --- \emph{the projection is what the inversion reads, so an orthogonal error can only reach residualPerp}
\end{itemize}
\item[densityInversionSpeed] {\small\emph{speed \textperiodcentered\ \texttt{app.\allowbreak{}densityInversion}}} --- The cost is Model.accel and nothing else; the inversion itself is a division per epoch. So chunking must cost what the extra Model.accel calls cost and no more.
\begin{itemize}[leftmargin=1.2em,itemsep=0pt,topsep=1pt]\small
\item densityInversion, 4 chunks / 1 chunk over 20 epochs: measured 1.6987e+00, bound 4.0 --- \emph{four Model.accel calls instead of one; the divisions are free by comparison}
\end{itemize}
\end{description}

\subsection{\texttt{tNgfAttitude}}
\noindent Tests of N-12 (\#69): io.readIMU1B, geom.attitudeQuality and geom.AttitudeProvider.fromItsg.

\begin{description}[leftmargin=1em,style=nextline,itemsep=2pt,font=\normalfont\ttfamily\small]
\item[readIMU1BContract] {\small\emph{contract \textperiodcentered\ \texttt{io.\allowbreak{}readIMU1B}}}
\item[readIMU1BReference] {\small\emph{reference \textperiodcentered\ \texttt{io.\allowbreak{}readIMU1B}}} --- The raw text, parsed here: record k carries rcvtime\_intg, rcvtime\_frac (microseconds), time\_ref, satellite, gyro\_id and FiltAng (degrees). Epoch = 2000-01-01 12:00 GPS + intg + frac, i.e. GPS day 51544 + floor(...), the angle deg2rad(FiltAng).
\begin{itemize}[leftmargin=1.2em,itemsep=0pt,topsep=1pt]\small
\item readIMU1B first epoch vs 2020-01-01 00:00:00.116294 GPS: measured 2.6375e-08, bound 1e-6 s --- \emph{the tag's resolution is 1 us; CI 40978582699 found it off by more than 1e-9 s (value not printed)}
\end{itemize}
\item[readIMU1BPrecision] {\small\emph{precision \textperiodcentered\ \texttt{io.\allowbreak{}readIMU1B}}} --- The conversion is one multiplication by pi/180: back in degrees the published 15 digits return.
\begin{itemize}[leftmargin=1.2em,itemsep=0pt,topsep=1pt]\small
\item readIMU1B rad2deg(angle) vs the published FiltAng, record 1: measured 0.0000e+00, bound 4 * eps rel --- \emph{one multiplication there and back}
\end{itemize}
\item[readIMU1BVectorization] {\small\emph{vectorization \textperiodcentered\ \texttt{io.\allowbreak{}readIMU1B}}} --- Two days read together are the two days read alone, concatenated, bitwise.
\item[readIMU1BRobustness] {\small\emph{robustness \textperiodcentered\ \texttt{io.\allowbreak{}readIMU1B}}} --- A record not tagged in GPS time is refused, an empty TimeRange is refused, a TimeRange inside the excerpt returns exactly the records it covers.
\item[readIMU1BMetamorphic] {\small\emph{metamorphic \textperiodcentered\ \texttt{io.\allowbreak{}readIMU1B}}} --- Each gyro is sampled at 8 Hz, its records strictly increasing in time, and its accumulated angle turns at no more than the orbital rate projected on its axis (\textbar{}rate\textbar{} \textless{}= \textbar{}omega\textbar{} \textasciitilde{} 1.1e-3 rad/s).
\item[readIMU1BSpeed] {\small\emph{speed \textperiodcentered\ \texttt{io.\allowbreak{}readIMU1B}}} --- The reader adds a column lookup, a conversion and a GPS check to gracefoReading's read: ratio near 1.
\begin{itemize}[leftmargin=1.2em,itemsep=0pt,topsep=1pt]\small
\item readIMU1B / readL1b time ratio: measured 1.0042e+00, bound 5 --- \emph{a lookup and a conversion over the read}
\end{itemize}
\item[attitudeQualityContract] {\small\emph{contract \textperiodcentered\ \texttt{geom.\allowbreak{}attitudeQuality}}}
\item[attitudeQualityReference] {\small\emph{reference \textperiodcentered\ \texttt{geom.\allowbreak{}attitudeQuality}}} --- tools/pyxcheck/attitude\_quality\_check.py 2020-01-01 (standard library, raw text, the definitions): omega over the 101st interval and each gyro's residual RMS after the projection.
\begin{itemize}[leftmargin=1.2em,itemsep=0pt,topsep=1pt]\small
\item attitudeQuality omega(101) vs the stdlib oracle: measured 2.7789e-14, bound 1e-9 --- \emph{the same rotation vector of R\_k' R\_k+1 in two languages; the oracle's convention (R\textasciicircum{}T of the active matrix) is the one the gyros fit}
\item attitudeQuality gyro residual RMS vs the stdlib oracle (3 gyros): measured 3.8585e-10, bound 1e-5 rel --- \emph{both are exact projections; the epochs differ by referenceSystem's rounding of the GPS seconds}
\item attitudeQuality gyro RMS rate vs the stdlib oracle (3 gyros): measured 2.0261e-11, bound 1e-6 rel --- \emph{the accumulated angle interpolated to the same SCA1B epochs and differenced}
\end{itemize}
\item[attitudeQualityPrecision] {\small\emph{precision \textperiodcentered\ \texttt{geom.\allowbreak{}attitudeQuality}}} --- The residual is orthogonal to every column of the design, to rounding: the projection is exact however ill-conditioned the fit (condition numbers in the thousands on a nominal arc).
\begin{itemize}[leftmargin=1.2em,itemsep=0pt,topsep=1pt]\small
\item \emph{(computed label)}: measured ---, bound 1e-12 --- \emph{an orthogonal projection by economy QR}
\end{itemize}
\item[attitudeQualityVectorization] {\small\emph{vectorization \textperiodcentered\ \texttt{geom.\allowbreak{}attitudeQuality}}} --- The gyros are processed one at a time: passing only gyro 2's records gives gyro 2's result bitwise.
\item[attitudeQualityRobustness] {\small\emph{robustness \textperiodcentered\ \texttt{geom.\allowbreak{}attitudeQuality}}} --- A gyro reset (a jump larger than the threshold) splits that gyro's series: no interval spans it, and every other interval keeps its rate.
\item[attitudeQualityMetamorphic] {\small\emph{metamorphic \textperiodcentered\ \texttt{geom.\allowbreak{}attitudeQuality}}} --- 1. Physics: an Earth-pointing satellite turns about body y at the orbital rate, so the SCA1B rates are \textasciitilde{}(0, 2 pi / P, 0), and against the nominal frame the attitude is off by about a degree.
\begin{itemize}[leftmargin=1.2em,itemsep=0pt,topsep=1pt]\small
\item attitudeQuality residual invariant to a gyro bias of 1e-5 rad/s: measured 1.3539e-08, bound 1e-6 --- \emph{the bias lies in the span of the design; rounding of a 1e-5 rad/s rate over 1 s}
\end{itemize}
\item[attitudeQualitySpeed] {\small\emph{speed \textperiodcentered\ \texttt{geom.\allowbreak{}attitudeQuality}}} --- 1800 rotation vectors and three interpolations against the read of the IMU excerpt itself.
\begin{itemize}[leftmargin=1.2em,itemsep=0pt,topsep=1pt]\small
\item attitudeQuality / readIMU1B time ratio (30 min): measured 8.8522e-03, bound 5 --- \emph{vectorised rotations and interpolations, cheaper than parsing the text}
\end{itemize}
\item[attitudeProviderItsgContract] {\small\emph{contract \textperiodcentered\ \texttt{geom.\allowbreak{}AttitudeProvider}}}
\item[attitudeProviderItsgReference] {\small\emph{reference \textperiodcentered\ \texttt{geom.\allowbreak{}AttitudeProvider}}} --- Graz's quaternion is SCA1B's in the other convention: conjugated once, the provider returns SCA1B's attitude at the shared 10 s epochs, the rotation to rounding.
\begin{itemize}[leftmargin=1.2em,itemsep=0pt,topsep=1pt]\small
\item fromItsg vs readSCA1B at the shared epochs, rotation angle: measured 2.1073e-08, bound 1e-7 rad --- \emph{bit-identical quaternions conjugated once on both paths; acos near 1 resolves only sqrt(eps)}
\end{itemize}
\end{description}

\subsection{\texttt{tNgfCalibrate}}
\noindent Tests of the calibration application (N-02c): app.calibrate over timeSeriesPreparation.run.

\begin{description}[leftmargin=1em,style=nextline,itemsep=2pt,font=\normalfont\ttfamily\small]
\item[calibrateContract] {\small\emph{contract \textperiodcentered\ \texttt{app.\allowbreak{}calibrate}}} --- a day after the real excerpt, so that the orbit of 2020-01-01 cannot cover it
\item[calibrateCampaignContract] {\small\emph{contract \textperiodcentered\ \texttt{app.\allowbreak{}calibrate}}} --- N-09: an Arc array is a campaign -- run per day, one joint fit, results per day.
\item[calibrateReference] {\small\emph{reference \textperiodcentered\ \texttt{app.\allowbreak{}calibrate}}} --- The composition is the oracle: run's arc and the raw reference through fitCalibration with run's kernel give the same calibration, bit for bit.
\begin{itemize}[leftmargin=1.2em,itemsep=0pt,topsep=1pt]\small
\item \emph{(computed label)}: measured ---, bound 1 --- \emph{error / computed bound}
\end{itemize}
\item[calibrateRealDay] {\small\emph{reference \textperiodcentered\ \texttt{app.\allowbreak{}calibrate}}} --- Integration: ACT1B of 2020-01-01 against the Model along GNV1B with SCA1B attitude.
\begin{itemize}[leftmargin=1.2em,itemsep=0pt,topsep=1pt]\small
\item \emph{(computed label)}: measured ---, bound 1 + 1e-12 --- \emph{least squares with a constant in the basis cannot do worse than the mean}
\end{itemize}
\item[calibrateCampaignRealDays] {\small\emph{reference \textperiodcentered\ \texttt{app.\allowbreak{}calibrate}}} --- N-09 on real data: GRACE-C 2020-01-01..03 (30 min of each, GFZ ISDC RL04, fetched 2026-10-15), ACT1B against the fast Model along each day's GNV1B with SCA1B attitude. Completion and consistency, not truth: one common scale, finite per-day biases that agree to a few per cent between consecutive days.
\begin{itemize}[leftmargin=1.2em,itemsep=0pt,topsep=1pt]\small
\item calibrate campaign, real x bias day-to-day spread / mean: measured 7.9303e-02, bound 0.2 --- \emph{three consecutive days of one instrument; measured 0.08 when written (-3.38, -3.32, -3.12e-7)}
\end{itemize}
\item[calibrateRegressorsRealDay] {\small\emph{reference \textperiodcentered\ \texttt{app.\allowbreak{}calibrate}}} --- N-11a on real data: 30 min of 2020-01-01 GRACE-C, ACT1B against the fast Model, with the sensor-unit temperature TSU\_Y+ (AHK1B) and the three net magnetorquer currents (MAG1B) as measured regressors, both on the MAG1B grid (the temperature linearly). Completion and finite coefficients. TICUR is recorded, not refused: it reads absolute zero with a 3 mK quantisation wobble, so the constancy check cannot tell it dead (found writing this test).
\begin{itemize}[leftmargin=1.2em,itemsep=0pt,topsep=1pt]\small
\item calibrate real day, TICUR distance from absolute zero: measured 5.4834e-02, bound 0.1 C --- \emph{a dead channel: -273.2 C on every record, varying by its 3 mK quantisation (choose live channels)}
\item calibrate real day, TSU\_Y+ span over the 30 min: measured 6.5012e-03, bound 0.05 C --- \emph{recorded: the thermally controlled sensor unit barely moves (0.022 C over the whole day)}
\end{itemize}
\item[calibrateReferenceInterpolationPrecision] {\small\emph{precision \textperiodcentered\ \texttt{app.\allowbreak{}calibrate}}} --- Part I, N-02c, the owner's question of 2026-10-15: calibrate evaluates the Model at the 1 s orbit epochs and interpolates by pchip onto the 10 Hz raw grid. Measured here against the Model evaluated AT 10 Hz, across a shadow entry -- the hardest case, the radiation dropping in seconds -- through calibrate's own route (diagnostics.rawReference), both then CRN-filtered as the fit sees them. The orbit is circular at 490 km in the plane holding the Sun direction, so it crosses the umbra; states at both rates are the closed form, so no orbit is interpolated.
\begin{itemize}[leftmargin=1.2em,itemsep=0pt,topsep=1pt]\small
\item calibrate reference, pchip vs 10 Hz Model, raw grid, / signal: measured 3.8097e-03, bound 0.05 --- \emph{recorded: unfiltered, the 1 s interpolation is visible at the shadow boundary (measured 3.9e-3)}
\item calibrate reference, pchip vs 10 Hz Model, CRN-filtered, / signal: measured 3.2482e-05, bound 1e-3 --- \emph{100x below the Model's own 10 \% (density) error; measured 3.2e-5, i.e. 9.3e-13 m/s\textasciicircum{}2, against a filtered noise floor ESTIMATED at \textasciitilde{}2e-11 (1e-10/rtHz over the 0.035 Hz band)}
\end{itemize}
\item[calibratePrecision] {\small\emph{precision \textperiodcentered\ \texttt{app.\allowbreak{}calibrate}}} --- Coloured noise, spikes and a gap through run's QC: the diagonal scale within 5 bootstrap sigma.
\begin{itemize}[leftmargin=1.2em,itemsep=0pt,topsep=1pt]\small
\item calibrate scale error in bootstrap sigma: measured 3.0306e+00, bound 5 sigma --- \emph{seeded noise, 50 replicates, run's QC}
\end{itemize}
\item[calibrateRobustness] {\small\emph{robustness \textperiodcentered\ \texttt{app.\allowbreak{}calibrate}}} --- A jump in the bias, declared: the step column recovers it and leaves the scale alone.
\begin{itemize}[leftmargin=1.2em,itemsep=0pt,topsep=1pt]\small
\item calibrate recovers a declared jump: measured 3.5293e-15, bound 1e-6 rel --- \emph{noise-free; the step is one more exact column}
\item calibrate scale beside a jump: measured 0.0000e+00, bound 1e-6 --- \emph{noise-free}
\item calibrate jump epochs vs UTC-\textgreater{}GPS arithmetic: measured 0.0000e+00, bound 2 * eps(max(sec)) * 1e9 + 1 ns --- \emph{oracle in double seconds: eps of \textasciitilde{}6.3e8 s, plus the 1 ns rounding}
\end{itemize}
\item[calibrateMetamorphic] {\small\emph{metamorphic \textperiodcentered\ \texttt{app.\allowbreak{}calibrate}}} --- Two entry points, one code path: the Model reference handed back in as an Arc gives the same calibration bitwise; only the attribution, which needs the effects, differs.
\item[calibrateSpeed] {\small\emph{speed \textperiodcentered\ \texttt{app.\allowbreak{}calibrate}}} --- calibrate's OWN cost: its parts are run (whose change-point QC is quadratic in the worst case) and fitCalibration, each with its own speed claim; what calibrate adds -- the jump conversion, the merge of options, the attribution -- must stay small beside them. A ratio of three timings on the same host (rule T5).
\begin{itemize}[leftmargin=1.2em,itemsep=0pt,topsep=1pt]\small
\item calibrate / (run + fitCalibration): measured 1.1373e+00, bound 1.5 --- \emph{calibrate adds conversions and a merge; an Arc reference has no attribution to fit}
\end{itemize}
\end{description}

\subsection{\texttt{tNgfCalibration}}
\noindent Tests of the calibration truth and its result object (N-02a): app.syntheticPair, app.Calibration.

\begin{description}[leftmargin=1em,style=nextline,itemsep=2pt,font=\normalfont\ttfamily\small]
\item[syntheticPairContract] {\small\emph{contract \textperiodcentered\ \texttt{app.\allowbreak{}syntheticPair}}}
\item[calibrationContract] {\small\emph{contract \textperiodcentered\ \texttt{app.\allowbreak{}Calibration}}}
\item[calibrationTimeVaryingContract] {\small\emph{contract \textperiodcentered\ \texttt{app.\allowbreak{}Calibration}}} --- N-03a: a scale curve -- shapes, the one-authority check of scale against the curve, the refused inversion, and the drift report's fields.
\item[calibrationScaleTemperatureContract] {\small\emph{contract \textperiodcentered\ \texttt{app.\allowbreak{}Calibration}}} --- N-11b: the temperature-dependent scale travels with the calibration; apply needs the temperature, invert refuses, and a scale that depends on time AND temperature is refused.
\item[calibrationRegressorContract] {\small\emph{contract \textperiodcentered\ \texttt{app.\allowbreak{}Calibration}}} --- N-11a: measured regressor terms travel with the calibration and apply refuses without them.
\item[syntheticPairReference] {\small\emph{reference \textperiodcentered\ \texttt{app.\allowbreak{}syntheticPair}}} --- The observation IS the CRN chain on S*model + b(t): filter the model with crnApply, scale it, and add dcGain * b, which is exact for a linear b and a symmetric kernel.
\begin{itemize}[leftmargin=1.2em,itemsep=0pt,topsep=1pt]\small
\item syntheticPair vs crnApply(S m) + dcGain b: measured 1.2176e-21, bound tc.chainBound(truth, ref, 4) m/s\textasciicircum{}2 --- \emph{bound 4 Nf eps sum\textbar{}F\textbar{} max\textbar{}raw\textbar{}, computed}
\end{itemize}
\item[calibrationReference] {\small\emph{reference \textperiodcentered\ \texttt{app.\allowbreak{}Calibration}}} --- Applying the planted calibration to the noise-free observation returns the filtered model, up to the one approximation the help names: the bias is evaluated where the chain filtered it, which leaves exactly (dcGain - 1) S \textbackslash{} b(t).
\begin{itemize}[leftmargin=1.2em,itemsep=0pt,topsep=1pt]\small
\item Calibration.apply on the truth vs the filtered model: measured 1.3235e-21, bound cond(truth.scale) * tc.chainBound(truth, ref, 8) m/s\textasciicircum{}2 --- \emph{bound cond(S) 8 Nf eps sum\textbar{}F\textbar{} max\textbar{}raw\textbar{}}
\item Calibration.basis pspline vs closed form: measured 4.9960e-16, bound 1e3 * eps --- \emph{Cox-de Boor against the per-segment cubic, values \textless{}= 2/3}
\end{itemize}
\item[syntheticPairPrecision] {\small\emph{precision \textperiodcentered\ \texttt{app.\allowbreak{}syntheticPair}}} --- The spike table is the truth: the run with spikes minus the run without (same Seed, so the same noise and gap) is the CRN filter of the 1 s boxcars the table lists.
\begin{itemize}[leftmargin=1.2em,itemsep=0pt,topsep=1pt]\small
\item syntheticPair spikes vs crnApply(boxcars): measured 2.0142e-22, bound tc.chainBound(truth, ref, 4) m/s\textasciicircum{}2 --- \emph{two runs, each one Nf-term dot product}
\end{itemize}
\item[calibrationPrecision] {\small\emph{precision \textperiodcentered\ \texttt{app.\allowbreak{}Calibration}}} --- Both directions are the same algebra: the correction-direction apply of invert() equals the observation-direction apply to cond(S) eps; invert twice is the identity.
\begin{itemize}[leftmargin=1.2em,itemsep=0pt,topsep=1pt]\small
\item Calibration observation vs correction apply: measured 2.5965e-16, bound 16 * c * eps * (1 + max(abs(cal.biasAt(t)), [], 'all') / max(abs(y), [], 'all')) rel --- \emph{16 cond(S) eps (1 + \textbar{}b\textbar{}/\textbar{}y\textbar{})}
\item Calibration invert twice, scale: measured 1.1102e-16, bound 16 * c\textasciicircum{}2 * eps --- \emph{16 cond(S)\textasciicircum{}2 eps}
\item Calibration invert twice, bias: measured 1.3235e-16, bound 16 * c\textasciicircum{}2 * eps rel --- \emph{16 cond(S)\textasciicircum{}2 eps}
\item Calibration pspline partition of unity: measured 3.3307e-16, bound 16 * eps --- \emph{four non-zero terms}
\end{itemize}
\item[calibrationTimeVaryingPrecision] {\small\emph{precision \textperiodcentered\ \texttt{app.\allowbreak{}Calibration}}} --- apply against an independent per-epoch loop of 3x3 solves of y = S(t) a + b(t), with S(t) written out by hand; and the Wald statistic of the linear basis against its closed form, the squared slope over its variance.
\begin{itemize}[leftmargin=1.2em,itemsep=0pt,topsep=1pt]\small
\item Calibration apply with S(t) vs a per-epoch loop of solves: measured 0.0000e+00, bound 100 * eps --- \emph{pages vs a loop; both one LU of the same 3x3, differing only in the order S(t) is formed}
\item Calibration drift statistic vs slope\textasciicircum{}2 / var(slope): measured 2.0464e-16, bound 1e-12 --- \emph{pinv of a 1x1 is its reciprocal}
\end{itemize}
\item[calibrationScaleTemperaturePrecision] {\small\emph{precision \textperiodcentered\ \texttt{app.\allowbreak{}Calibration}}} --- apply against a per-epoch loop of y = S(T(t)) a + b(t) + C (z(t) - m), every series interpolated linearly.
\begin{itemize}[leftmargin=1.2em,itemsep=0pt,topsep=1pt]\small
\item Calibration apply with a temperature-dependent scale vs the observation equation: measured 3.9705e-16, bound 1e-9 --- \emph{one 3x3 solve per epoch on S(T(t)); \textasciitilde{}1e-15 expected, bound as for the regressor terms}
\end{itemize}
\item[calibrationRegressorPrecision] {\small\emph{precision \textperiodcentered\ \texttt{app.\allowbreak{}Calibration}}} --- apply against a per-epoch loop of y = S a + b(t) + C (z(t) - m), z interpolated linearly.
\begin{itemize}[leftmargin=1.2em,itemsep=0pt,topsep=1pt]\small
\item Calibration apply with regressor terms vs the observation equation: measured 5.2940e-16, bound 1e-9 --- \emph{a 3x3 solve after subtracting the regressor term; the 1e-7 signal beside a 2e-7 bias leaves \textasciitilde{}1e-15, bound generous}
\end{itemize}
\item[syntheticPairVectorization] {\small\emph{vectorization \textperiodcentered\ \texttt{app.\allowbreak{}syntheticPair}}} --- The arc length is the batch: a shorter arc is the start of a longer one wherever both have a complete window, to the rounding of the chain (the noise is 1e-30).
\begin{itemize}[leftmargin=1.2em,itemsep=0pt,topsep=1pt]\small
\item syntheticPair half arc vs whole arc: measured 0.0000e+00, bound tc.chainBound(truth, ref1, 4) m/s\textasciicircum{}2 --- \emph{the same sums, the BLAS product associated by batch}
\end{itemize}
\item[calibrationVectorization] {\small\emph{vectorization \textperiodcentered\ \texttt{app.\allowbreak{}Calibration}}} --- An arc applied whole equals the arc applied in two halves (not claimed bitwise: the solve and the bias product are batched, CLAUDE.md pitfall).
\begin{itemize}[leftmargin=1.2em,itemsep=0pt,topsep=1pt]\small
\item Calibration.apply whole vs halves: measured 0.0000e+00, bound 16 * cond(cal.scale) * eps rel --- \emph{16 cond(S) eps}
\item Calibration.biasAt batch vs per epoch: measured 0.0000e+00, bound 8 * eps rel --- \emph{8 eps, four basis terms}
\end{itemize}
\item[syntheticPairRobustness] {\small\emph{robustness \textperiodcentered\ \texttt{app.\allowbreak{}syntheticPair}}} --- NaN exactly where the 140.7 s window leaves the arc or meets the gap, nowhere else.
\item[calibrationRobustness] {\small\emph{robustness \textperiodcentered\ \texttt{app.\allowbreak{}Calibration}}}
\item[syntheticPairMetamorphic] {\small\emph{metamorphic \textperiodcentered\ \texttt{app.\allowbreak{}syntheticPair}}} --- The Seed is the only randomness: the same Seed is bitwise the same arcs, another is not.
\begin{itemize}[leftmargin=1.2em,itemsep=0pt,topsep=1pt]\small
\item syntheticPair bias offset -\textgreater{} dcGain * offset: measured 7.7093e-22, bound tc.chainBound(tr1, r1, 4) m/s\textasciicircum{}2 --- \emph{two runs}
\end{itemize}
\item[calibrationMetamorphic] {\small\emph{metamorphic \textperiodcentered\ \texttt{app.\allowbreak{}Calibration}}}
\item[calibrationTimeVaryingMetamorphic] {\small\emph{metamorphic \textperiodcentered\ \texttt{app.\allowbreak{}Calibration}}} --- A curve with zero slope applies as the constant scale does, to a few eps (pages vs one matrix solve); a curve survives save/load; a schema-1 file loads as a constant.
\begin{itemize}[leftmargin=1.2em,itemsep=0pt,topsep=1pt]\small
\item Calibration flat curve vs constant scale, apply: measured 2.1692e-16, bound 100 * eps --- \emph{pages vs one 3x3 solve}
\end{itemize}
\item[syntheticPairSpeed] {\small\emph{speed \textperiodcentered\ \texttt{app.\allowbreak{}syntheticPair}}} --- synthesis and CRN outputs both grow with the arc: twice the arc, about twice the cost.
\begin{itemize}[leftmargin=1.2em,itemsep=0pt,topsep=1pt]\small
\item syntheticPair, 2x arc / 1x arc: measured 1.9673e+00, bound 4 --- \emph{n log n would be about 2.1; quadratic 4}
\end{itemize}
\item[calibrationSpeed] {\small\emph{speed \textperiodcentered\ \texttt{app.\allowbreak{}Calibration}}}
\begin{itemize}[leftmargin=1.2em,itemsep=0pt,topsep=1pt]\small
\item Calibration.apply, 20000 / 5000 epochs: measured 2.2056e+00, bound 8 --- \emph{linear would be 4; the basis is evaluated column by column}
\end{itemize}
\end{description}

\subsection{\texttt{tNgfCompareCalibrations}}
\noindent Tests of the calibration model comparison (N-03b): app.compareCalibrations.

\begin{description}[leftmargin=1em,style=nextline,itemsep=2pt,font=\normalfont\ttfamily\small]
\item[compareCalibrationsContract] {\small\emph{contract \textperiodcentered\ \texttt{app.\allowbreak{}compareCalibrations}}}
\item[compareCalibrationsReference] {\small\emph{reference \textperiodcentered\ \texttt{app.\allowbreak{}compareCalibrations}}} --- The composition, re-derived by hand for one interior fold: fit app.fitCalibration on the rows outside the test block and its buffers, predict the test block, square and average. Bitwise the fold loss compareCalibrations reports.
\item[compareCalibrationsPrecision] {\small\emph{precision \textperiodcentered\ \texttt{app.\allowbreak{}compareCalibrations}}} --- The planted bias order (linear on x and z, constant on y) against a flexible spline: CV must pick the truth on every axis of every seeded realisation, and AIC is recorded beside it -- it picked the spline on 9 of 9 axes when this test was written, which is the v3 plan's reason for reporting it as IC\_unreliable.
\begin{itemize}[leftmargin=1.2em,itemsep=0pt,topsep=1pt]\small
\item compareCalibrations CV misses of the planted bias order, fraction: measured 0.0000e+00, bound 0 --- \emph{3 seeds x 3 axes, Days 0.25, 5 folds; CV must find the truth every time}
\item compareCalibrations AIC picks the overfit spline, fraction: measured 1.0000e+00, bound 1 --- \emph{recorded, not bounded: the v3 plan's claim is that it exceeds 0.5 (IC\_unreliable)}
\end{itemize}
\item[compareCalibrationsVectorization] {\small\emph{vectorization \textperiodcentered\ \texttt{app.\allowbreak{}compareCalibrations}}} --- Candidates compared together equal candidates compared one at a time, bitwise: the folds and the buffer depend on the observation only, and each candidate's fit on a fold is independent of the others.
\item[compareCalibrationsRobustness] {\small\emph{robustness \textperiodcentered\ \texttt{app.\allowbreak{}compareCalibrations}}} --- The default pair carries spikes and a mid-arc gap: every fold still scores finitely.
\item[compareCalibrationsMetamorphic] {\small\emph{metamorphic \textperiodcentered\ \texttt{app.\allowbreak{}compareCalibrations}}} --- Reordering the candidates permutes the rows, bitwise, and leaves the winners alone; a 'Fit' pair a spec repeats is overridden by the spec.
\item[compareCalibrationsSpeed] {\small\emph{speed \textperiodcentered\ \texttt{app.\allowbreak{}compareCalibrations}}} --- Each candidate costs Folds + 1 fits: two candidates, about twice one.
\begin{itemize}[leftmargin=1.2em,itemsep=0pt,topsep=1pt]\small
\item compareCalibrations, 2 candidates / 1: measured 1.9742e+00, bound 4 --- \emph{linear in candidates would be 2; the shared buffer computation pulls it lower}
\end{itemize}
\end{description}

\subsection{\texttt{tNgfCore}}
\noindent Tests of L0: const, toolboxVersion, fail, provenance, Options, dependencyReport.

\begin{description}[leftmargin=1em,style=nextline,itemsep=2pt,font=\normalfont\ttfamily\small]
\item[toolboxVersionContract] {\small\emph{contract \textperiodcentered\ \texttt{toolboxVersion}}}
\item[constContract] {\small\emph{contract \textperiodcentered\ \texttt{const}}}
\item[failContract] {\small\emph{contract \textperiodcentered\ \texttt{fail}}}
\item[provenanceContract] {\small\emph{contract \textperiodcentered\ \texttt{provenance}}}
\item[optionsContract] {\small\emph{contract \textperiodcentered\ \texttt{Options}}}
\item[dependencyReportContract] {\small\emph{contract \textperiodcentered\ \texttt{dependencyReport}}}
\item[constReference] {\small\emph{reference \textperiodcentered\ \texttt{const}}} --- 2018 exact values and IAU 2015 B3 nominal solar radius.
\item[constPrecision] {\small\emph{precision \textperiodcentered\ \texttt{const}}} --- Solar radiation pressure at 1 au: S0/c. Bound 1e-15: pure arithmetic.
\item[constRobustness] {\small\emph{robustness \textperiodcentered\ \texttt{const}}} --- The constants are data, so robustness is about the data being well formed: nothing non-finite, nothing non-positive that is physically positive, the two species arrays the same length, and no hidden state between calls.
\item[toolboxVersionRobustness] {\small\emph{robustness \textperiodcentered\ \texttt{toolboxVersion}}} --- The numeric fields are the authority and `semantic` is composed from them (rule R1): the robustness claim is that the two can never drift, and that a bad format refuses.
\item[optionsRobustness] {\small\emph{robustness \textperiodcentered\ \texttt{Options}}} --- Empty override list and repeated set() are well defined.
\item[failRobustness] {\small\emph{robustness \textperiodcentered\ \texttt{fail}}} --- Message with a literal percent sign must survive sprintf.
\item[provenanceRobustness] {\small\emph{robustness \textperiodcentered\ \texttt{provenance}}}
\item[dependencyReportRobustness] {\small\emph{robustness \textperiodcentered\ \texttt{dependencyReport}}} --- Must not error whatever is on the path; only warn.
\item[optionsMetamorphic] {\small\emph{metamorphic \textperiodcentered\ \texttt{Options}}} --- Setting a field to its own value is the identity; two orders of independent overrides give equal structs.
\item[provenanceMetamorphic] {\small\emph{metamorphic \textperiodcentered\ \texttt{provenance}}} --- Hash depends on content only, not on field order of the struct.
\item[failMetamorphic] {\small\emph{metamorphic \textperiodcentered\ \texttt{fail}}} --- The identifier is a pure function of (unit, mnemonic).
\end{description}

\subsection{\texttt{tNgfDemos}}
\noindent Every demo in demos/ runs to completion under the gate.

\begin{description}[leftmargin=1em,style=nextline,itemsep=2pt,font=\normalfont\ttfamily\small]
\item[demoRuns] {\small\emph{contract \textperiodcentered\ no unit}} --- Demos print by design -- they are top-level scripts, not library code -- and the output is swallowed so the gate's own report stays readable. evalc also means a demo that starts printing megabytes is a memory problem here rather than a wall of text in the run log.
\end{description}

\subsection{\texttt{tNgfEnv}}
\noindent Tests of the env layer (LP-3): density, eclipseFactor, earthRadiance, radiationSources, and buildMex.

\begin{description}[leftmargin=1em,style=nextline,itemsep=2pt,font=\normalfont\ttfamily\small]
\item[windContract] {\small\emph{contract \textperiodcentered\ \texttt{env.\allowbreak{}wind}}}
\item[densityDtm2020Contract] {\small\emph{contract \textperiodcentered\ \texttt{env.\allowbreak{}density}}}
\item[densityContract] {\small\emph{contract \textperiodcentered\ \texttt{env.\allowbreak{}density}}}
\item[magneticFieldContract] {\small\emph{contract \textperiodcentered\ \texttt{env.\allowbreak{}magneticField}}}
\item[eclipseFactorContract] {\small\emph{contract \textperiodcentered\ \texttt{env.\allowbreak{}eclipseFactor}}} --- at 490 km the penumbra spans 67.97..68.51 deg from the anti-Sun point (b -/+ a); 68.25 is inside it
\item[earthRadianceContract] {\small\emph{contract \textperiodcentered\ \texttt{env.\allowbreak{}earthRadiance}}}
\item[radiationSourcesDe440Reference] {\small\emph{contract \textperiodcentered\ \texttt{env.\allowbreak{}radiationSources}}} --- PRECISE PRESET'S OWN DEFAULT, evaluated at last. Until 2026-09-16 the only test naming rs.de440 asserted that it REFUSES without kernels, so a shipped preset's default sun.source had never produced a number (HANDOVER O-7, O-15).  Oracle: JPL's DE440 through referenceSystem's MICE provider -- an ephemeris this project did not compute. The claim is NOT that the two agree to rounding; they are different ephemerides. It is that referenceSystem's analytic Sun is accurate enough for radiation pressure, which is what justifies rs.analytic being a preset value at all.
\begin{itemize}[leftmargin=1.2em,itemsep=0pt,topsep=1pt]\small
\item DE440 vs the analytic Sun: direction: measured 3.1446e-03, bound 0.05 arcsec --- \emph{measured 0.0031 arcsec 2026-09-16; 1 arcsec of Sun direction moves a panel's cosine by 5e-6}
\item DE440 vs the analytic Sun: distance: measured 2.0083e-09, bound 1e-7 --- \emph{measured 2.0e-9 (0.3 km) 2026-09-16}
\item DE440 vs the analytic Sun: solar radiation pressure: measured 4.0165e-09, bound 1e-7 --- \emph{twice the relative distance difference (1/r\textasciicircum{}2); measured 4.0e-9, far below the 1e-3 uncertainty of the solar constant itself}
\end{itemize}
\item[radiationSourcesContract] {\small\emph{contract \textperiodcentered\ \texttt{env.\allowbreak{}radiationSources}}}
\item[buildMexContract] {\small\emph{contract \textperiodcentered\ \texttt{buildMex}}} --- 'Target' names one gateway and the report is scalar; the default "all" returns one element per target and is asserted at the end of this method.
\item[windReference] {\small\emph{reference \textperiodcentered\ \texttt{env.\allowbreak{}wind}}} --- An independent route to the same number: the nominal rigid rotation omega x r, with omega from referenceSystem's own IERS constant and the position in ITRS. The two agree to the polar motion and length-of-day terms the nominal rate omits.
\begin{itemize}[leftmargin=1.2em,itemsep=0pt,topsep=1pt]\small
\item co-rotation vs a nominal omega x r: measured 1.4123e-06, bound 5e-3 --- \emph{Rdot carries polar motion and the length-of-day term that a nominal rate omits}
\item co-rotation speed vs omega * sqrt(x\textasciicircum{}2 + y\textasciicircum{}2) over the arc: measured 1.5651e-06, bound 5e-3 --- \emph{the residual is the polar motion and the length-of-day term; measured 2e-6 over 600 epochs spanning 256 to 467 m/s}
\end{itemize}
\item[windHwm14Reference] {\small\emph{reference \textperiodcentered\ \texttt{env.\allowbreak{}wind}}} --- Oracle: NRL's own check output for HWM14 (Check/gfortran.txt, read from the package the MEX was built from), printed to 1 mm/s, for the height, latitude, local-time, longitude, day-of-year and ap profiles of checkhwm14.f90, whose fixed inputs are transcribed below from that driver. Through THIS unit: epoch, GCRS position, rotation, geodetic conversion and the rotation of the wind back are all in the loop. Quiet columns use 'Disturbed', false; total columns 'Ap'. Rows the unit cannot pose are left out and named: the two poles (north and east undefined by a position), day 0 (no such date), local times whose UT the driver pushes past 86400 s (a date interface normalises them into the next ...
\begin{itemize}[leftmargin=1.2em,itemsep=0pt,topsep=1pt]\small
\item \emph{(computed label)}: measured ---, bound 6e-4 m/s --- \emph{half a printed unit (5e-4) plus the model's single-precision response to a geodetic round trip; disturbed = total - quiet carries two roundings, halved}
\end{itemize}
\item[densityDtm2020Reference] {\small\emph{reference \textperiodcentered\ \texttt{env.\allowbreak{}density}}} --- CNES's benchmark, every printed number: tz and tinf to 3 decimals, ro to 5 significant digits.
\begin{itemize}[leftmargin=1.2em,itemsep=0pt,topsep=1pt]\small
\item dtm2020 vs CNES benchmark: tz (4 cases): measured 3.7876e-04, bound 1e-3 K --- \emph{printed to 3 decimals (half a unit 5e-4); first run 2026-10-27}
\item dtm2020 vs CNES benchmark: tinf (4 cases): measured 4.7384e-04, bound 1e-3 K --- \emph{printed to 3 decimals (half a unit 5e-4); first run 2026-10-27}
\item dtm2020 vs CNES benchmark: ro, relative (4 cases): measured 1.4794e-05, bound 1e-4 --- \emph{printed to 5 significant digits (half a unit \textless{}= 1.9e-5 relative); first run 2026-10-27}
\end{itemize}
\item[densityReference] {\small\emph{reference \textperiodcentered\ \texttt{env.\allowbreak{}density}}} --- The 17 published cases of DOCUMENTATION \S{}7 (printed with \%E, 7 significant digits) must be reproduced digit for digit in every column: HE O N2 O2 AR RHO H N ANOM TINF TG.
\begin{itemize}[leftmargin=1.2em,itemsep=0pt,topsep=1pt]\small
\item msis00 vs DOCUMENTATION table (17 cases): measured 4.8028e-07, bound 5e-7 rel --- \emph{printed 7 digits; measured 4.8e-7}
\item msis21 vs NRL msis2.1\_test\_ref\_dp (200 cases, 10 columns), in half last printed digits: measured 9.9859e-01, bound 1 --- \emph{the package's own statement: double precision reproduces its reference exactly; rounding of the print only}
\end{itemize}
\item[radiationSourcesCeresReference] {\small\emph{reference \textperiodcentered\ \texttt{env.\allowbreak{}radiationSources}}} --- Oracle: tools/pyxcheck/ceres\_check.py integrates the SAME continuous CERES field (geodetic bilinear in space, linear in time, albedo = SWup/SWdown on the instantaneous insolation, longwave as given, on the 20 km reference sphere) over the visible cap by Gauss-Legendre quadrature to \textless{} 2e-5, written from a specification by an agent that did not open this toolbox (T9). Eight cases: sub-solar, terminator, polar summer and night, over the pole, the longitude seam, a node time, the night side. THE CLAIM IS IDENTITY OF THE MODEL, and the tile sum converges to it at second order in the ring width. So the Richardson extrapolation of [24 48] and [48 96] must meet the oracle to its own quadrature ...
\begin{itemize}[leftmargin=1.2em,itemsep=0pt,topsep=1pt]\small
\item ceres Earth sources vs the oracle, tiling error at [24 48] (precise preset): measured 3.5387e-03, bound 1e-2 --- \emph{second-order tiling error of the cap; measured 3.5e-3 of the total Earth irradiance on 2026-09-15}
\item ceres Earth sources, Richardson-extrapolated tile sum vs the oracle: measured 5.6365e-05, bound 1e-4 --- \emph{the oracle's own quadrature error (1.7e-5) plus the next term of the tiling series; measured 5.6e-5 on 2026-09-15}
\item ceres longwave tiling error ratio [24 48] / [48 96], departure from 4: measured 1.9593e-01, bound 0.5 --- \emph{second-order convergence to the same integral: the ratios measured 3.80 to 4.20 on 2026-09-15}
\end{itemize}
\item[magneticFieldReference] {\small\emph{reference \textperiodcentered\ \texttt{env.\allowbreak{}magneticField}}} --- Oracle 1: tools/pyxcheck/igrf\_check.py, an independent standard-library synthesis written from the IGRF description by an agent that did not open this toolbox (T9), for 12 cases: both poles exactly, the dipole alone, the time rule inside a five-year step and in the secular-variation extrapolation, 1900, and a geostationary radius. Each case takes its coefficients from io.igrfCoefficients at the case's year.
\begin{itemize}[leftmargin=1.2em,itemsep=0pt,topsep=1pt]\small
\item magneticField vs the independent synthesis (12 cases, Cartesian; NED off the poles): measured 8.6537e-16, bound 1e-13 --- \emph{two routes to the same series of at most 104 terms; measured 8.7e-16 on 2026-09-15}
\item GRACE-C MAG1B intensity / IGRF-14 intensity, largest departure from 1: measured 3.9020e-03, bound 0.012 --- \emph{main field only against a raw magnetometer; measured 3.9e-3 on 2026-09-15, bound 3x}
\end{itemize}
\item[eclipseFactorReference] {\small\emph{reference \textperiodcentered\ \texttt{env.\allowbreak{}eclipseFactor}}} --- Oracle: pure-Python numerical disc-overlap integration (eclipse\_check.py); cylindrical bitwise.
\begin{itemize}[leftmargin=1.2em,itemsep=0pt,topsep=1pt]\small
\item eclipse conical vs Python disc integration: measured 5.6066e-14, bound 1e-11 --- \emph{oracle \textasciitilde{}1e-12; atan2 form (acos form measured 6.4e-10)}
\item SOLAARS-CF vs Python restatement of eq. 3.86: measured 2.2204e-16, bound 1e-14 --- \emph{same coefficients, tanh/exp rounding}
\end{itemize}
\item[earthRadianceReference] {\small\emph{reference \textperiodcentered\ \texttt{env.\allowbreak{}earthRadiance}}} --- Oracle: knocke\_coefficients.csv (independent restatement of Table 1) and the formula evaluated here from it.
\begin{itemize}[leftmargin=1.2em,itemsep=0pt,topsep=1pt]\small
\item Knocke albedo tile values vs Table 1 formula: measured 5.5511e-17, bound 1e-15 --- \emph{same arithmetic from an independent coefficient file}
\item Knocke emissivity tile values vs Table 1 formula: measured 0.0000e+00, bound 1e-15
\end{itemize}
\item[radiationSourcesReference] {\small\emph{reference \textperiodcentered\ \texttt{env.\allowbreak{}radiationSources}}} --- Oracle: knocke\_check.py cap integration at 0.25 deg. Errors relative to the longwave flux, per grid. Measured 2026-09-07 (worst over the six zonal cases): [12 24] 2.3e-2, [24 48] 5.4e-3, [48 96] 1.4e-3; first-order convergence in the ring count. Bounds 3x-ish: 4e-2, 1e-2, 3e-3.
\begin{itemize}[leftmargin=1.2em,itemsep=0pt,topsep=1pt]\small
\item \emph{(computed label)}: measured ---, bound bounds.(names(gi)) rel --- \emph{relative to the longwave flux, worst of 6 cases}
\item pointSource albedo vs uniform sphere at 100 R: measured 7.4436e-03, bound 2e-2 rel --- \emph{far-field law, O(R/r)}
\item pointSource longwave vs uniform sphere at 100 R: measured 2.5547e-05, bound 5e-3 rel --- \emph{exact law; oracle rim tiles \textasciitilde{}2e-3}
\item SRP at 1 au vs tsi/c: measured 0.0000e+00, bound 1e-15 rel --- \emph{one division}
\end{itemize}
\item[windPrecision] {\small\emph{precision \textperiodcentered\ \texttt{env.\allowbreak{}wind}}} --- The corotation branch subtracts exactly Rdot*(R'*r) and nothing else.
\begin{itemize}[leftmargin=1.2em,itemsep=0pt,topsep=1pt]\small
\item co-rotation term vs Rdot*(R'*r): measured 0.0000e+00, bound 0 m/s --- \emph{the same two products in the same order: bitwise}
\end{itemize}
\item[windHwm14Precision] {\small\emph{precision \textperiodcentered\ \texttt{env.\allowbreak{}wind}}} --- The model wind is HORIZONTAL at the geodetic point: its GCRS vector, rotated back to the ITRS, has no component along the ellipsoid normal and reproduces the north and east components exactly.
\begin{itemize}[leftmargin=1.2em,itemsep=0pt,topsep=1pt]\small
\item hwm14 wind along the ellipsoid normal: measured 4.2633e-14, bound 1e-12 m/s --- \emph{two rotations of a 100 m/s vector}
\item hwm14 north and east recovered from the GCRS vector: measured 9.9476e-14, bound 1e-12 m/s --- \emph{the same two rotations}
\end{itemize}
\item[densityDtm2020Precision] {\small\emph{precision \textperiodcentered\ \texttt{env.\allowbreak{}density}}} --- forms ro as the plain sum of its six partial densities; env.density converts each with the toolbox's molar masses, so n\_j m\_j summed must return rho to a few ulps.
\begin{itemize}[leftmargin=1.2em,itemsep=0pt,topsep=1pt]\small
\item dtm2020: sum of n\_j m\_j over rho - 1: measured 3.3307e-16, bound 8 * eps --- \emph{six conversions and a sum against the model's own sum: rounding only}
\end{itemize}
\item[densitySpeciesSumPrecision] {\small\emph{precision \textperiodcentered\ \texttt{env.\allowbreak{}density}}} --- kernel.dragPanelForce re-forms the mass density from the species and const.speciesMass, while NRLMSISE-00 forms its own total from ITS molecular masses. Its ACCURACY block quotes the difference, so the difference has to be measured here rather than asserted.
\begin{itemize}[leftmargin=1.2em,itemsep=0pt,topsep=1pt]\small
\item msis00 total vs the species sum over const.speciesMass: measured 5.6922e-04, bound 1e-3 --- \emph{the backend's molecular masses differ from CODATA/IUPAC values; measured 5.7e-4 on 2026-09-07 and quoted by kernel.dragPanelForce}
\item msis21 total vs the species sum over the model's own masses: measured 3.3307e-16, bound 8 * eps --- \emph{eight products summed in another order}
\item msis21 total vs the species sum over const.speciesMass: measured 3.0618e-04, bound 1e-3 --- \emph{the model's molar masses against this toolbox's; the kernel re-forms the density from the latter}
\end{itemize}
\item[densityPrecision] {\small\emph{precision \textperiodcentered\ \texttt{env.\allowbreak{}density}}} --- Exponential model: rho(h0) = rho0 exactly; rho(h0 + H) = rho0/e to 4 eps (one exp, one multiply).
\begin{itemize}[leftmargin=1.2em,itemsep=0pt,topsep=1pt]\small
\item exponential density one scale height: measured 0.0000e+00, bound 4 * eps rel --- \emph{exp and a multiply}
\end{itemize}
\item[magneticFieldPrecision] {\small\emph{precision \textperiodcentered\ \texttt{env.\allowbreak{}magneticField}}} --- Degree 1 IS a dipole: B = mu0/(4 pi) [3 (m . rhat) rhat - m] / r\textasciicircum{}3 with the moment vector m = (4 pi a\textasciicircum{}3 / mu0) [g11; h11; g10], so B = (a/r)\textasciicircum{}3 [3 (d . rhat) rhat - d], d = [g11; h11; g10].
\begin{itemize}[leftmargin=1.2em,itemsep=0pt,topsep=1pt]\small
\item magneticField degree 1 vs the closed-form dipole: measured 1.5230e-16, bound 1e-14 --- \emph{the same three coefficients through two formulas, one of them the exact pole}
\end{itemize}
\item[eclipseFactorPrecision] {\small\emph{precision \textperiodcentered\ \texttt{env.\allowbreak{}eclipseFactor}}} --- Exact limits of the conical model: full light 1, umbra 0, annular 1 - b\textasciicircum{}2/a\textasciicircum{}2 (Sun larger than the body).
\begin{itemize}[leftmargin=1.2em,itemsep=0pt,topsep=1pt]\small
\item eclipse annular limit 1 - b\textasciicircum{}2/a\textasciicircum{}2: measured 0.0000e+00, bound 1e-15 --- \emph{same asin, one division}
\item SOLAARS-CF umbra limit: measured 0.0000e+00, bound 1e-12 --- \emph{three saturated sigmoids}
\item SOLAARS-CF full-light limit: measured 0.0000e+00, bound 1e-12
\end{itemize}
\item[earthRadianceClimatologyPrecision] {\small\emph{precision \textperiodcentered\ \texttt{env.\allowbreak{}earthRadiance}}} --- A field LINEAR in the grid's own latitude is reproduced exactly by bilinear interpolation, so every tile must carry the closed form at the GEOCENTRIC latitude it was read at. The ceres test above forms the GEODETIC latitude for the same reason and the two differ by up to 0.19 deg: if this unit ever converted, this bound would break by four orders of magnitude rather than subtly.
\begin{itemize}[leftmargin=1.2em,itemsep=0pt,topsep=1pt]\small
\item ceresClimatology albedo of a latitude-linear field vs its closed form: measured 5.5511e-17, bound 1e-13 --- \emph{bilinear reproduces a linear field; rounding only}
\item ceresClimatology longwave of a constant field: measured 0.0000e+00, bound 1e-12 W/m\textasciicircum{}2 --- \emph{every node equal}
\item ceresClimatology tile latitude is GEOCENTRIC: measured 1.4211e-14, bound 1e-12 deg --- \emph{the grid's own convention, not converted}
\item ceresClimatology reference height measured from the grid: measured 3.7107e-10, bound 1e-6 m --- \emph{the constant the synthetic radii were built with}
\end{itemize}
\item[earthRadianceCeresPrecision] {\small\emph{precision \textperiodcentered\ \texttt{env.\allowbreak{}earthRadiance}}} --- A field LINEAR in latitude is reproduced exactly by bilinear interpolation, and one linear in time by the time blend, so every tile must carry the closed form at the geodetic latitude it was read at -- which this test forms itself from the tile normal with tan(phi\_d) = tan(phi\_c)/(1 - f)\textasciicircum{}2 (poleward of +-89.5 deg the grid is held constant, and so is the closed form).
\begin{itemize}[leftmargin=1.2em,itemsep=0pt,topsep=1pt]\small
\item ceres albedo of a latitude-linear field vs its closed form: measured 5.5511e-17, bound 1e-13 --- \emph{bilinear and linear blends reproduce a linear field; rounding of single-precision nodes exactly representable}
\item ceres longwave of a constant field: measured 2.8422e-14, bound 1e-12 W/m\textasciicircum{}2 --- \emph{a blend of equal values}
\item ceres tile latitude vs the geodetic formula: measured 1.4211e-14, bound 1e-12 deg --- \emph{the same formula}
\end{itemize}
\item[earthRadiancePrecision] {\small\emph{precision \textperiodcentered\ \texttt{env.\allowbreak{}earthRadiance}}} --- Cap area: sum of the tiles equals 2 pi R\textasciicircum{}2 (1 - cos alpha\_max) (measured 4e-16); tile normals unit.
\begin{itemize}[leftmargin=1.2em,itemsep=0pt,topsep=1pt]\small
\item ERP cap tiles area closure: measured 4.4409e-16, bound 1e-14 rel --- \emph{measured 4e-16}
\item ERP tile normals unit norm: measured 2.2204e-16, bound 1e-14
\end{itemize}
\item[radiationSourcesPrecision] {\small\emph{precision \textperiodcentered\ \texttt{env.\allowbreak{}radiationSources}}} --- Anti-solar geometry: every albedo tile is unlit, so the albedo pressure is exactly zero while longwave is not; the Sun direction in the body frame with the identity attitude is the GCRS unit vector to 1e-15.
\begin{itemize}[leftmargin=1.2em,itemsep=0pt,topsep=1pt]\small
\item Sun direction body vs GCRS for identity q: measured 0.0000e+00, bound 1e-15 --- \emph{identity rotation}
\end{itemize}
\item[windVectorization] {\small\emph{vectorization \textperiodcentered\ \texttt{env.\allowbreak{}wind}}}
\item[windHwm14Vectorization] {\small\emph{vectorization \textperiodcentered\ \texttt{env.\allowbreak{}wind}}} --- Batched equals alone, bitwise: the gateway loops per point and the model keeps no state that changes a value.
\item[densityDtm2020Vectorization] {\small\emph{vectorization \textperiodcentered\ \texttt{env.\allowbreak{}density}}} --- The gateway loops over epochs, so a batch equals the same epochs one at a time, bitwise.
\item[densityVectorization] {\small\emph{vectorization \textperiodcentered\ \texttt{env.\allowbreak{}density}}} --- The 15 daily-ap published cases in one call equal the 15 single calls, bitwise (the MEX loop is per epoch).
\item[magneticFieldVectorization] {\small\emph{vectorization \textperiodcentered\ \texttt{env.\allowbreak{}magneticField}}}
\item[eclipseFactorVectorization] {\small\emph{vectorization \textperiodcentered\ \texttt{env.\allowbreak{}eclipseFactor}}}
\item[earthRadianceClimatologyVectorization] {\small\emph{vectorization \textperiodcentered\ \texttt{env.\allowbreak{}earthRadiance}}}
\item[earthRadianceCeresVectorization] {\small\emph{vectorization \textperiodcentered\ \texttt{env.\allowbreak{}earthRadiance}}}
\item[earthRadianceVectorization] {\small\emph{vectorization \textperiodcentered\ \texttt{env.\allowbreak{}earthRadiance}}}
\item[radiationSourcesVectorization] {\small\emph{vectorization \textperiodcentered\ \texttt{env.\allowbreak{}radiationSources}}}
\item[windRobustness] {\small\emph{robustness \textperiodcentered\ \texttt{env.\allowbreak{}wind}}}
\begin{itemize}[leftmargin=1.2em,itemsep=0pt,topsep=1pt]\small
\item co-rotation on the rotation axis: measured 7.0273e-04, bound 1e-2 m/s --- \emph{omega x r vanishes on the axis; what remains is the polar-motion and nutation rate}
\end{itemize}
\item[windHwm14Robustness] {\small\emph{robustness \textperiodcentered\ \texttt{env.\allowbreak{}wind}}} --- Ground level, the model's top (the quiet wind is held constant above its highest level) and a storm ap all evaluate; the current folder is the caller's afterwards (the gateway loads its tables from its own folder).
\item[densityDtm2020Robustness] {\small\emph{robustness \textperiodcentered\ \texttt{env.\allowbreak{}density}}} --- The edges of the drivers and of the position all evaluate to finite positive densities, and a local time of 24 h is 0 h.
\item[densityRobustness] {\small\emph{robustness \textperiodcentered\ \texttt{env.\allowbreak{}density}}}
\begin{itemize}[leftmargin=1.2em,itemsep=0pt,topsep=1pt]\small
\item msis21 / msis00 at 490 km on 2020-01-01, largest departure of the ratio from 1: measured 1.0398e-01, bound log(2) --- \emph{two empirical models of the same thermosphere: a plausibility bound on the wiring (species order, units), not a model comparison}
\end{itemize}
\item[magneticFieldRobustness] {\small\emph{robustness \textperiodcentered\ \texttt{env.\allowbreak{}magneticField}}} --- The geographic poles are ordinary points: the field there is the limit from a microradian away.
\begin{itemize}[leftmargin=1.2em,itemsep=0pt,topsep=1pt]\small
\item \emph{(computed label)}: measured ---, bound 1e-5 --- \emph{continuity: the horizontal gradient over 7 m of arc is below 1e-6 of the field}
\end{itemize}
\item[eclipseFactorRobustness] {\small\emph{robustness \textperiodcentered\ \texttt{env.\allowbreak{}eclipseFactor}}}
\item[earthRadianceClimatologyRobustness] {\small\emph{robustness \textperiodcentered\ \texttt{env.\allowbreak{}earthRadiance}}}
\item[earthRadianceCeresRobustness] {\small\emph{robustness \textperiodcentered\ \texttt{env.\allowbreak{}earthRadiance}}}
\item[earthRadianceRobustness] {\small\emph{robustness \textperiodcentered\ \texttt{env.\allowbreak{}earthRadiance}}} --- Over the pole the tangent basis needs its fallback; a single tile; just above the surface.
\item[radiationSourcesRobustness] {\small\emph{robustness \textperiodcentered\ \texttt{env.\allowbreak{}radiationSources}}}
\item[buildMexRobustness] {\small\emph{robustness \textperiodcentered\ \texttt{buildMex}}} --- A tampered digest makes the state STALE by content; restoring it makes it current again.
\item[windMetamorphic] {\small\emph{metamorphic \textperiodcentered\ \texttt{env.\allowbreak{}wind}}}
\begin{itemize}[leftmargin=1.2em,itemsep=0pt,topsep=1pt]\small
\item co-rotation doubles with the position: measured 0.0000e+00, bound 1e-15 --- \emph{an exact power of two through two matrix products}
\item co-rotation removed and added back: measured 1.2857e-16, bound 1e-15 --- \emph{two roundings of a 7.6 km/s velocity; measured 1.3e-16}
\end{itemize}
\item[windHwm14Metamorphic] {\small\emph{metamorphic \textperiodcentered\ \texttt{env.\allowbreak{}wind}}} --- The model wind is additive: hwm14 minus corotation is exactly the model wind; quiet plus DWM07 is the total (to the model's single precision); and at ap = 0 the disturbance is not zero but DWM07's quiet level.
\begin{itemize}[leftmargin=1.2em,itemsep=0pt,topsep=1pt]\small
\item hwm14 minus corotation is the model wind: measured 4.1211e-13, bound 1e-9 m/s --- \emph{(v - c) - (v - c - w) = w, three roundings of a 7.6 km/s velocity}
\end{itemize}
\item[densityDtm2020Metamorphic] {\small\emph{metamorphic \textperiodcentered\ \texttt{env.\allowbreak{}density}}} --- More solar flux, more density; higher, less. And beside NRLMSISE-00 on the reference day at GRACE-FO's altitude -- the number O-19 asks about: CNES fitted DTM2020 to the TU Delft CHAMP, GOCE and Swarm densities derived with high-fidelity satellite models.
\begin{itemize}[leftmargin=1.2em,itemsep=0pt,topsep=1pt]\small
\item dtm2020 / msis00 at 490 km, 2020-01-01: \textbar{}log\textbar{} of the median ratio: measured 7.3723e-01, bound log(3) --- \emph{first run 2026-10-27: recorded; the solar-minimum density level O-19 asks about}
\item dtm2020 / msis21 at 490 km, 2020-01-01: \textbar{}log\textbar{} of the median ratio: measured 6.5067e-01, bound log(3) --- \emph{first run 2026-10-27: recorded}
\end{itemize}
\item[densityMetamorphic] {\small\emph{metamorphic \textperiodcentered\ \texttt{env.\allowbreak{}density}}} --- gtd7d - gtd7: the mass density gains exactly 16 * 1.66e-24 g times the anomalous-oxygen number density (the reference's own arithmetic, nrlmsise-00.c gtd7d), in SI 16 * 1.66e-24 * 1e-3 * 1e6 kg per m\textasciicircum{}-3.
\begin{itemize}[leftmargin=1.2em,itemsep=0pt,topsep=1pt]\small
\item gtd7d - gtd7 vs 16 amu x anomalous O: measured 2.0195e-16, bound 1e-12 rel --- \emph{two rounding operations, per case}
\item msis21 rho - rho(no anomalous O) vs n(O*) m\_O: measured 8.5441e-17, bound 1e-14 rel --- \emph{a nine-term and an eight-term dot product; a few eps of the total}
\item exponential scale-height ratio: measured 1.6337e-16, bound 1e-14 rel
\end{itemize}
\item[magneticFieldMetamorphic] {\small\emph{metamorphic \textperiodcentered\ \texttt{env.\allowbreak{}magneticField}}} --- A dipole is even in position and falls as r\textasciicircum{}-3; the rotation option is the rotation of the Earth-fixed field.
\begin{itemize}[leftmargin=1.2em,itemsep=0pt,topsep=1pt]\small
\item dipole field is even in position: measured 1.5994e-16, bound 1e-15 --- \emph{cos and sin of phi + pi and of pi - theta}
\item dipole field falls as r\textasciicircum{}-3: measured 0.0000e+00, bound 1e-15 --- \emph{an exact power of two}
\item rotation preserves the intensity: measured 1.3419e-16, bound 1e-15 --- \emph{an orthogonal matrix}
\end{itemize}
\item[eclipseFactorMetamorphic] {\small\emph{metamorphic \textperiodcentered\ \texttt{env.\allowbreak{}eclipseFactor}}} --- Mirror symmetry across the Sun line (bitwise), rotation of the whole geometry (1e-15), monotone through the penumbra along the oracle sweep.
\begin{itemize}[leftmargin=1.2em,itemsep=0pt,topsep=1pt]\small
\item eclipse rotation invariance: measured 2.5424e-14, bound 1e-12 --- \emph{HANDOVER accuracy figure; measured see last\_run}
\item SOLAARS-CF tail 100 km inside the limb: measured 1.3147e-06, bound 1e-4
\item SOLAARS-CF tail 100 km outside the limb: measured 8.3711e-07, bound 1e-4
\item SOLAARS-CF derivative jump (C1 proxy, 100 m steps): measured 4.8176e-03, bound 1e-2 --- \emph{smooth tanh sum: second difference small against the slope}
\item SOLAARS-CF oblateness scaling relation: measured 4.3854e-14, bound 1e-9 --- \emph{rE/sO*sO rounding through steep sigmoids (slope \textasciitilde{}1e-4 per m)}
\item SOLAARS-CF vs conical half-shadow radius shift: measured 2.2500e+04, bound 5e4 m --- \emph{refraction and oblateness; measured see last\_run}
\item SOLAARS-CF limb height above b (normal) vs above a (ecliptic): measured 1.5314e+01, bound 500 m --- \emph{same atmosphere both ways; 100 m grid plus the height scaled by b/a (\textasciitilde{}0.1 km); the inverted mapping gave 42.5 km}
\end{itemize}
\item[earthRadianceClimatologyMetamorphic] {\small\emph{metamorphic \textperiodcentered\ \texttt{env.\allowbreak{}earthRadiance}}} --- CONSTANT IN TIME is the property that separates this backend from ceres, and it is asserted as an identity rather than described in a comment: two epochs inside one month must give bitwise equal tiles, and the step must land exactly at the node.
\begin{itemize}[leftmargin=1.2em,itemsep=0pt,topsep=1pt]\small
\item ceresClimatology steps by exactly one month's difference at the node: measured 2.7756e-17, bound 4 * eps --- \emph{0.1 per month by construction; a linear blend would give a fraction of it}
\item ceresClimatology of a constant field is that constant: measured 0.0000e+00, bound 4 * eps --- \emph{bilinear weights sum to one}
\end{itemize}
\item[earthRadianceCeresMetamorphic] {\small\emph{metamorphic \textperiodcentered\ \texttt{env.\allowbreak{}earthRadiance}}} --- Turning the Earth by exactly one cell and shifting the fields by one column back leaves every tile value: the ITRS longitudes all move by 1 deg and read the same numbers (rounding of atan2 and of the weights).
\begin{itemize}[leftmargin=1.2em,itemsep=0pt,topsep=1pt]\small
\item ceres: Earth turned by one cell, fields shifted by one column: measured 9.5746e-13, bound 1e-9 --- \emph{a longitude moved by 1 deg through atan2 and back into a weight: 1e-13 in the weight, times a field gradient of at most a few W/m\textasciicircum{}2 per cell}
\item ceres albedo is invariant to a common scale of SWup and SWdown: measured 0.0000e+00, bound 1e-14 --- \emph{an exact power of two in numerator and denominator}
\end{itemize}
\item[earthRadianceMetamorphic] {\small\emph{metamorphic \textperiodcentered\ \texttt{env.\allowbreak{}earthRadiance}}} --- One period later the seasonal term repeats (1e-14); mirrored latitudes: the odd term cancels in the sum.
\begin{itemize}[leftmargin=1.2em,itemsep=0pt,topsep=1pt]\small
\item Knocke annual periodicity: measured 0.0000e+00, bound 1e-14 --- \emph{cos of 2 pi}
\item Knocke odd term cancels over mirrored caps: measured 2.2204e-16, bound 1e-14
\end{itemize}
\item[radiationSourcesMetamorphic] {\small\emph{metamorphic \textperiodcentered\ \texttt{env.\allowbreak{}radiationSources}}} --- Rotating the attitude rotates every direction and leaves every pressure (bitwise); scaling tsi scales every pressure by the same factor.
\begin{itemize}[leftmargin=1.2em,itemsep=0pt,topsep=1pt]\small
\item source directions rotate with the attitude: measured 5.5511e-16, bound 1e-14 --- \emph{two rotations}
\item pressures scale with tsi: measured 4.4409e-16, bound 1e-14 --- \emph{one division}
\item ERP convergence ratio (err 12x24 / err 24x48): measured 4.1747e+00, bound 6 --- \emph{first order expects \textasciitilde{}4; measured 4.1}
\end{itemize}
\item[windSpeed] {\small\emph{speed \textperiodcentered\ \texttt{env.\allowbreak{}wind}}} --- Its whole cost is the frame rotation it delegates, plus two page products.
\begin{itemize}[leftmargin=1.2em,itemsep=0pt,topsep=1pt]\small
\item wind vs the rotationItrsToGcrs it calls: measured 2.1482e+00, bound 4.5 --- \emph{budget 3x the expected 1.4: one rotation call plus two page products; 5 repeats, median}
\end{itemize}
\item[windHwm14Speed] {\small\emph{speed \textperiodcentered\ \texttt{env.\allowbreak{}wind}}} --- over the co-rotation for the 600-epoch arc: the geodetic conversion and 600 model calls on top of the same rotation. Measured see last\_run; budget 3x the 2026-09-15 figure.
\begin{itemize}[leftmargin=1.2em,itemsep=0pt,topsep=1pt]\small
\item wind hwm14 / corotation time ratio, 600 epochs: measured 1.4569e+00, bound 6 --- \emph{measured 1.76 on 2026-09-15; budget 3x}
\end{itemize}
\item[densitySpeed] {\small\emph{speed \textperiodcentered\ \texttt{env.\allowbreak{}density}}} --- Linear scaling of the MEX loop: 86400 epochs cost at most 2x ten times 8640 epochs (per-call overhead amortised).
\begin{itemize}[leftmargin=1.2em,itemsep=0pt,topsep=1pt]\small
\item msis00 86400 / 8640 epochs time ratio: measured 9.6786e-01, bound 2 --- \emph{linear scaling; measured see last\_run}
\item msis21 / msis00 time ratio, 8640 epochs: measured 4.8943e+00, bound 15 --- \emph{measured 4.8 on 2026-09-15; budget 3x}
\end{itemize}
\item[magneticFieldSpeed] {\small\emph{speed \textperiodcentered\ \texttt{env.\allowbreak{}magneticField}}} --- Degree 13 over degree 1 for 86400 positions: 104 terms against 3 plus the shared Legendre setup. Measured see last\_run; budget 3x the 2026-09-15 figure.
\begin{itemize}[leftmargin=1.2em,itemsep=0pt,topsep=1pt]\small
\item magneticField degree 13 / degree 1 time ratio, 86400 positions: measured 5.5072e+01, bound 150 --- \emph{measured 50.8 on 2026-09-15; budget 3x}
\end{itemize}
\item[eclipseFactorSpeed] {\small\emph{speed \textperiodcentered\ \texttt{env.\allowbreak{}eclipseFactor}}} --- Conical over cylindrical for 86400 epochs: two asin, an acos and the lens area; measured see last\_run. Bound 8.
\begin{itemize}[leftmargin=1.2em,itemsep=0pt,topsep=1pt]\small
\item eclipse conical / cylindrical time ratio: measured 1.0747e+00, bound 8
\end{itemize}
\item[earthRadianceClimatologySpeed] {\small\emph{speed \textperiodcentered\ \texttt{env.\allowbreak{}earthRadiance}}} --- Against knocke on the same cap and epochs: the rotation, a latitude SEARCH per tile and 8 gathers. The search is what makes this dearer than ceres's arithmetic indexing.
\begin{itemize}[leftmargin=1.2em,itemsep=0pt,topsep=1pt]\small
\item earthRadiance ceresClimatology / knocke time ratio [12 24], 600 epochs: measured 3.2838e+00, bound 30 --- \emph{see last\_run; budget set above the measured figure, ratio not absolute}
\end{itemize}
\item[earthRadianceCeresSpeed] {\small\emph{speed \textperiodcentered\ \texttt{env.\allowbreak{}earthRadiance}}} --- The CERES tiles over Knocke's on the same grid and 600 epochs: the rotation, the geodetic latitude and 8 x 3 gathers per tile. Measured see last\_run; budget 3x the 2026-09-15 figure.
\begin{itemize}[leftmargin=1.2em,itemsep=0pt,topsep=1pt]\small
\item earthRadiance ceres / knocke time ratio [12 24], 600 epochs: measured 4.9874e+00, bound 15 --- \emph{measured 5.0 on 2026-09-15; budget 3x}
\end{itemize}
\item[earthRadianceSpeed] {\small\emph{speed \textperiodcentered\ \texttt{env.\allowbreak{}earthRadiance}}} --- Tile count scales the cost: [24 48] over [12 24] for 600 epochs, four times the tiles, bound 8.
\begin{itemize}[leftmargin=1.2em,itemsep=0pt,topsep=1pt]\small
\item earthRadiance [24 48] / [12 24] time ratio: measured 2.4495e+00, bound 8 --- \emph{4x the tiles}
\end{itemize}
\item[radiationSourcesSpeed] {\small\emph{speed \textperiodcentered\ \texttt{env.\allowbreak{}radiationSources}}} --- The full source list over the bare tiling for the 600-epoch sample: pressures, visibility, rotation. Bound 6.
\begin{itemize}[leftmargin=1.2em,itemsep=0pt,topsep=1pt]\small
\item radiationSources / earthRadiance time ratio (600 epochs): measured 3.6628e+00, bound 6
\end{itemize}
\item[densityScaleContract] {\small\emph{contract \textperiodcentered\ \texttt{env.\allowbreak{}densityScale}}}
\item[densityScaleReference] {\small\emph{reference \textperiodcentered\ \texttt{env.\allowbreak{}densityScale}}} --- Oracle: an INDEPENDENT density model. CNES's DTM2020 (fitted to accelerometer densities, among them Swarm A's) against NRLMSISE-00 at Swarm A's own positions must show the same ratio ESA's POD densities show, to the 10-30 \% both are uncertain by; the month gave 0.53 (DTM2020) against 0.49 (Swarm).
\begin{itemize}[leftmargin=1.2em,itemsep=0pt,topsep=1pt]\small
\item Swarm A / msis00 vs DTM2020 / msis00 on 2020-01-01, \textbar{}log\textbar{}: measured 2.0377e-02, bound log(1.3) --- \emph{two independent sources of the same ratio}
\end{itemize}
\item[densityScalePrecision] {\small\emph{precision \textperiodcentered\ \texttt{env.\allowbreak{}densityScale}}} --- At the bracket's own heights the interpolation returns the bracket's factors: an identity up to exp(log(x)) rounding, a few ulp.
\begin{itemize}[leftmargin=1.2em,itemsep=0pt,topsep=1pt]\small
\item densityScale.at at the bracket heights / bracket factors - 1: measured 2.2204e-16, bound 8 * eps --- \emph{exp(log(x)) round trip}
\end{itemize}
\item[densityScaleVectorization] {\small\emph{vectorization \textperiodcentered\ \texttt{env.\allowbreak{}densityScale}}} --- S.at is element-wise: 300 epochs over three days at once equal 300 single calls, bitwise.
\item[densityScaleRobustness] {\small\emph{robustness \textperiodcentered\ \texttt{env.\allowbreak{}densityScale}}} --- A day with too few records is refused, not filled; one lower satellite suffices, and is then the bracket.
\item[densityScaleMetamorphic] {\small\emph{metamorphic \textperiodcentered\ \texttt{env.\allowbreak{}densityScale}}} --- log f is linear in height: the midpoint of the bracket gets the geometric mean of its ends; the order of the files changes nothing; a denser stride moves the daily factor by less than the product's own noise.
\begin{itemize}[leftmargin=1.2em,itemsep=0pt,topsep=1pt]\small
\item densityScale Stride 2 vs 10, max \textbar{}log ratio\textbar{}: measured 2.4391e-03, bound 0.05 --- \emph{5-min sampling of a \textasciitilde{}20-min product}
\end{itemize}
\item[densityScaleSpeed] {\small\emph{speed \textperiodcentered\ \texttt{env.\allowbreak{}densityScale}}} --- Three satellites x three days: three reads and 2592 msis00 points. Against the same number of msis00 points alone; the reads and the bookkeeping are what the ratio measures. Bound 20, first run.
\begin{itemize}[leftmargin=1.2em,itemsep=0pt,topsep=1pt]\small
\item densityScale / msis00 on as many points, time ratio: measured 7.3899e+00, bound 20 --- \emph{first run}
\end{itemize}
\item[densityScaleFallbackContract] {\small\emph{contract \textperiodcentered\ \texttt{env.\allowbreak{}densityScale}}}
\item[densityScaleFallbackReference] {\small\emph{reference \textperiodcentered\ \texttt{env.\allowbreak{}densityScale}}} --- Oracle: the two scales it combines. On Swarm's days it IS the plain Swarm scale, on every other day it IS the driver scale, bitwise -- epochs of both kinds interleaved in one call.
\item[densityScaleFallbackVectorization] {\small\emph{vectorization \textperiodcentered\ \texttt{env.\allowbreak{}densityScale}}} --- 300 epochs over six days, three with Swarm and three without, at once equal 300 single calls, bitwise.
\item[densityScaleFallbackRobustness] {\small\emph{robustness \textperiodcentered\ \texttt{env.\allowbreak{}densityScale}}} --- A Swarm day with too few records hands the whole day to the fallback, never half of it; a day neither covers (2019-10-01: F30's 81 days before it are not in the excerpt) is refused by the fallback.
\item[densityDriverScaleContract] {\small\emph{contract \textperiodcentered\ \texttt{env.\allowbreak{}densityDriverScale}}}
\item[densityDriverScaleReference] {\small\emph{reference \textperiodcentered\ \texttt{env.\allowbreak{}densityDriverScale}}} --- Oracle: tools/densityDriverFit.m's own drivers and prediction for 2020-01-01 on the runner's full pool (run 40676249479, F30 coefficient record). Independent code -- the tool builds Fbar with a trailing movmean over the file's rows, the unit with a calendar window -- over the same bytes.
\begin{itemize}[leftmargin=1.2em,itemsep=0pt,topsep=1pt]\small
\item densityDriverScale drivers vs densityDriverFit's on 2020-01-01, max rel: measured 0.0000e+00, bound 1e-12 --- \emph{an 81-term mean summed in another order}
\item densityDriverScale log f vs densityDriverFit's on 2020-01-01: measured 0.0000e+00, bound 1e-12 --- \emph{transcribed to 17 digits}
\end{itemize}
\item[densityDriverScalePrecision] {\small\emph{precision \textperiodcentered\ \texttt{env.\allowbreak{}densityDriverScale}}} --- 'Offset' adds to log f and nothing else: an identity up to exp/log rounding of a number near 0.3.
\begin{itemize}[leftmargin=1.2em,itemsep=0pt,topsep=1pt]\small
\item densityDriverScale Offset -0.041: log ratio + 0.041: measured 7.6328e-17, bound 1e-14 --- \emph{exp then log of numbers below 2}
\end{itemize}
\item[densityDriverScaleVectorization] {\small\emph{vectorization \textperiodcentered\ \texttt{env.\allowbreak{}densityDriverScale}}} --- S.at over 300 epochs on three days at once equals 300 single calls, bitwise.
\item[densityDriverScaleRobustness] {\small\emph{robustness \textperiodcentered\ \texttt{env.\allowbreak{}densityDriverScale}}} --- A missing F30 day inside the 81-day window or a missing Ap is refused, not averaged over; an epoch in GPS time is put on its UTC day.
\item[densityDriverScaleMetamorphic] {\small\emph{metamorphic \textperiodcentered\ \texttt{env.\allowbreak{}densityDriverScale}}} --- The factor is the same at every height of a day, and raising the day's Ap by 10 multiplies it by exp(10 c\_ap) exactly; the day before's Ap does not enter.
\item[densityDriverScaleSpeed] {\small\emph{speed \textperiodcentered\ \texttt{env.\allowbreak{}densityDriverScale}}} --- A day of 10 s epochs: one driver evaluation per UTC day, then one exp per epoch. Against msis00 on as many points, which the factor multiplies anyway. Bound 1, first run.
\begin{itemize}[leftmargin=1.2em,itemsep=0pt,topsep=1pt]\small
\item densityDriverScale.at / msis00 on 8640 points, time ratio: measured 4.0498e-02, bound 1 --- \emph{first run}
\end{itemize}
\end{description}

\subsection{\texttt{tNgfFetchCatalog}}
\noindent Tests of NonGravForces.fetchCatalog and of setupNonGravForces' 'Items' against it.

\begin{description}[leftmargin=1em,style=nextline,itemsep=2pt,font=\normalfont\ttfamily\small]
\item[fetchCatalogContract] {\small\emph{contract \textperiodcentered\ \texttt{fetchCatalog}}}
\item[fetchCatalogReference] {\small\emph{reference \textperiodcentered\ \texttt{fetchCatalog}}} --- The spans against an authority not written here: the HAPI headers of the committed TOLEOS and Swarm excerpts (startDate, stopDate), as the server returned them.
\item[fetchCatalogPrecision] {\small\emph{precision \textperiodcentered\ \texttt{fetchCatalog}}} --- Counts derived by hand for one selection; each line says how.
\item[fetchCatalogRobustness] {\small\emph{robustness \textperiodcentered\ \texttt{fetchCatalog}}} --- no selection: the dated sources count nothing, the undated keep their fixed files
\item[fetchCatalogMetamorphic] {\small\emph{metamorphic \textperiodcentered\ \texttt{fetchCatalog}}} --- (1) a range is the list of its days, and splitting it changes nothing
\end{description}

\subsection{\texttt{tNgfFigures}}
\noindent Tests of the calibration figures (N-05): app.plotCalibration, app.plotAttribution.

\begin{description}[leftmargin=1em,style=nextline,itemsep=2pt,font=\normalfont\ttfamily\small]
\item[plotCalibrationContract] {\small\emph{contract \textperiodcentered\ \texttt{app.\allowbreak{}plotCalibration}}}
\item[plotAttributionContract] {\small\emph{contract \textperiodcentered\ \texttt{app.\allowbreak{}plotAttribution}}}
\item[plotCalibrationPrecision] {\small\emph{precision \textperiodcentered\ \texttt{app.\allowbreak{}plotCalibration}}} --- The figure draws the calibration and nothing else: bitwise its methods' values.
\item[plotAttributionPrecision] {\small\emph{precision \textperiodcentered\ \texttt{app.\allowbreak{}plotAttribution}}} --- The per-revolution points are T's, and each fitted line is regressDrivers' coefficients.
\item[plotCalibrationRobustness] {\small\emph{robustness \textperiodcentered\ \texttt{app.\allowbreak{}plotCalibration}}} --- A constant calibration with no covariance (a hand-built one) draws flat lines and no bands; a residual with a gap still whitens, run by run.
\item[plotAttributionRobustness] {\small\emph{robustness \textperiodcentered\ \texttt{app.\allowbreak{}plotAttribution}}} --- A gated revolution is drawn as a gap (NaN), never as a point; an unresolved parameter (no regression) still gets its time panel and banner.
\item[plotCalibrationSpeed] {\small\emph{speed \textperiodcentered\ \texttt{app.\allowbreak{}plotCalibration}}} --- Curves of twice the epochs: about twice the drawing, the fixed figure cost pulling it lower.
\begin{itemize}[leftmargin=1.2em,itemsep=0pt,topsep=1pt]\small
\item plotCalibration, 2x epochs / 1x: measured 1.3690e+00, bound 4 --- \emph{linear would be 2}
\end{itemize}
\item[plotAttributionSpeed] {\small\emph{speed \textperiodcentered\ \texttt{app.\allowbreak{}plotAttribution}}} --- One driver panel against three: the panels dominate, so about three times, never quadratic.
\begin{itemize}[leftmargin=1.2em,itemsep=0pt,topsep=1pt]\small
\item plotAttribution, 3 driver panels / 1: measured 1.1982e+00, bound 6 --- \emph{panels are drawn once each: about 3 at most}
\end{itemize}
\end{description}

\subsection{\texttt{tNgfFitCalibration}}
\noindent Tests of the calibration estimator (N-02b): app.fitCalibration.

\begin{description}[leftmargin=1em,style=nextline,itemsep=2pt,font=\normalfont\ttfamily\small]
\item[fitCalibrationContract] {\small\emph{contract \textperiodcentered\ \texttt{app.\allowbreak{}fitCalibration}}}
\item[fitCalibrationTimeVaryingContract] {\small\emph{contract \textperiodcentered\ \texttt{app.\allowbreak{}fitCalibration}}} --- N-03a: 'ScaleBasis' makes every estimated entry a curve the Calibration carries.
\item[fitCalibrationReference] {\small\emph{reference \textperiodcentered\ \texttt{app.\allowbreak{}fitCalibration}}} --- Noise-free, a full S with off-diagonals and a drifting bias: the truth to rounding.
\begin{itemize}[leftmargin=1.2em,itemsep=0pt,topsep=1pt]\small
\item \emph{(computed label)}: measured ---, bound 1 --- \emph{error / computed bound}
\end{itemize}
\item[fitCalibrationTimeVaryingReference] {\small\emph{reference \textperiodcentered\ \texttt{app.\allowbreak{}fitCalibration}}} --- Noise-free with a planted linear scale drift: S(t) a is exactly in the span of the columns H([1, t] .* a), so the curve returns to the same computed rounding bound.
\item[fitCalibrationPrecision] {\small\emph{precision \textperiodcentered\ \texttt{app.\allowbreak{}fitCalibration}}} --- Coloured noise, spikes and a gap: every diagonal scale within 5 bootstrap sigma of the truth, and the bootstrap sigma no smaller than half the formal one (the filtered residual is correlated, so it should inflate, never collapse).
\begin{itemize}[leftmargin=1.2em,itemsep=0pt,topsep=1pt]\small
\item fitCalibration scale error in bootstrap sigma: measured 3.1104e+00, bound 5 sigma --- \emph{seeded noise, 100 replicates}
\end{itemize}
\item[fitCalibrationTimeVaryingPrecision] {\small\emph{precision \textperiodcentered\ \texttt{app.\allowbreak{}fitCalibration}}} --- Coloured noise and a 1 \%/day along-track drift (v3 WP-14's case): the recovered slope within 4 bootstrap sigma, the drift flagged, and the pointwise 2-sigma band covering the true S\_xx(t) at more than 90 \% of the fitted epochs.
\begin{itemize}[leftmargin=1.2em,itemsep=0pt,topsep=1pt]\small
\item fitCalibration scale drift error in bootstrap sigma: measured 1.0221e+00, bound 4 sigma --- \emph{seeded noise, 40 replicates; the slope of a linear S\_xx(t)}
\item fitCalibration S(t) outside its 2-sigma band, fraction of epochs: measured 0.0000e+00, bound 0.1 --- \emph{the v3 WP-14 criterion: the band covers the truth at more than 90 \% of epochs}
\end{itemize}
\item[fitCalibrationRobustness] {\small\emph{robustness \textperiodcentered\ \texttt{app.\allowbreak{}fitCalibration}}} --- Spikes and a gap: the fit runs over the windows the filter left, nothing else.
\item[fitCalibrationTimeVaryingRobustness] {\small\emph{robustness \textperiodcentered\ \texttt{app.\allowbreak{}fitCalibration}}} --- Knotted scale bases over a gap and spikes: every curve finite where it is evaluated, and a cubic B-spline curve on an arc shorter than its knot spacing still solves.
\item[fitCalibrationMetamorphic] {\small\emph{metamorphic \textperiodcentered\ \texttt{app.\allowbreak{}fitCalibration}}}
\begin{itemize}[leftmargin=1.2em,itemsep=0pt,topsep=1pt]\small
\item fitCalibration prediction + residual vs observation: measured 1.3156e-01, bound 1 --- \emph{two subtractions of the same fitted value; bound 4 eps max\textbar{}y\textbar{}}
\end{itemize}
\item[fitCalibrationTimeVaryingMetamorphic] {\small\emph{metamorphic \textperiodcentered\ \texttt{app.\allowbreak{}fitCalibration}}}
\begin{itemize}[leftmargin=1.2em,itemsep=0pt,topsep=1pt]\small
\item fitCalibration S(t), hat with one span vs linear: measured 3.0868e-03, bound 1 --- \emph{two parameterisations of one column space; bound 100 eps cond max\textbar{}S\textbar{}}
\end{itemize}
\item[fitCalibrationDayBootstrapSpeed] {\small\emph{speed \textperiodcentered\ \texttt{app.\allowbreak{}fitCalibration}}} --- With the per-day contributions the data pass is paid once and a replicate is a small solve, so 50 replicates must add little to a 10-day joint fit. The replicate-by-refit it replaced cost about one full fit per replicate (measured 2026-10-17: 870 s of bootstrap beside 350 s of fit on 10 real days at 100 replicates). Budget: 50 replicates within twice the fit without them.
\begin{itemize}[leftmargin=1.2em,itemsep=0pt,topsep=1pt]\small
\item fitCalibration, 10 days: 50 day-bootstrap replicates / no bootstrap: measured 1.0337e+00, bound 2 --- \emph{a replicate is lsqvce.solveContributions, O(N p\textasciicircum{}2); a refit per replicate would read \textasciitilde{}50}
\end{itemize}
\item[fitCalibrationSpeed] {\small\emph{speed \textperiodcentered\ \texttt{app.\allowbreak{}fitCalibration}}} --- Filter, gate and three solves all grow with the arc: twice the arc, about twice the cost.
\begin{itemize}[leftmargin=1.2em,itemsep=0pt,topsep=1pt]\small
\item fitCalibration, 2x arc / 1x arc: measured 1.8468e+00, bound 4 --- \emph{linear would be 2 and a fixed overhead pulls it lower; quadratic 4}
\end{itemize}
\item[fitCalibrationMultiArcContract] {\small\emph{contract \textperiodcentered\ \texttt{app.\allowbreak{}fitCalibration}}}
\item[fitCalibrationMultiArcReference] {\small\emph{reference \textperiodcentered\ \texttt{app.\allowbreak{}fitCalibration}}} --- Noise-free, three days, ONE planted scale and three different biases: the joint fit returns them all. The scale bound is the single-arc one over all rows (the common cofactor info.Qxx times the row rounding); a bias is a constant per day, recovered to the same rounding relative to the signal it sits beside.
\begin{itemize}[leftmargin=1.2em,itemsep=0pt,topsep=1pt]\small
\item \emph{(computed label)}: measured ---, bound 1 --- \emph{error / (2 sqrt(Q n) eRow + 100 eps cond sqrt(Q n) \textbar{}signal\textbar{}), n over all days}
\item fitCalibration multi-arc per-day biases, relative to the signal: measured 9.5291e-15, bound 1e-9 --- \emph{rounding of the CRN-filtered solve; measured \textasciitilde{}1e-15 when written}
\end{itemize}
\item[fitCalibrationMultiArcDayBootstrapPrecision] {\small\emph{precision \textperiodcentered\ \texttt{app.\allowbreak{}fitCalibration}}} --- A day-bootstrap replicate is lsqvce.solveContributions on the drawn multiplicities. Its oracle is the public joint fit over the EXPANDED day list -- a day drawn twice entered twice, as a copy of its data generated at another day's time (syntheticPair is shift-invariant) -- for the very draws timeSeriesPreparation.bootstrap makes (same seed, same stage name). Under white noise the two are one estimator, so the bootstrap sigma must agree to the solves' rounding.
\begin{itemize}[leftmargin=1.2em,itemsep=0pt,topsep=1pt]\small
\item day bootstrap sigma S\_xx: solveContributions vs full refits of the expanded day lists: measured 2.2286e-12, bound 1e-6 rel --- \emph{lsqvce: solveContributions equals a fit over the expanded days to 1e-9 in x; the spread of 8 replicates of \textasciitilde{}1e-4 keeps \textasciitilde{}5 digits of that}
\end{itemize}
\item[fitCalibrationMultiArcPrecision] {\small\emph{precision \textperiodcentered\ \texttt{app.\allowbreak{}fitCalibration}}} --- The algebra of pre-elimination: lsqvce's joint solve with each day's bias as a local block equals the DENSE joint least squares -- a block-diagonal bias design beside the common scale column, solved here by QR -- at the rounding of that solve, cond * eps.
\begin{itemize}[leftmargin=1.2em,itemsep=0pt,topsep=1pt]\small
\item fitCalibration multi-arc scale vs dense joint LS (relative): measured 3.8149e-08, bound 1 --- \emph{two solvers of one problem; bound 100 eps cond(X)}
\item fitCalibration multi-arc biases vs dense joint LS (relative to max bias): measured 7.2281e-08, bound 1 --- \emph{same solve, the local block}
\end{itemize}
\item[fitCalibrationVectorization] {\small\emph{vectorization \textperiodcentered\ \texttt{app.\allowbreak{}fitCalibration}}} --- 'ScaleMode' "perArc" is N single-arc fits, bitwise, arc k with seed Seed + k - 1.
\item[fitCalibrationMultiArcRobustness] {\small\emph{robustness \textperiodcentered\ \texttt{app.\allowbreak{}fitCalibration}}} --- Spikes and a gap inside every day (several sub-arcs per day), the held S\_yy axis that has no common column (it decouples into per-day fits), and the bootstrap over days: all finite.
\item[fitCalibrationMultiArcMetamorphic] {\small\emph{metamorphic \textperiodcentered\ \texttt{app.\allowbreak{}fitCalibration}}} --- Doubling every day's observation doubles the common scale and every bias of the free axes, bitwise (the equilibration takes a power of two out exactly); and the day's own reference is what matters -- passing one reference arc per day equals passing them as an array.
\item[fitCalibrationScaleTemperatureReference] {\small\emph{reference \textperiodcentered\ \texttt{app.\allowbreak{}fitCalibration}}} --- N-11b. Noise-free with a planted dS/dT on x and z: the coefficients return to the solve's rounding, at the planted reference (the mean over the raw grid on both sides). Over three days whose temperatures differ only by a daily OFFSET (5 mK of in-day swing), the common S1 returns too -- which a per-day reference could not give -- and the scale is S at the common T0. And apply undoes the coupling to the bound of exchanging the filter with S(T(t)) over its window.
\begin{itemize}[leftmargin=1.2em,itemsep=0pt,topsep=1pt]\small
\item fitCalibration dS/dT, one noise-free arc, relative: measured 3.5421e-13, bound 1e-9 --- \emph{rounding of the solve; measured 2.8e-13 on 2026-10-20}
\item Calibration apply of a fitted temperature-dependent scale vs H(model), relative: measured 2.2117e-04, bound bound --- \emph{the filter and S(T(t)) exchanged over one window: 2 \textbar{}S1\textbar{} max dT(window); measured 2.2e-4 at 0.5 C per orbit}
\item fitCalibration dS/dT, three days, common, relative: measured 2.3492e-11, bound 1e-8 --- \emph{noise-free; the daily offsets alone carry S1; measured 2.4e-11 on 2026-10-20}
\end{itemize}
\item[fitCalibrationRegressorReference] {\small\emph{reference \textperiodcentered\ \texttt{app.\allowbreak{}fitCalibration}}} --- Noise-free with a planted coupling to two measured series: every parameter, the coefficients included, returns to the computed rounding bound; over two days in common mode the coefficients are one set and return too.
\begin{itemize}[leftmargin=1.2em,itemsep=0pt,topsep=1pt]\small
\item fitCalibration multi-arc regressor coefficients, relative: measured 9.9874e-14, bound 1e-9 --- \emph{noise-free; rounding of the joint solve, measured \textasciitilde{}1e-13 when written}
\end{itemize}
\item[fitCalibrationScaleTemperatureContract] {\small\emph{contract \textperiodcentered\ \texttt{app.\allowbreak{}fitCalibration}}} --- The column names, the refusals, and the gate seeing the new columns.
\item[fitCalibrationScaleTemperatureRobustness] {\small\emph{robustness \textperiodcentered\ \texttt{app.\allowbreak{}fitCalibration}}} --- A housekeeping gap: the temperature series is missing over 20 minutes mid-arc. The rows it leaves undefined are not fitted, and the coefficients still return (noise-free).
\begin{itemize}[leftmargin=1.2em,itemsep=0pt,topsep=1pt]\small
\item fitCalibration dS/dT across a temperature gap, relative: measured 1.5613e-11, bound 1e-8 --- \emph{noise-free; fewer rows, same algebra}
\end{itemize}
\item[fitCalibrationRegressorRobustness] {\small\emph{robustness \textperiodcentered\ \texttt{app.\allowbreak{}fitCalibration}}} --- A constant channel is refused before any fit, paired with the live series accepted; a series covering only part of the arc fits the rows it covers.
\item[fitCalibrationScaleTemperatureMetamorphic] {\small\emph{metamorphic \textperiodcentered\ \texttt{app.\allowbreak{}fitCalibration}}} --- The temperature in another unit, T' = 10 T + 5: dS/dT' = dS/dT / 10, T0' = 10 T0 + 5, and the scale at the reference is the same -- one physical coupling, two descriptions.
\item[fitCalibrationRegressorMetamorphic] {\small\emph{metamorphic \textperiodcentered\ \texttt{app.\allowbreak{}fitCalibration}}} --- A series in other units (x 10) has its coefficients divided by 10, and an offset (+ 5) changes nothing but the stored mean: the column is demeaned.
\begin{itemize}[leftmargin=1.2em,itemsep=0pt,topsep=1pt]\small
\item fitCalibration regressor x10 -\textgreater{} coefficient / 10: measured 2.7802e-14, bound 1e-9 --- \emph{a rescaled column in an equilibrated solve; rounding only}
\item fitCalibration regressor + 5 -\textgreater{} same coefficients: measured 1.6481e-14, bound 1e-9 --- \emph{the column is demeaned; only the rounding of the interpolation differs}
\end{itemize}
\end{description}

\subsection{\texttt{tNgfGeom}}
\noindent Tests of geom.MacroModel, geom.graceFO and geom.grace.

\begin{description}[leftmargin=1em,style=nextline,itemsep=2pt,font=\normalfont\ttfamily\small]
\item[macroModelContract] {\small\emph{contract \textperiodcentered\ \texttt{geom.\allowbreak{}MacroModel}}}
\item[graceFOContract] {\small\emph{contract \textperiodcentered\ \texttt{geom.\allowbreak{}graceFO}}}
\item[graceFOReference] {\small\emph{reference \textperiodcentered\ \texttt{geom.\allowbreak{}graceFO}}} --- D-56935 Table 5: total area and the nadir-plate optics.
\item[macroModelReference] {\small\emph{reference \textperiodcentered\ \texttt{geom.\allowbreak{}MacroModel}}} --- Side normals use sin/cos 50 deg as printed (panels), renormalised to unit length in N (O-2, 2026-09-07).
\begin{itemize}[leftmargin=1.2em,itemsep=0pt,topsep=1pt]\small
\item renormalised normal vs printed digits: measured 7.3135e-07, bound 1e-6 --- \emph{six printed digits: their norm is 1 - 7.3e-7 (measured 2026-09-07; the record had declared 5e-7)}
\end{itemize}
\item[graceFOPrecision] {\small\emph{precision \textperiodcentered\ \texttt{geom.\allowbreak{}graceFO}}} --- The silhouette defect is the sum of the PRINTED area vectors, which do not close on each other (a closed body, convex or not, gives zero): in Z the flanks' 2*A*c50 terms plus nadir minus zenith; in Y the inner and outer flank normals cancel. Bound 1e-9 m\textasciicircum{}2: arithmetic. c50 is taken from the packed unit normal (O-2, 2026-09-07), not from the printed 0.642787 (norm 1 - 7.3e-7).
\begin{itemize}[leftmargin=1.2em,itemsep=0pt,topsep=1pt]\small
\item graceFO silhouette defect z: measured 0.0000e+00, bound 1e-9 m\textasciicircum{}2
\item graceFO silhouette defect xy: measured 1.3878e-16, bound 1e-9 m\textasciicircum{}2
\end{itemize}
\item[macroModelPrecision] {\small\emph{precision \textperiodcentered\ \texttt{geom.\allowbreak{}MacroModel}}} --- Energy closure is enforced at 1e-6; the handbook numbers close to 1e-12.
\begin{itemize}[leftmargin=1.2em,itemsep=0pt,topsep=1pt]\small
\item MacroModel energy closure vis: measured 0.0000e+00, bound 1e-12
\item MacroModel energy closure ir: measured 0.0000e+00, bound 1e-12
\end{itemize}
\item[macroModelRobustness] {\small\emph{robustness \textperiodcentered\ \texttt{geom.\allowbreak{}MacroModel}}} --- Unknown centroids (all NaN) are accepted; they only disable torques.
\item[graceFORobustness] {\small\emph{robustness \textperiodcentered\ \texttt{geom.\allowbreak{}graceFO}}} --- Non-operating array closes the energy budget by construction.
\item[macroModelMetamorphic] {\small\emph{metamorphic \textperiodcentered\ \texttt{geom.\allowbreak{}MacroModel}}} --- Rotating every normal and centroid by R rotates the silhouette defect by R.
\begin{itemize}[leftmargin=1.2em,itemsep=0pt,topsep=1pt]\small
\item MacroModel rotation covariance: measured 1.0155e-15, bound 1e-12 m\textasciicircum{}2
\end{itemize}
\item[graceFOMetamorphic] {\small\emph{metamorphic \textperiodcentered\ \texttt{geom.\allowbreak{}graceFO}}} --- C and D differ only in mass properties, never in geometry.
\item[graceContract] {\small\emph{contract \textperiodcentered\ \texttt{geom.\allowbreak{}grace}}}
\item[graceReference] {\small\emph{reference \textperiodcentered\ \texttt{geom.\allowbreak{}grace}}} --- R-5 Sec. I.5.3.3 (launch mass) and I.5.3.5 (surface table).
\item[graceReferenceDimensions] {\small\emph{reference \textperiodcentered\ \texttt{geom.\allowbreak{}grace}}} --- The polygons against numbers R-5 prints and geom.grace does not use: - Fig 3's 56 mm: depth of the overhang below the nadir plane (geom.grace uses the printed underside AREA, 0.2282913 m\textasciicircum{}2, and the 50 deg slant); - Sec. I.5.3.2: the orbit thrusters O11/O21 at X = -1561 mm on the rear face and the attitude thrusters A14/A22 at Z = +275 mm on the nadir face, A12/A24 at Z = -444 mm; - Fig 4's projected side area, 2.4633921 m\textasciicircum{}2, which includes the boom (0.0461901 m\textasciicircum{}2, Sec. I.5.3.5 note 4) and the dome, notch and recesses the plates do not model. Dimensions are printed to the millimetre, so 1 mm (plus the half-millimetre of the odd 3123 mm length) is the bound; the side-area bound is the 1 \% ...
\begin{itemize}[leftmargin=1.2em,itemsep=0pt,topsep=1pt]\small
\item grace overhang depth vs Fig 3 (56 mm): measured 2.1836e-06, bound 1e-3 m
\item grace nadir plane vs thrusters A14/A22 (Z = +275 mm): measured 0.0000e+00, bound 1e-3 m
\item grace zenith plane vs thrusters A12/A24 (Z = -444 mm): measured 1.0000e-03, bound 2e-3 m --- \emph{measured 1 mm on CI run 35773012179; millimetre rounding of three printed values allows 1.5 mm}
\item grace rear face vs orbit thrusters O11/O21 (X = -1561 mm): measured 5.0000e-04, bound 1e-3 m
\item grace side silhouette vs Fig 4 C.o.A. area: measured 2.5328e-03, bound 1e-2 --- \emph{measured 2.53e-3 in Python on 2026-09-22 before any MATLAB run}
\end{itemize}
\item[gracePrecision] {\small\emph{precision \textperiodcentered\ \texttt{geom.\allowbreak{}grace}}} --- Same identity as graceFOPrecision, on R-5's table: inner and outer flank normals cancel in Y, the Z defect is 2*A\_inner*c50 - 2*A\_outer*c50 + A\_nadir - A\_zenith. Arithmetic: 1e-9 m\textasciicircum{}2.
\begin{itemize}[leftmargin=1.2em,itemsep=0pt,topsep=1pt]\small
\item grace silhouette defect z: measured 0.0000e+00, bound 1e-9 m\textasciicircum{}2
\item grace silhouette defect xy: measured 1.3878e-16, bound 1e-9 m\textasciicircum{}2
\item grace energy closure: measured 0.0000e+00, bound 1e-12
\end{itemize}
\item[graceRobustness] {\small\emph{robustness \textperiodcentered\ \texttt{geom.\allowbreak{}grace}}} --- Note 2's non-operating array closes by construction (the printed 0.72/0.05/0.30 would sum to 1.07).
\item[graceMetamorphic] {\small\emph{metamorphic \textperiodcentered\ \texttt{geom.\allowbreak{}grace}}} --- A and B are one table (R-5 Sec. I.5.3.3): everything but the name is identical.
\item[graceMetamorphicAcrossDocuments] {\small\emph{metamorphic \textperiodcentered\ \texttt{geom.\allowbreak{}grace}}} --- R-5 Sec. I.5.3.5 and JPL D-56935 Table 5 print the same surface table. The two factories transcribe them independently, so this compares two transcriptions of one fact; the geometry, mass and origin differ by document and are NOT compared.
\end{description}

\subsection{\texttt{tNgfGeometry}}
\noindent Tests of the LP-4 geom units: boxWing, selfShadow, thermalState, AttitudeProvider, checkAttitudeSense.

\begin{description}[leftmargin=1em,style=nextline,itemsep=2pt,font=\normalfont\ttfamily\small]
\item[flowShadowTableContract] {\small\emph{contract \textperiodcentered\ \texttt{geom.\allowbreak{}flowShadow}}} --- Method "table": the table itself (both cones), the struct it returns, and what it refuses.
\item[flowShadowContract] {\small\emph{contract \textperiodcentered\ \texttt{geom.\allowbreak{}flowShadow}}}
\item[boxWingContract] {\small\emph{contract \textperiodcentered\ \texttt{geom.\allowbreak{}boxWing}}}
\item[selfShadowContract] {\small\emph{contract \textperiodcentered\ \texttt{geom.\allowbreak{}selfShadow}}}
\item[thermalStateContract] {\small\emph{contract \textperiodcentered\ \texttt{geom.\allowbreak{}thermalState}}}
\item[attitudeProviderContract] {\small\emph{contract \textperiodcentered\ \texttt{geom.\allowbreak{}AttitudeProvider}}}
\item[checkAttitudeSenseContract] {\small\emph{contract \textperiodcentered\ \texttt{geom.\allowbreak{}checkAttitudeSense}}}
\item[flowShadowTableReference] {\small\emph{reference \textperiodcentered\ \texttt{geom.\allowbreak{}flowShadow}}} --- The table against the computation it tabulates (Method lut), over GRACE-FO's stream: yaw +-4 deg, pitch +-1 deg about both +x and -x, s\_O 7..10. Measured 2026-10-17 on 81952 epochs of ten real days (C and D): 1.2e-3 at most, which moved the drag by 6.7e-13 m/s\textasciicircum{}2 of 2.1e-8. Bound 4x.
\begin{itemize}[leftmargin=1.2em,itemsep=0pt,topsep=1pt]\small
\item flowShadow table vs Method lut, every ratio, GRACE-FO cone (400 epochs): measured 1.1079e-03, bound 5e-3 --- \emph{trilinear at 0.5 deg and 0.5 in s\_O; measured 1.2e-3 on 81952 real epochs (2026-10-17)}
\end{itemize}
\item[flowShadowReference] {\small\emph{reference \textperiodcentered\ \texttt{geom.\allowbreak{}flowShadow}}} --- Oracle: tools/pyxcheck/tpmc\_check.py, a standard-library test-particle Monte Carlo of the incident number, normal-momentum and tangential-momentum fluxes on a tiny plate screened by a wall, written from a physics statement by an agent that did not open this toolbox (T9). The model is the flow ratios times Sentman's incident closed forms. Measured 2026-09-15 at quadrature [96 128] with an exact ray cast: worst \textbar{}z\textbar{} = 2.58 over the components the Monte Carlo resolves (standard error below 10 \% of the value). The occlusion is a step in direction, which the quadrature resolves to its node spacing: at [24 32] the same cases missed by up to 20 standard errors and converged as the nodes were ...
\begin{itemize}[leftmargin=1.2em,itemsep=0pt,topsep=1pt]\small
\item flowShadow vs test-particle Monte Carlo, worst \textbar{}z\textbar{} (resolved components): measured 2.5823e+00, bound 4 --- \emph{measured 2.58 on 2026-09-15; 18 cases, s = 2, 5, 8, elevation -10 to 90 deg}
\end{itemize}
\item[boxWingReference] {\small\emph{reference \textperiodcentered\ \texttt{geom.\allowbreak{}boxWing}}} --- The bus of a 1 m cube is the unit cube: six unit areas, face centres at +-1/2, closed silhouette.
\item[selfShadowReference] {\small\emph{reference \textperiodcentered\ \texttt{geom.\allowbreak{}selfShadow}}} --- Oracle: the pure-Python ray cast (160 x 160 per plate). Measured 2026-09-07: 40 samples 1.25e-2, 100 samples 3.8e-3; the sampling error bound is \textasciitilde{}2/m at a shadow edge: 5e-2 and 2e-2. Bounds 3e-2 and 1e-2. The oracle writes 0 for a plate facing away (the pre-LP-C convention); occlusion is compared where the plate faces the direction, which is the only place any consumer weights it.
\begin{itemize}[leftmargin=1.2em,itemsep=0pt,topsep=1pt]\small
\item selfShadow raycast 40 vs Python 160 (lit plates): measured 1.2500e-02, bound 3e-2 --- \emph{measured 1.25e-2}
\item selfShadow raycast 100 vs Python 160 (lit plates): measured 3.7500e-03, bound 1e-2 --- \emph{measured 3.75e-3}
\end{itemize}
\item[thermalStateReference] {\small\emph{reference \textperiodcentered\ \texttt{geom.\allowbreak{}thermalState}}} --- Oracle: Stefan-Boltzmann in closed form. One plate normal to the Sun at 1 au with alpha 0.34, eps 0.62: T = (0.34 * 1361 / (0.62 sigma))\textasciicircum{}(1/4).
\begin{itemize}[leftmargin=1.2em,itemsep=0pt,topsep=1pt]\small
\item steady-state temperature vs closed form: measured 0.0000e+00, bound 1e-9 K --- \emph{one fourth root}
\end{itemize}
\item[attitudeProviderReference] {\small\emph{reference \textperiodcentered\ \texttt{geom.\allowbreak{}AttitudeProvider}}} --- Oracle: the real 2020-01-01 attitude. The nominal Earth-pointing frame with +X against the velocity is within 1 deg of the measured SCA1B attitude (measured 0.84 deg mean, 0.86 max); with +X along it, 180 deg.
\begin{itemize}[leftmargin=1.2em,itemsep=0pt,topsep=1pt]\small
\item nominal attitude (+X against v) vs SCA1B: measured 8.6093e-01, bound 3 deg --- \emph{pointing offsets \textasciitilde{}1 deg; measured 0.86}
\end{itemize}
\item[checkAttitudeSenseReference] {\small\emph{reference \textperiodcentered\ \texttt{geom.\allowbreak{}checkAttitudeSense}}} --- The measured GRACE-C attitude: nadir within 1 deg, flying with +X against the velocity (JPL D-56935: GRACE-C trails GRACE-D with its K-band horn towards it on 2020-01-01, measured in tests/tNgfIo.m as 179.2 deg).
\begin{itemize}[leftmargin=1.2em,itemsep=0pt,topsep=1pt]\small
\item checkAttitudeSense max nadir angle, real day: measured 8.6018e-01, bound 3 deg --- \emph{measured 0.86}
\end{itemize}
\item[thermalStateOcclusionPrecision] {\small\emph{precision \textperiodcentered\ \texttt{geom.\allowbreak{}thermalState}}} --- (LP-C): the absorbed flux with per-source occlusion is the flux radiationPanelForce forms from the same sources and the same occlusion, source by source; and a panel facing away from the Sun keeps its Earth heating.
\begin{itemize}[leftmargin=1.2em,itemsep=0pt,topsep=1pt]\small
\item thermalState absorbed flux vs radiationPanelForce's, per-source occlusion: measured 1.8086e-16, bound 1e-13 --- \emph{the same product a P c gamma o summed over 577 sources in two association orders}
\item nadir absorbed flux: per-source minus Sun-only occlusion (F2 regression): measured -5.3882e+01, bound 0 W/m\textasciicircum{}2 --- \emph{per-source occlusion never absorbs less than the Sun column forced onto the Earth tiles}
\end{itemize}
\item[flowShadowTablePrecision] {\small\emph{precision \textperiodcentered\ \texttt{geom.\allowbreak{}flowShadow}}} --- At its own nodes the table returns what it stored, and what it stored is Method lut at that node: a node read back differs from the exact evaluation there only by the rounding of the node's coordinates.
\begin{itemize}[leftmargin=1.2em,itemsep=0pt,topsep=1pt]\small
\item flowShadow table at its nodes vs Method lut there: measured 0.0000e+00, bound 1e-10 --- \emph{node coordinates rebuilt from v/\textbar{}v\textbar{} carry a few ulp, times a slope below 1e2 per unit of tan}
\end{itemize}
\item[flowShadowPrecision] {\small\emph{precision \textperiodcentered\ \texttt{geom.\allowbreak{}flowShadow}}} --- Without occlusion the directional integrals ARE Sentman's incident coefficients (number flux, normal and tangential momentum): an identity of kinetic theory, so the comparison measures the quadrature alone. Measured 2026-09-15 at the default [12 16]: worst 1.6e-6 in coefficient units over s = 1..15 and incidence 0..180 deg (a ram plate's coefficient is \textasciitilde{}2). Bound 5e-6. On two plates that cannot see each other every ratio is exactly 1.
\begin{itemize}[leftmargin=1.2em,itemsep=0pt,topsep=1pt]\small
\item flowShadow integrals vs Sentman incident closed forms: measured 2.2656e-07, bound 5e-6 --- \emph{measured 1.6e-6 on 2026-09-15 at [12 16]; [8 16] gives 3e-4}
\end{itemize}
\item[boxWingPrecision] {\small\emph{precision \textperiodcentered\ \texttt{geom.\allowbreak{}boxWing}}} --- Polygon corners reproduce the face areas exactly (products of two edge lengths, bitwise).
\item[selfShadowPrecision] {\small\emph{precision \textperiodcentered\ \texttt{geom.\allowbreak{}selfShadow}}} --- Analytic two-plate case: a unit square half covered by a parallel unit square shifted by 1/2 has exposed fraction exactly 1/2 for a beam along +z; with an even sample count the grid straddles the edge exactly.
\begin{itemize}[leftmargin=1.2em,itemsep=0pt,topsep=1pt]\small
\item selfShadow two plates half covered: measured 3.3307e-16, bound 1e-12 --- \emph{edge on a grid line: exact}
\item selfShadow two plates, odd grid: measured 1.2821e-02, bound 1/39 + 1e-12 --- \emph{one straddling row}
\end{itemize}
\item[thermalStatePrecision] {\small\emph{precision \textperiodcentered\ \texttt{geom.\allowbreak{}thermalState}}} --- Transient against the linearised relaxation of a 1 K excursion on the Front plate (constant source): T - Ts = exp(-t/tau), tau = C/(4 eps sigma Ts\textasciicircum{}3). The nonlinearity of T\textasciicircum{}4 for 1 K at 300 K is \textasciitilde{}1 \%; RK4 with dt \textless{}= tau/10 adds \textless{} 1e-5. Bound 5e-2 K (measured see last\_run). Untouched plates stay exactly.
\begin{itemize}[leftmargin=1.2em,itemsep=0pt,topsep=1pt]\small
\item transient vs linearised relaxation (1 K): measured 1.1063e-03, bound 5e-3 K --- \emph{T\textasciicircum{}4 nonlinearity of a 1 K step; measured 1.1e-3}
\item transient: steady plates stay put: measured 0.0000e+00, bound 1e-12 K --- \emph{fixed point of RK4}
\end{itemize}
\item[attitudeProviderPrecision] {\small\emph{precision \textperiodcentered\ \texttt{geom.\allowbreak{}AttitudeProvider}}} --- A slerp query exactly at a node returns the node bitwise (referenceSystem reproduces nodes to 1e-16 and the node value is the input), and the nominal attitude round-trips through its rotation matrix to 1e-15.
\begin{itemize}[leftmargin=1.2em,itemsep=0pt,topsep=1pt]\small
\item AttitudeProvider at the nodes: measured 0.0000e+00, bound 2e-16 --- \emph{renormalisation of a unit quaternion}
\item nominal attitude +Z is nadir: measured 3.3307e-16, bound 1e-14 --- \emph{rotation round trip}
\end{itemize}
\item[checkAttitudeSensePrecision] {\small\emph{precision \textperiodcentered\ \texttt{geom.\allowbreak{}checkAttitudeSense}}} --- The nominal frame points +Z exactly at nadir: the angle is rounding only. With atan2 (not acos) the rounding is 1e-14 deg (acos would give sqrt(eps) \textasciitilde{} 1e-6 deg, measured 1.2e-6 before the change).
\begin{itemize}[leftmargin=1.2em,itemsep=0pt,topsep=1pt]\small
\item checkAttitudeSense nadir angle of the nominal frame: measured 2.1572e-14, bound 1e-9 deg --- \emph{atan2 of cross and dot: rounding only}
\item checkAttitudeSense along-track angle of the nominal frame: measured 7.5903e-02, bound 1 deg --- \emph{flight-path angle of the near-circular orbit (\textless{} 0.2 deg)}
\end{itemize}
\item[flowShadowTableVectorization] {\small\emph{vectorization \textperiodcentered\ \texttt{geom.\allowbreak{}flowShadow}}} --- Element by element, inside and outside the table alike: a batch equals its epochs one at a time, bitwise.
\item[flowShadowVectorization] {\small\emph{vectorization \textperiodcentered\ \texttt{geom.\allowbreak{}flowShadow}}} --- five epochs at once equal five single epochs, bitwise: every sum runs along a dimension whose length does not depend on the number of epochs
\item[selfShadowVectorization] {\small\emph{vectorization \textperiodcentered\ \texttt{geom.\allowbreak{}selfShadow}}}
\item[thermalStateVectorization] {\small\emph{vectorization \textperiodcentered\ \texttt{geom.\allowbreak{}thermalState}}}
\item[attitudeProviderVectorization] {\small\emph{vectorization \textperiodcentered\ \texttt{geom.\allowbreak{}AttitudeProvider}}}
\item[checkAttitudeSenseVectorization] {\small\emph{vectorization \textperiodcentered\ \texttt{geom.\allowbreak{}checkAttitudeSense}}}
\item[flowShadowTableRobustness] {\small\emph{robustness \textperiodcentered\ \texttt{geom.\allowbreak{}flowShadow}}} --- Outside the table -- a stream 20 deg off the axis, side-on, or a speed ratio beyond [5, 12] -- the epoch is computed exactly, bitwise as Method lut computes it, and flagged; the epochs inside are not disturbed.
\item[flowShadowRobustness] {\small\emph{robustness \textperiodcentered\ \texttt{geom.\allowbreak{}flowShadow}}} --- the stream exactly along a plate normal (no tangential direction), exactly edge-on, exactly leeward, a very cold and a very hot gas, and no epochs at all: finite ratios in [0, 1], and exactly 1 on GRACE-FO plates nothing screens
\item[boxWingRobustness] {\small\emph{robustness \textperiodcentered\ \texttt{geom.\allowbreak{}boxWing}}} --- A flat box (thin plate), a tiny box, an array without a name, and a polygon check on a hand-built model.
\item[selfShadowRobustness] {\small\emph{robustness \textperiodcentered\ \texttt{geom.\allowbreak{}selfShadow}}} --- Exactly edge-on beams, a beam along the body axis, and the two-plate model with no overlap and with full overlap. Along +x nothing obstructs any plate: the Front plate faces it, the Rear plate's mirror direction is -x, and the six edge-on plates take their front-side limit (1e-6 rad out of plane), which clears the Front plate's edge.
\item[thermalStateRobustness] {\small\emph{robustness \textperiodcentered\ \texttt{geom.\allowbreak{}thermalState}}} --- No sources: 0 K steady, and a transient cooling from 300 K decays monotonically; an exposed-fraction array scales the absorbed flux; NaN pressures refuse.
\item[attitudeProviderRobustness] {\small\emph{robustness \textperiodcentered\ \texttt{geom.\allowbreak{}AttitudeProvider}}} --- Two nodes suffice; a query in another time scale is converted; squad and slerp agree at the nodes and differ between them only at the interpolation level for this slowly turning attitude.
\begin{itemize}[leftmargin=1.2em,itemsep=0pt,topsep=1pt]\small
\item AttitudeProvider UTC vs GPS query: measured 0.0000e+00, bound 1e-12 --- \emph{epoch arithmetic rounding}
\item squad vs slerp between 1 Hz nodes: measured 3.0705e-08, bound 1e-6 --- \emph{orbital rate 1.1e-3 rad/s: cubic vs linear over 0.5 s}
\end{itemize}
\item[checkAttitudeSenseRobustness] {\small\emph{robustness \textperiodcentered\ \texttt{geom.\allowbreak{}checkAttitudeSense}}} --- A tolerance below the real pointing offset turns ok false without an error; a single epoch works; mixed sense reported.
\item[flowShadowTableMetamorphic] {\small\emph{metamorphic \textperiodcentered\ \texttt{geom.\allowbreak{}flowShadow}}} --- Switching geom.flowShadow from exact to table may move the DRAG and nothing else: radiation and thermal recoil are bitwise unchanged, and the drag moves, but by no more than the table's error carries.
\begin{itemize}[leftmargin=1.2em,itemsep=0pt,topsep=1pt]\small
\item Model drag, flowShadow table vs exact (sample arc): measured 6.7914e-06, bound 1e-4 rel --- \emph{measured 6.7e-13 m/s\textasciicircum{}2 of 4.5e-8 on real days (1.5e-5); bound \textasciitilde{}7x}
\end{itemize}
\item[flowShadowMetamorphic] {\small\emph{metamorphic \textperiodcentered\ \texttt{geom.\allowbreak{}flowShadow}}} --- (1) the beam field IS selfShadow along v\_rel, bitwise; (2) mirroring the stream in y swaps starboard and port; (3) as the gas cools (s -\textgreater{} inf) every ratio converges on the beam value: the error falls as s rises.
\begin{itemize}[leftmargin=1.2em,itemsep=0pt,topsep=1pt]\small
\item flowShadow y-mirror swaps starboard and port: measured 5.5511e-16, bound 1e-12 --- \emph{the cube map is symmetric in y to the ray cast's sampling, which is symmetric; the flow frame is not, so rounding enters}
\item flowShadow normal ratio vs beam at s = 400: measured 5.4173e-05, bound 2e-2 --- \emph{spread 1/s = 2.5e-3 rad across a cube cell of 0.2 rad with a value jump \textless{}= 1}
\end{itemize}
\item[boxWingMetamorphic] {\small\emph{metamorphic \textperiodcentered\ \texttt{geom.\allowbreak{}boxWing}}} --- Scaling every dimension by s scales every area by s\textasciicircum{}2 and every centroid by s; permuting the dimensions permutes the faces.
\begin{itemize}[leftmargin=1.2em,itemsep=0pt,topsep=1pt]\small
\item boxWing area scaling: measured 0.0000e+00, bound 1e-14 m\textasciicircum{}2 --- \emph{products of two lengths}
\end{itemize}
\item[selfShadowMetamorphic] {\small\emph{metamorphic \textperiodcentered\ \texttt{geom.\allowbreak{}selfShadow}}} --- (1) sampling convergence: the 40- and 100-sample casts differ by less than the sampling bound 2/40; (2) the table read at its own nodes returns the stored values (bilinear weights (1,0,0,0)); (3) the table error decreases with resolution; (4) mirror symmetry of GRACE-FO in y swaps starboard and port; (5) a back-facing direction returns exactly the value of its mirror direction in the plate's plane.
\begin{itemize}[leftmargin=1.2em,itemsep=0pt,topsep=1pt]\small
\item selfShadow 40 vs 100 samples: measured 1.0000e-02, bound 2 / 40 --- \emph{one edge row of the coarser grid}
\item cube map at its nodes vs the table: measured 1.0408e-17, bound 1e-12 --- \emph{a node is read back through two arctangents; one ulp in the face angle moves the weight by 1e-16, times a jump of at most 1 between neighbours}
\item cube map 16 mean error vs ray cast (40 random directions): measured 3.0586e-03, bound 1e-2 --- \emph{measured 3.1e-3 on 2026-09-15 (the 642-vertex icosphere it replaces measured 2.6e-2: the zeroed back faces smeared the horizon)}
\item cube map 16 / cube map 4 mean error ratio: measured 1.1210e-01, bound 0.5 --- \emph{first-order convergence in the cell size gives 0.25; measured 0.11}
\item back-facing direction vs its mirror (nadir plate): measured 0.0000e+00, bound 0 --- \emph{the same rays are cast: the mirror of the mirror is the direction itself to rounding, and the occlusion test is discrete}
\end{itemize}
\item[thermalStateMetamorphic] {\small\emph{metamorphic \textperiodcentered\ \texttt{geom.\allowbreak{}thermalState}}} --- Energy balance in steady state (emitted = absorbed + internal) exactly; doubling every pressure raises every steady temperature by 2\textasciicircum{}(1/4); a transient run from the steady state stays there.
\begin{itemize}[leftmargin=1.2em,itemsep=0pt,topsep=1pt]\small
\item steady-state energy balance: measured 5.8077e-16, bound 1e-13 rel --- \emph{fourth root and fourth power}
\item steady temperature scales with P\textasciicircum{}(1/4): measured 1.1369e-13, bound 1e-9 K
\item transient from the steady state stays: measured 0.0000e+00, bound 1e-9 K --- \emph{fixed point}
\end{itemize}
\item[attitudeProviderMetamorphic] {\small\emph{metamorphic \textperiodcentered\ \texttt{geom.\allowbreak{}AttitudeProvider}}} --- (1) the nominal attitude of a state rotated by R is the rotated nominal attitude (R q); (2) reversing the velocity sense flips the +X choice: nominal(r, v, true) == nominal(r, -v, false) up to sign; (3) a provider built from q reproduces the nominal frame it was built from at every node.
\begin{itemize}[leftmargin=1.2em,itemsep=0pt,topsep=1pt]\small
\item nominal attitude rotates with the state: measured 4.9960e-16, bound 1e-14 --- \emph{one rotation}
\item sense flag vs reversed velocity: measured 1.1102e-16, bound 1e-14 --- \emph{same rotation up to the quaternion sign}
\item provider reproduces its own nominal nodes: measured 1.1102e-16, bound 1e-15 --- \emph{renormalisation of unit quaternions}
\end{itemize}
\item[checkAttitudeSenseMetamorphic] {\small\emph{metamorphic \textperiodcentered\ \texttt{geom.\allowbreak{}checkAttitudeSense}}} --- Rotating the whole scene (state and attitude) by R leaves every angle; swapping the sense flag of a nominal attitude turns alongVelocity into againstVelocity with along-track angles summing to 180 deg.
\begin{itemize}[leftmargin=1.2em,itemsep=0pt,topsep=1pt]\small
\item sense angles invariant under a scene rotation: measured 2.2649e-14, bound 1e-9 deg --- \emph{two rotations and a quaternion round trip}
\item along-track angles of the two senses sum to 180: measured 2.8422e-14, bound 1e-9 deg
\end{itemize}
\item[flowShadowTableSpeed] {\small\emph{speed \textperiodcentered\ \texttt{geom.\allowbreak{}flowShadow}}} --- The table read against the computation it replaces, 200 epochs, after the session's table is built. Measured 2026-10-17: 2-4 us against 450 us per epoch. Budget 5x the measured 0.01.
\begin{itemize}[leftmargin=1.2em,itemsep=0pt,topsep=1pt]\small
\item flowShadow table / Method lut time ratio (200 epochs): measured 1.6350e-02, bound 0.05 --- \emph{measured \textasciitilde{}0.01 on 2026-10-17; budget 5x}
\end{itemize}
\item[flowShadowSpeed] {\small\emph{speed \textperiodcentered\ \texttt{geom.\allowbreak{}flowShadow}}} --- flowShadow for 32 epochs against the radiation look-up the same Model epoch needs (577 source directions x 32): the gas quadrature must stay within an order of magnitude of the radiation occlusion it sits beside. Measured 2026-09-15: see last\_run; budget 3x the measured ratio.
\begin{itemize}[leftmargin=1.2em,itemsep=0pt,topsep=1pt]\small
\item flowShadow (32 epochs) / radiation occlusion look-up (577 x 32): measured 4.8731e+00, bound 16 --- \emph{measured 5.3 on 2026-09-15; budget 3x}
\end{itemize}
\item[checkAttitudeSenseSpeed] {\small\emph{speed \textperiodcentered\ \texttt{geom.\allowbreak{}checkAttitudeSense}}} --- Its whole cost is two rotations of the same quaternion series plus two atan2 angles, so it must stay within a small multiple of one rotationFromQuaternion over the same K.
\begin{itemize}[leftmargin=1.2em,itemsep=0pt,topsep=1pt]\small
\item checkAttitudeSense vs one rotationFromQuaternion over 600 epochs: measured 7.7262e+00, bound 12 --- \emph{budget 3x the expected 4: two rotations, one quaternion conjugation and two atan2 angles; 5 repeats, median}
\end{itemize}
\item[selfShadowSpeed] {\small\emph{speed \textperiodcentered\ \texttt{geom.\allowbreak{}selfShadow}}} --- A table query against the ray cast of the same directions. The cube map is O(1) per direction (two arctangents and four reads); measured 2026-09-15: 3.6 us per direction against 1.5 ms per direction for the 40-sample cast, ratio 2.4e-3. Budget 3x: 1e-2. (The icosphere it replaced scanned every face per direction.)
\begin{itemize}[leftmargin=1.2em,itemsep=0pt,topsep=1pt]\small
\item selfShadow cube-map query / ray cast time ratio (600 directions): measured 9.8321e-04, bound 1e-2 --- \emph{O(1) look-up; measured 2.4e-3}
\end{itemize}
\item[thermalStateSpeed] {\small\emph{speed \textperiodcentered\ \texttt{geom.\allowbreak{}thermalState}}} --- Transient over the steady state for the 600-epoch sample: the RK4 loop with dt \textless{}= tau/10 substeps. Measured 2026-09-07: 0.02 s for 400 epochs; bound 20x the steady state.
\begin{itemize}[leftmargin=1.2em,itemsep=0pt,topsep=1pt]\small
\item thermalState transient / steady time ratio (600 epochs): measured 1.5024e+00, bound 20 --- \emph{RK4 substeps}
\end{itemize}
\item[attitudeProviderSpeed] {\small\emph{speed \textperiodcentered\ \texttt{geom.\allowbreak{}AttitudeProvider}}} --- at() over referenceSystem.interpolateQuaternions called directly for 5000 queries: the wrapper adds a scale check and a struct field. Measured 2026-09-07: 5.6 in a cold probe (the three-output call of the dependency computes its diagnostics); bound 8, recorded per run.
\begin{itemize}[leftmargin=1.2em,itemsep=0pt,topsep=1pt]\small
\item AttitudeProvider.at / interpolateQuaternions time ratio: measured 1.0283e+00, bound 8 --- \emph{three-output form of the dependency}
\end{itemize}
\end{description}

\subsection{\texttt{tNgfIntegration}}
\noindent Tests of app.DynamicsSource: the integrator interface and its native C evaluator (LP-C, WP-C9).

\begin{description}[leftmargin=1em,style=nextline,itemsep=2pt,font=\normalfont\ttfamily\small]
\item[DynamicsSourceContract] {\small\emph{contract \textperiodcentered\ \texttt{app.\allowbreak{}DynamicsSource}}}
\item[DynamicsSourceReference] {\small\emph{reference \textperiodcentered\ \texttt{app.\allowbreak{}DynamicsSource}}} --- Model.accel is the reference and the native evaluator is the port: at a table node the two are two implementations of the same expressions, and nothing else differs. Seven configurations, chosen to reach every branch the native path serves: exponential density with constantCd and no shadowing (fast); msis00, sentman, Knocke tiles, the cube-map shadow and the transient walls (standard); msis21, CERES, SOLAARS-CF, a 32-cell table and transient walls (precise); HWM14 winds, maxwell, the point source and the lunar shadow (hwm14). The last two exist only to reach the accommodation enum: MATLAB and C number it independently, so a value no configuration selects would let both paths run and ...
\begin{itemize}[leftmargin=1.2em,itemsep=0pt,topsep=1pt]\small
\item native evaluator vs Model.accel at the table nodes (7 configurations): measured 1.1398e-15, bound 1e-13 --- \emph{accumulated rounding of two implementations of one expression set, written in the same association order where it decides the digits; measured 2.504e-14 on 2026-09-16}
\item native evaluator vs Model.accel between the table nodes (60 s tables): measured 8.1162e-11, bound 1e-9 --- \emph{the cubic tables' interpolation error, not the model's: measured 8.124e-11 at 60 s spacing, and it falls as the fourth power of the step}
\item native evaluator vs Model.accel between the table nodes, transient walls (60 s tables): measured 2.7676e-05, bound 1e-4 --- \emph{the wall-temperature table's cubic between nodes; measured 1.564e-05 (CI 41152904075); shrink TableStep\_s where it matters}
\end{itemize}
\item[DynamicsSourceThermalTableReference] {\small\emph{reference \textperiodcentered\ \texttt{app.\allowbreak{}DynamicsSource}}} --- thermal.model transient on the native path (2026-10-30, the owner's ruling that standard goes transient and stays integrable): prepare integrates the walls ONCE along 'ThermalOrbit' with Model.accel, and the C reads the temperatures by time. On that orbit and at a table node (an orbit epoch) the native path and Model.accel of the SAME orbit evaluate the same temperatures, so they are again two implementations of one expression set; the thermal recoil is compared on its own through the diagnostics, so a parity carried by the drag cannot pass for one of the walls.
\begin{itemize}[leftmargin=1.2em,itemsep=0pt,topsep=1pt]\small
\item native transient table vs Model.accel at the table nodes, on the reference orbit: measured 8.9358e-16, bound 1e-13 --- \emph{the same parity as the seven configurations: the temperatures are the same numbers, the recoil the same expression}
\item native transient thermal recoil vs Model's at the table nodes: measured 1.5161e-16, bound 1e-13 rel --- \emph{kernel.thermalRecoil and its C port on identical temperatures}
\item native transient table vs Model.accel between the nodes (60 s tables): measured 1.1598e-05, bound 1e-2 --- \emph{the cubic reads the wall temperatures as smooth between nodes; recorded, not claimed (shrink TableStep\_s where it matters)}
\end{itemize}
\item[DynamicsSourceDensityScaleReference] {\small\emph{reference \textperiodcentered\ \texttt{app.\allowbreak{}DynamicsSource}}} --- drag.densityScale on the native path (2026-10-29, O-19): the daily table against Model.accel at the table nodes, for Swarm's bracket (interpolated in height in C), the driver scale (one factor per day) and a combined scale forced onto its fallback (MinRecords above any day's count). The factor itself is compared through the diagnostics, so a parity that holds only because the drag is small against the radiation cannot pass for one that holds on the scale.
\begin{itemize}[leftmargin=1.2em,itemsep=0pt,topsep=1pt]\small
\item native evaluator vs Model.accel at the table nodes, drag.densityScale swarm / driver / combined: measured 1.0185e-15, bound 1e-13 --- \emph{the same parity as the seven configurations: the scale is one multiplication after the density}
\item native density scale vs Model's: measured 0.0000e+00, bound 1e-14 rel --- \emph{Swarm's bracket: the same log-linear expression in the same order; driver days: the factor itself, computed by S.at}
\end{itemize}
\item[DynamicsSourcePrecision] {\small\emph{precision \textperiodcentered\ \texttt{app.\allowbreak{}DynamicsSource}}} --- The two PATHS of one source: the native entry and Model.accel, at a node, with everything else identical. And the gravity source is added, not folded in: with it the acceleration is the sum.
\begin{itemize}[leftmargin=1.2em,itemsep=0pt,topsep=1pt]\small
\item the source's two paths at a table node: measured 6.1849e-16, bound 1e-12 --- \emph{the C evaluator against Model.accel through the same source}
\item native density vs Model's: measured 0.0000e+00, bound 1e-13 rel --- \emph{the same backend called with the same drivers}
\item native relative velocity vs Model's: measured 0.0000e+00, bound 1e-12 --- \emph{the co-rotation from the same rotation tables}
\item native shadow factor vs Model's: measured 0.0000e+00, bound 1e-12 --- \emph{the same eclipse model}
\end{itemize}
\item[DynamicsSourceVectorization] {\small\emph{vectorization \textperiodcentered\ \texttt{app.\allowbreak{}DynamicsSource}}} --- An n-by-m block equals m single-column calls, BITWISE: the evaluator loops over columns and carries nothing between them, which is what makes the integrator's ensembles trustworthy.
\item[DynamicsSourceRobustness] {\small\emph{robustness \textperiodcentered\ \texttt{app.\allowbreak{}DynamicsSource}}}
\item[DynamicsSourceMetamorphicFlowTable] {\small\emph{metamorphic \textperiodcentered\ \texttt{app.\allowbreak{}DynamicsSource}}} --- The reference test's parity covers the flow-shadow table only if BOTH paths read it: the standard preset's native acceleration agrees with Model.accel at a node, and the exact flow shadow moves the result by far more than the two paths differ.
\begin{itemize}[leftmargin=1.2em,itemsep=0pt,topsep=1pt]\small
\item native vs Model.accel, both on the flow-shadow table, at a node: measured 5.6190e-16, bound 1e-13 --- \emph{the C trilinear is MATLAB's in the same order}
\end{itemize}
\item[DynamicsSourceMetamorphic] {\small\emph{metamorphic \textperiodcentered\ \texttt{app.\allowbreak{}DynamicsSource}}} --- 1. cdScale multiplies the gas part and nothing else, through the native evaluator.
\begin{itemize}[leftmargin=1.2em,itemsep=0pt,topsep=1pt]\small
\item native drag doubles with cdScale: measured 0.0000e+00, bound 1e-15 --- \emph{an exact power of two through a linear factor}
\item propagated arc: native binding vs vectorized binding over 300 s: measured 0.0000e+00, bound 1e-6 m --- \emph{one model through two stepping kernels; the acceleration differs by 1e-14 relative and the arc is 300 s (measured 0 m)}
\end{itemize}
\item[DynamicsSourceDensityScaleMetamorphic] {\small\emph{metamorphic \textperiodcentered\ \texttt{app.\allowbreak{}DynamicsSource}}} --- The gas force is linear in the densities, so the native drag with the driver scale is the unscaled native drag times the day's factor, and radiation and thermal do not notice it.
\begin{itemize}[leftmargin=1.2em,itemsep=0pt,topsep=1pt]\small
\item native drag with the driver scale / (factor x unscaled native drag): measured 3.9496e-16, bound 1e-14 --- \emph{a linear factor through the per-species force sums; fixed accommodation does not read the density}
\end{itemize}
\item[DynamicsSourceDensityScaleRobustness] {\small\emph{robustness \textperiodcentered\ \texttt{app.\allowbreak{}DynamicsSource}}} --- A Swarm-only scale on days without a factor prepares (the table marks them missing) and is REFUSED when a step reaches one, as Model.accel refuses it; the combined scale's fallback serves the same day.
\item[DynamicsSourceSpeed] {\small\emph{speed \textperiodcentered\ \texttt{app.\allowbreak{}DynamicsSource}}} --- The native entry exists to be cheap in an inner loop: one evaluation against the same source's MATLAB path, both timed inside one repeat on this host (rule T5, a ratio and never an absolute).
\begin{itemize}[leftmargin=1.2em,itemsep=0pt,topsep=1pt]\small
\item native accel / MATLAB accel, one state: measured 4.1108e-02, bound 0.08 --- \emph{measured 0.0234 (43x) on 2026-09-16, R2026a, standard preset; budget 3x the measured ratio}
\end{itemize}
\end{description}

\subsection{\texttt{tNgfIo}}
\noindent Tests of the io layer (LP-2): readACC1B, readSCA1B, readGNV1B, readTHR1B, readMAS1B, readMAG1B, igrfCoefficients, importCERES, massModel, spaceWeather, readSwarmDensity.

\begin{description}[leftmargin=1em,style=nextline,itemsep=2pt,font=\normalfont\ttfamily\small]
\item[importCERESContract] {\small\emph{contract \textperiodcentered\ \texttt{io.\allowbreak{}importCERES}}}
\item[igrfCoefficientsContract] {\small\emph{contract \textperiodcentered\ \texttt{io.\allowbreak{}igrfCoefficients}}}
\item[readMAG1BContract] {\small\emph{contract \textperiodcentered\ \texttt{io.\allowbreak{}readMAG1B}}}
\item[readACC1BContract] {\small\emph{contract \textperiodcentered\ \texttt{io.\allowbreak{}readACC1B}}}
\item[readSCA1BContract] {\small\emph{contract \textperiodcentered\ \texttt{io.\allowbreak{}readSCA1B}}}
\item[readGNV1BContract] {\small\emph{contract \textperiodcentered\ \texttt{io.\allowbreak{}readGNV1B}}}
\item[readTHR1BContract] {\small\emph{contract \textperiodcentered\ \texttt{io.\allowbreak{}readTHR1B}}}
\item[readMAS1BContract] {\small\emph{contract \textperiodcentered\ \texttt{io.\allowbreak{}readMAS1B}}}
\item[massModelContract] {\small\emph{contract \textperiodcentered\ \texttt{io.\allowbreak{}massModel}}}
\item[spaceWeatherContract] {\small\emph{contract \textperiodcentered\ \texttt{io.\allowbreak{}spaceWeather}}}
\item[importCERESReference] {\small\emph{reference \textperiodcentered\ \texttt{io.\allowbreak{}importCERES}}} --- Oracle: tools/pyxcheck/ceres\_check.py reads the same netCDF-3 files with its own parser (T9). Cell values at eight cells per field and month, chosen at the corners, the equator, both hemispheres and the longitude seam, pin the axis order and orientation: they must agree BITWISE (one float, printed with 17 digits). The area-weighted global means are the whole grid in one number each.
\begin{itemize}[leftmargin=1.2em,itemsep=0pt,topsep=1pt]\small
\item importCERES cells vs the independent netCDF-3 reader (48 values): measured 0.0000e+00, bound 0 W/m\textasciicircum{}2 --- \emph{the same float read twice: bitwise}
\item importCERES global means vs the oracle (6 means): measured 1.2879e-14, bound 1e-13 --- \emph{64800 products summed in another order; rounding only}
\end{itemize}
\item[igrfCoefficientsReference] {\small\emph{reference \textperiodcentered\ \texttt{io.\allowbreak{}igrfCoefficients}}} --- ORACLE IS THE PUBLISHED FILE, read a second way. These four numbers are transcribed from igrf14coeffs.txt by eye; if a column or a degree/order index moved, they are the first thing that fails.
\begin{itemize}[leftmargin=1.2em,itemsep=0pt,topsep=1pt]\small
\item g\_1\textasciicircum{}0 at 2020.0 against the published table: measured 0.0000e+00, bound 0 nT --- \emph{the axial dipole, read off the file}
\item g\_1\textasciicircum{}1 at 2020.0 against the published table: measured 0.0000e+00, bound 0 nT --- \emph{the equatorial dipole}
\item h\_1\textasciicircum{}1 at 2020.0 against the published table: measured 0.0000e+00, bound 0 nT --- \emph{h is a separate table and must not be confused with g}
\item g\_1\textasciicircum{}0 at 2025.0 against the published table: measured 0.0000e+00, bound 0 nT --- \emph{the most recent definitive epoch}
\item secular variation of g\_1\textasciicircum{}0: measured 0.0000e+00, bound 0 nT/yr --- \emph{positive: the axial dipole is weakening, since g\_1\textasciicircum{}0 is negative}
\item quadrupole over dipole at 2025.0: measured 1.7937e-01, bound 0.3 --- \emph{the Earth's field is dipole-dominated; a ratio near one would mean the degrees were transposed}
\end{itemize}
\item[readMAG1BReference] {\small\emph{reference \textperiodcentered\ \texttt{io.\allowbreak{}readMAG1B}}} --- ORACLE IS THE FILE, read independently: the first data line is parsed here by hand and compared against what the reader produced, so a column shifted by one cannot pass. The published header names MfvX/Y/Z\_RAW as columns 5 to 7 in microtesla and the six torquer currents as columns 8 to 13 in milliampere.
\begin{itemize}[leftmargin=1.2em,itemsep=0pt,topsep=1pt]\small
\item readMAG1B field vs the first data line, microtesla to tesla: measured 0.0000e+00, bound 0 T --- \emph{one multiplication by 1e-6; anything but zero means a column moved}
\item readMAG1B branch currents vs the first data line, mA to A: measured 0.0000e+00, bound 0 A --- \emph{one multiplication by 1e-3, six columns in the published order}
\item measured field magnitude over the excerpt: measured 3.6751e+01, bound 100 uT --- \emph{the Earth's field at 490 km reaches about 50 microtesla; a bound of 100 catches a misread column without asserting the geomagnetic model}
\end{itemize}
\item[readACC1BHybridReference] {\small\emph{reference \textperiodcentered\ \texttt{io.\allowbreak{}readACC1B}}} --- ACH1B, GRACE-D's hybrid transplant (N-10b): read as product ACH1B, and its linear accelerations are the sum of the four .ACX components ACU + ACS + ACM + AC0 (CONCEPT V-16), read here with gracefoReading's generic reader -- an independent path. 120 real records of 2020-01-01 in which the thruster model is non-zero; the files print 16 significant digits, so the sum agrees to a few units of that rounding.
\begin{itemize}[leftmargin=1.2em,itemsep=0pt,topsep=1pt]\small
\item ACH1B = ACU + ACS + ACM + AC0 (120 real records), relative to the largest value: measured 7.5155e-16, bound 4 * 5e-16 + 4 * eps --- \emph{four 16-digit printed terms and their sum: 4 x 5e-16 printing plus 4 eps of summation}
\end{itemize}
\item[readACC1BReference] {\small\emph{reference \textperiodcentered\ \texttt{io.\allowbreak{}readACC1B}}} --- Oracle: the file text itself, parsed independently (regexp tokens of the first data record).
\item[readSCA1BReference] {\small\emph{reference \textperiodcentered\ \texttt{io.\allowbreak{}readSCA1B}}} --- Physics oracle: the body +Z axis (graceFO "Nadir" plate normal) rotated into the GCRS must point to nadir (-r/\textbar{}r\textbar{} from GNV1B). Measured 2026-09-06: 0.84 deg mean, 0.86 deg max when conjugated; 96 deg mean when NOT conjugated. Bound 3 deg: the GRACE-FO pitch offset and attitude deadband are \textasciitilde{}1 deg.
\begin{itemize}[leftmargin=1.2em,itemsep=0pt,topsep=1pt]\small
\item SCA1B body +Z vs nadir (conjugated): measured 8.6018e-01, bound 3 deg --- \emph{measured 0.86; 96 without the conjugation}
\end{itemize}
\item[readGNV1BReference] {\small\emph{reference \textperiodcentered\ \texttt{io.\allowbreak{}readGNV1B}}} --- Oracle: Greenwich mean sidereal time from the IAU 1982 polynomial (Aoki et al. 1982; Meeus 12.4), evaluated here without referenceSystem. The GCRS-ITRS longitude difference of a point equals the hour angle of the equinox at Greenwich up to precession-nutation since J2000 (\textasciitilde{}0.28 deg in 2020) and the CIO/equinox difference. Measured 2026-09-06: -0.19 deg after the 18 s GPS-UTC offset is applied; bound 0.5 deg. A wrong rotation sense gives \textasciitilde{}160 deg.
\begin{itemize}[leftmargin=1.2em,itemsep=0pt,topsep=1pt]\small
\item GNV1B GCRS-ITRS longitude vs IAU1982 GMST: measured 1.9067e-01, bound 0.5 deg --- \emph{precession since J2000 dominates (measured 0.19)}
\end{itemize}
\item[readTHR1BReference] {\small\emph{reference \textperiodcentered\ \texttt{io.\allowbreak{}readTHR1B}}} --- Oracle: the raw first record (time\_intg time\_frac time\_ref id 14 counts 14 on\_times ...).
\item[readMAS1BReference] {\small\emph{reference \textperiodcentered\ \texttt{io.\allowbreak{}readMAS1B}}} --- Oracle: the values printed in the first record of the real file (also pinned by gracefoReading's realdata test).
\item[massModelReference] {\small\emph{reference \textperiodcentered\ \texttt{io.\allowbreak{}massModel}}} --- Oracle: JPL D-56935 masses (GRACE-C 601.214 kg, GRACE-D 601.209 kg).
\item[spaceWeatherReference] {\small\emph{reference \textperiodcentered\ \texttt{io.\allowbreak{}spaceWeather}}} --- Oracle: the raw CSV / TXT lines parsed here by hand (split on comma / whitespace).
\item[importCERESPrecision] {\small\emph{precision \textperiodcentered\ \texttt{io.\allowbreak{}importCERES}}} --- The month nodes are exact: the file says 7228.25 and 7259 days after 2000-03-01 00:00 UTC, which are 2019-12-15 06:00 and 2020-01-15 00:00, MJD 58832.25 and 58863 (the CF time units parsed, not assumed).
\item[igrfCoefficientsPrecision] {\small\emph{precision \textperiodcentered\ \texttt{io.\allowbreak{}igrfCoefficients}}} --- Interpolation is one linear blend, so the two ends must be EXACT and the midpoint must be the mean. Anything else means the weight or the bracketing index is wrong.
\begin{itemize}[leftmargin=1.2em,itemsep=0pt,topsep=1pt]\small
\item interpolating to a published epoch returns that epoch: measured 0.0000e+00, bound 0 nT --- \emph{a weight of zero must reproduce the node, not approximate it}
\item the midpoint is the mean of the two nodes: measured 0.0000e+00, bound 1e-11 nT --- \emph{one linear blend at w = 0.5; the bound is the rounding of numbers of order 3e4}
\item extrapolation is the last epoch plus SV times dt: measured 0.0000e+00, bound 1e-11 nT --- \emph{two years of secular variation, applied once}
\end{itemize}
\item[readMAG1BPrecision] {\small\emph{precision \textperiodcentered\ \texttt{io.\allowbreak{}readMAG1B}}} --- The net current is the sum of two branches, and the sum is only meaningful because the branches never overlap. That is asserted here rather than trusted: it was measured over the full day and it must hold on whatever excerpt is present.
\begin{itemize}[leftmargin=1.2em,itemsep=0pt,topsep=1pt]\small
\item torquer net equals A plus B exactly: measured 0.0000e+00, bound 0 A --- \emph{one addition, quoted back}
\end{itemize}
\item[readACC1BPrecision] {\small\emph{precision \textperiodcentered\ \texttt{io.\allowbreak{}readACC1B}}} --- Time-tag round trip is exact: epoch -\textgreater{} [day sec] equals gracefoReading's time bitwise (bound 0).
\begin{itemize}[leftmargin=1.2em,itemsep=0pt,topsep=1pt]\small
\item ACC1B time tag round trip: measured 0.0000e+00, bound 0 s --- \emph{exact by construction}
\end{itemize}
\item[readSCA1BPrecision] {\small\emph{precision \textperiodcentered\ \texttt{io.\allowbreak{}readSCA1B}}} --- invertQuaternion normalises: \textbar{}norm(q) - 1\textbar{} at the eps level (measured 2.2e-16); bound 1e-15 = 4.5 eps.
\begin{itemize}[leftmargin=1.2em,itemsep=0pt,topsep=1pt]\small
\item SCA1B unit norm of q\_body2gcrs: measured 2.2204e-16, bound 1e-15 --- \emph{measured 2.2e-16}
\end{itemize}
\item[readGNV1BPrecision] {\small\emph{precision \textperiodcentered\ \texttt{io.\allowbreak{}readGNV1B}}} --- The rotated velocity must contain the frame-rate term: a five-point central difference of the GCRS positions (1 Hz) reproduces it to 4.4e-5 m/s (measured; dominated by the GNV1B position noise \textasciitilde{}mm and the stencil truncation at 1 s). Bound 1e-3 m/s: 20x the measurement, 5 orders below the 470 m/s error of a dropped Rdot term.
\begin{itemize}[leftmargin=1.2em,itemsep=0pt,topsep=1pt]\small
\item GNV1B velocity vs 5-point FD of GCRS position: measured 4.3800e-05, bound 1e-3 m/s --- \emph{measured 4.4e-5; 4.7e2 without Rdot}
\end{itemize}
\item[readTHR1BPrecision] {\small\emph{precision \textperiodcentered\ \texttt{io.\allowbreak{}readTHR1B}}} --- Event duration equals the longest on-time of the record exactly (integer ms / 1000, one division).
\begin{itemize}[leftmargin=1.2em,itemsep=0pt,topsep=1pt]\small
\item THR1B event duration vs max on-time: measured 0.0000e+00, bound 1e-6 s --- \emph{datetime resolution \textasciitilde{}1e-7 s near 2020}
\end{itemize}
\item[readMAS1BPrecision] {\small\emph{precision \textperiodcentered\ \texttt{io.\allowbreak{}readMAS1B}}} --- total is the plain sum tank1 + tank2 (same operands, same order): bitwise.
\begin{itemize}[leftmargin=1.2em,itemsep=0pt,topsep=1pt]\small
\item MAS1B total vs tank1+tank2: measured 0.0000e+00, bound 0 kg --- \emph{bitwise}
\end{itemize}
\item[massModelPrecision] {\small\emph{precision \textperiodcentered\ \texttt{io.\allowbreak{}massModel}}} --- At the MAS1B sample epochs the model equals anchor + gas - gas(1): interp1 at a node is exact, the two additions round at 1e-13 kg (600 kg * eps); bound 1e-9 kg leaves 4 decades.
\begin{itemize}[leftmargin=1.2em,itemsep=0pt,topsep=1pt]\small
\item massModel mas1b at the sample epochs: measured 0.0000e+00, bound 1e-9 kg --- \emph{rounding 1e-13}
\end{itemize}
\item[spaceWeatherPrecision] {\small\emph{precision \textperiodcentered\ \texttt{io.\allowbreak{}spaceWeather}}} --- kp is the file integer / 10 (the only arithmetic): 10*kp is an integer to 4 eps.
\begin{itemize}[leftmargin=1.2em,itemsep=0pt,topsep=1pt]\small
\item spaceWeather 10*kp integer: measured 0.0000e+00, bound 1e-14 --- \emph{one division}
\end{itemize}
\item[importCERESVectorization] {\small\emph{vectorization \textperiodcentered\ \texttt{io.\allowbreak{}importCERES}}} --- A period that sits exactly on one node loads that month alone; each such single-month import is the corresponding page of the whole import, bitwise.
\item[igrfCoefficientsVectorization] {\small\emph{vectorization \textperiodcentered\ \texttt{io.\allowbreak{}igrfCoefficients}}} --- The tables are read once and indexed; asking for many epochs must give the same numbers as asking for each, because there is no state between calls.
\begin{itemize}[leftmargin=1.2em,itemsep=0pt,topsep=1pt]\small
\item the tables are identical whatever epoch is requested: measured 0.0000e+00, bound 0 nT --- \emph{interpolation reads the tables; it may not modify them}
\end{itemize}
\item[readMAG1BVectorization] {\small\emph{vectorization \textperiodcentered\ \texttt{io.\allowbreak{}readMAG1B}}} --- A TimeRange is a filter, not a different reader: the selected records must be the same numbers the full read produced, at the same epochs.
\begin{itemize}[leftmargin=1.2em,itemsep=0pt,topsep=1pt]\small
\item a TimeRange selects records without changing them: measured 0.0000e+00, bound 0 T --- \emph{a filter that alters a value is not a filter}
\end{itemize}
\item[readACC1BVectorization] {\small\emph{vectorization \textperiodcentered\ \texttt{io.\allowbreak{}readACC1B}}} --- A TimeRange excerpt equals the slice of the full read, bitwise (leading and interior).
\item[readSCA1BVectorization] {\small\emph{vectorization \textperiodcentered\ \texttt{io.\allowbreak{}readSCA1B}}}
\item[readGNV1BVectorization] {\small\emph{vectorization \textperiodcentered\ \texttt{io.\allowbreak{}readGNV1B}}} --- The rotation of a batch must equal the rotation of a subset, bitwise (measured 0 difference).
\item[readTHR1BVectorization] {\small\emph{vectorization \textperiodcentered\ \texttt{io.\allowbreak{}readTHR1B}}} --- Whole day equals the two half days concatenated (isequaln: amplitude and sigma are NaN by design).
\item[readMAS1BVectorization] {\small\emph{vectorization \textperiodcentered\ \texttt{io.\allowbreak{}readMAS1B}}}
\item[massModelVectorization] {\small\emph{vectorization \textperiodcentered\ \texttt{io.\allowbreak{}massModel}}}
\item[spaceWeatherVectorization] {\small\emph{vectorization \textperiodcentered\ \texttt{io.\allowbreak{}spaceWeather}}}
\item[importCERESRobustness] {\small\emph{robustness \textperiodcentered\ \texttt{io.\allowbreak{}importCERES}}}
\item[igrfCoefficientsRobustness] {\small\emph{robustness \textperiodcentered\ \texttt{io.\allowbreak{}igrfCoefficients}}} --- A truncated row must be refused with the mnemonic that names the cause.
\item[readMAG1BRobustness] {\small\emph{robustness \textperiodcentered\ \texttt{io.\allowbreak{}readMAG1B}}} --- A corrupted copy must be refused with the mnemonic that names the cause, not with whatever MATLAB happens to throw first.
\item[readACC1BRobustness] {\small\emph{robustness \textperiodcentered\ \texttt{io.\allowbreak{}readACC1B}}} --- A quality mask naming a bit no record sets removes nothing; two files of different products refuse (foreign id).
\item[readACC1BMultiEmptyRobustness] {\small\emph{robustness \textperiodcentered\ \texttt{io.\allowbreak{}readACC1B}}} --- O-12 discharge. Until gracefoReading 1.0.1, l1bLoad sent one file through readL1b and several through readL1bMulti, because the multi-file reader raised MATLAB:badsubscript on an all-empty TimeRange. dataEmpty was therefore guaranteed for one file and UNREACHABLE for several. The split is gone; this asserts the path that could not be tested before. The same fixture twice is enough: the range excludes every record, so the concatenation is empty and no ordering check is reached.
\item[readSCA1BRobustness] {\small\emph{robustness \textperiodcentered\ \texttt{io.\allowbreak{}readSCA1B}}} --- A quaternion off the unit sphere is refused, not repaired.
\item[readGNV1BRobustness] {\small\emph{robustness \textperiodcentered\ \texttt{io.\allowbreak{}readGNV1B}}} --- coord\_ref other than 'E' is refused; the whole 1962-now C04 series propagates referenceSystem's refusal unwrapped.
\item[readTHR1BRobustness] {\small\emph{robustness \textperiodcentered\ \texttt{io.\allowbreak{}readTHR1B}}}
\item[readMAS1BRobustness] {\small\emph{robustness \textperiodcentered\ \texttt{io.\allowbreak{}readMAS1B}}} --- 150 kg in a tank column is not propellant: refused.
\item[massModelRobustness] {\small\emph{robustness \textperiodcentered\ \texttt{io.\allowbreak{}massModel}}}
\item[spaceWeatherRobustness] {\small\emph{robustness \textperiodcentered\ \texttt{io.\allowbreak{}spaceWeather}}}
\item[importCERESMetamorphic] {\small\emph{metamorphic \textperiodcentered\ \texttt{io.\allowbreak{}importCERES}}} --- Any period strictly inside the two nodes needs both and returns the whole import unchanged; the three fields are independent, so the rsut of one import is the rsut of any other.
\item[igrfCoefficientsMetamorphic] {\small\emph{metamorphic \textperiodcentered\ \texttt{io.\allowbreak{}igrfCoefficients}}}
\begin{itemize}[leftmargin=1.2em,itemsep=0pt,topsep=1pt]\small
\item a quarter step moves a quarter of the difference: measured 9.0949e-13, bound 1e-11 nT --- \emph{linearity in the weight, checked away from the midpoint where a factor of two would hide}
\end{itemize}
\item[readMAG1BMetamorphic] {\small\emph{metamorphic \textperiodcentered\ \texttt{io.\allowbreak{}readMAG1B}}}
\begin{itemize}[leftmargin=1.2em,itemsep=0pt,topsep=1pt]\small
\item magnitude equals the norm of the components: measured 0.0000e+00, bound 0 T --- \emph{one vecnorm, quoted back}
\item microtesla to tesla is a pure scaling: measured 2.1176e-16, bound 1e-15 --- \emph{every non-zero component divided by its file value gives the same factor}
\item calibration columns are constant over this excerpt: measured 0.0000e+00, bound 0 --- \emph{measured [0, 65000, -65000] for all 172800 records of the full day; recorded because the values look like sentinels, not scale factors}
\end{itemize}
\item[readACC1BMetamorphic] {\small\emph{metamorphic \textperiodcentered\ \texttt{io.\allowbreak{}readACC1B}}} --- Split/merge: two adjoining TimeRanges concatenated equal the whole read, bitwise.
\item[readSCA1BMetamorphic] {\small\emph{metamorphic \textperiodcentered\ \texttt{io.\allowbreak{}readSCA1B}}} --- Inverting the output once more reproduces the (normalised) file quaternion at the eps level (two normalisations; measured 3.3e-16). Bound 1e-14.
\begin{itemize}[leftmargin=1.2em,itemsep=0pt,topsep=1pt]\small
\item SCA1B double conjugation vs file: measured 3.3307e-16, bound 1e-14 --- \emph{measured 3.3e-16}
\end{itemize}
\item[readGNV1BMetamorphic] {\small\emph{metamorphic \textperiodcentered\ \texttt{io.\allowbreak{}readGNV1B}}} --- (1) a rotation preserves \textbar{}r\textbar{} (measured 4e-16 relative; bound 1e-14); (2) round trip GCRS -\textgreater{} ITRS through applyRotation's inverse reproduces the file (measured 5e-9 m; bound 1e-6 m, 1000x the double rounding at 7e6 m and 1000x below the HANDOVER requirement of 1 mm); (3) an Eop struct and the same series as a file give bitwise the same state.
\begin{itemize}[leftmargin=1.2em,itemsep=0pt,topsep=1pt]\small
\item GNV1B \textbar{}r\textbar{} GCRS vs ITRS: measured 5.4082e-16, bound 1e-14 rel --- \emph{measured 4e-16}
\item GNV1B round trip GCRS-\textgreater{}ITRS position: measured 4.3930e-09, bound 1e-6 m --- \emph{measured 5e-9; requirement 1e-3}
\item GNV1B round trip GCRS-\textgreater{}ITRS velocity: measured 1.2974e-11, bound 1e-9 m/s --- \emph{measured 1e-11}
\end{itemize}
\item[readTHR1BMetamorphic] {\small\emph{metamorphic \textperiodcentered\ \texttt{io.\allowbreak{}readTHR1B}}} --- The event table and aux describe the same records: one event per fired record, in order, and the UTC epochs are the GPS epochs minus the leap-second count.
\item[readMAS1BMetamorphic] {\small\emph{metamorphic \textperiodcentered\ \texttt{io.\allowbreak{}readMAS1B}}} --- Split/merge over TimeRange equals the whole read (isequaln: NaN columns by design).
\item[massModelMetamorphic] {\small\emph{metamorphic \textperiodcentered\ \texttt{io.\allowbreak{}massModel}}} --- (1) shifting the anchor by D shifts every mass by D (one addition: 1e-12 kg); (2) the epoch scale of the query does not matter (converted inside): bitwise-equal day/sec after the 18 s GPS-UTC shift, so identical masses to 1e-12 kg.
\begin{itemize}[leftmargin=1.2em,itemsep=0pt,topsep=1pt]\small
\item massModel anchor shift: measured 0.0000e+00, bound 1e-12 kg --- \emph{one addition at 600 kg}
\item massModel query scale invariance: measured 0.0000e+00, bound 1e-12 kg --- \emph{exact 18 s shift}
\end{itemize}
\item[spaceWeatherMetamorphic] {\small\emph{metamorphic \textperiodcentered\ \texttt{io.\allowbreak{}spaceWeather}}} --- Daily indices are step functions: 00:00:00 and 23:59:59 of one day agree on every daily field; the JB2008 lags are relations between query and table: query(d).f10 == table(d-1), etc.
\item[importCERESSpeed] {\small\emph{speed \textperiodcentered\ \texttt{io.\allowbreak{}importCERES}}} --- Everything importCERES adds to three ncread calls (attribute checks, axis comparisons, the permutation and the fill screen) against those bare reads: measured see last\_run; budget 3x the 2026-09-15 figure.
\begin{itemize}[leftmargin=1.2em,itemsep=0pt,topsep=1pt]\small
\item importCERES vs three bare ncread calls: measured 7.9442e+00, bound 25 --- \emph{checks and one permutation on 1.2e5 floats per field; measured 8.3 on 2026-09-15, budget 3x}
\end{itemize}
\item[igrfCoefficientsSpeed] {\small\emph{speed \textperiodcentered\ \texttt{io.\allowbreak{}igrfCoefficients}}}
\begin{itemize}[leftmargin=1.2em,itemsep=0pt,topsep=1pt]\small
\item igrfCoefficients, 195 rows x 27 columns: measured 4.2379e+01, bound 2000 ms --- \emph{a smoke ceiling on a 42 kB parse, not a claim about this host}
\end{itemize}
\item[readMAG1BSpeed] {\small\emph{speed \textperiodcentered\ \texttt{io.\allowbreak{}readMAG1B}}} --- Reading is dominated by gracefoReading's parse; what this checks is that nothing quadratic sits in the marshalling above it.
\begin{itemize}[leftmargin=1.2em,itemsep=0pt,topsep=1pt]\small
\item readMAG1B, 600 records: measured 9.3114e+01, bound 5000 ms --- \emph{a smoke ceiling on a 600-record parse, not a claim about this host}
\end{itemize}
\item[readACC1BSpeed] {\small\emph{speed \textperiodcentered\ \texttt{io.\allowbreak{}readACC1B}}} --- Reader over raw gracefoReading.readL1b: measured 0.98 (median of 5, warm). Bound 5 = 3x a generous 1.5 expectation for the epoch struct, checks and provenance.
\begin{itemize}[leftmargin=1.2em,itemsep=0pt,topsep=1pt]\small
\item readACC1B / readL1b time ratio: measured 1.0130e+00, bound 5 --- \emph{measured 0.98}
\end{itemize}
\item[readSCA1BSpeed] {\small\emph{speed \textperiodcentered\ \texttt{io.\allowbreak{}readSCA1B}}} --- Measured 1.46 (norm check + conjugation). Bound 5.
\begin{itemize}[leftmargin=1.2em,itemsep=0pt,topsep=1pt]\small
\item readSCA1B / readL1b time ratio: measured 1.0524e+00, bound 5 --- \emph{measured 1.46}
\end{itemize}
\item[readGNV1BSpeed] {\small\emph{speed \textperiodcentered\ \texttt{io.\allowbreak{}readGNV1B}}} --- read over ITRS read (the rotation with Rdot for 600 epochs): measured 1.16 warm, 6.3 cold. Bound 5 after warm-up (3x the 1.5 expectation).
\begin{itemize}[leftmargin=1.2em,itemsep=0pt,topsep=1pt]\small
\item readGNV1B GCRS / ITRS time ratio: measured 1.2288e+00, bound 5 --- \emph{measured 1.16 warm}
\end{itemize}
\item[readTHR1BSpeed] {\small\emph{speed \textperiodcentered\ \texttt{io.\allowbreak{}readTHR1B}}} --- Measured 1.04 (event table build is cheap). Bound 5.
\begin{itemize}[leftmargin=1.2em,itemsep=0pt,topsep=1pt]\small
\item readTHR1B / readL1b time ratio: measured 1.0631e+00, bound 5 --- \emph{measured 1.04}
\end{itemize}
\item[readMAS1BSpeed] {\small\emph{speed \textperiodcentered\ \texttt{io.\allowbreak{}readMAS1B}}} --- Measured 1.47. Bound 5.
\begin{itemize}[leftmargin=1.2em,itemsep=0pt,topsep=1pt]\small
\item readMAS1B / readL1b time ratio: measured 1.2444e+00, bound 5 --- \emph{measured 1.47}
\end{itemize}
\item[massModelSpeed] {\small\emph{speed \textperiodcentered\ \texttt{io.\allowbreak{}massModel}}} --- at() over one day at 1 Hz against a bare interp1 of the same size: measured 2.24 (validation, clamping, the range check). Bound 7 = 3x.
\begin{itemize}[leftmargin=1.2em,itemsep=0pt,topsep=1pt]\small
\item massModel.at(86400) / interp1 time ratio: measured 2.7048e+00, bound 7 --- \emph{measured 2.24}
\end{itemize}
\item[spaceWeatherSpeed] {\small\emph{speed \textperiodcentered\ \texttt{io.\allowbreak{}spaceWeather}}} --- query of one day at 1 Hz (86400 epochs, 0.05 s measured) against an ismember day lookup of the same size: measured 65 (four day rows, 20 history picks, five NaN scans). Bound 200 = 3x.
\begin{itemize}[leftmargin=1.2em,itemsep=0pt,topsep=1pt]\small
\item spaceWeather.query(86400) / ismember time ratio: measured 5.5414e+01, bound 200 --- \emph{measured 65}
\end{itemize}
\item[importGroopsGridContract] {\small\emph{contract \textperiodcentered\ \texttt{io.\allowbreak{}importGroopsGrid}}}
\item[importGroopsGridReference] {\small\emph{reference \textperiodcentered\ \texttt{io.\allowbreak{}importGroopsGrid}}} --- The oracle is the format itself, asserted by the three facts that make a misreading visible; each is checked on a file whose contents are known independently.
\begin{itemize}[leftmargin=1.2em,itemsep=0pt,topsep=1pt]\small
\item GROOPS grid: solid angles sum to 4 pi: measured 6.6613e-16, bound 1e-12 --- \emph{the writer lays a cosine-weighted grid on the sphere; the reader must not rescale it}
\end{itemize}
\item[importGroopsGridPrecision] {\small\emph{precision \textperiodcentered\ \texttt{io.\allowbreak{}importGroopsGrid}}} --- float64 in, float64 out: a read that rescales or single-casts shows here and nowhere else.
\begin{itemize}[leftmargin=1.2em,itemsep=0pt,topsep=1pt]\small
\item GROOPS grid values survive the round trip: measured 0.0000e+00, bound 0 --- \emph{float64 written and read unchanged: bitwise or the reader is lossy}
\end{itemize}
\item[importGroopsGridVectorization] {\small\emph{vectorization \textperiodcentered\ \texttt{io.\allowbreak{}importGroopsGrid}}} --- Reading a subrange must give exactly what reading everything gives for those grids: the seek arithmetic is the only difference and it must change no number.
\item[importGroopsGridRobustness] {\small\emph{robustness \textperiodcentered\ \texttt{io.\allowbreak{}importGroopsGrid}}}
\item[importGroopsGridMetamorphic] {\small\emph{metamorphic \textperiodcentered\ \texttt{io.\allowbreak{}importGroopsGrid}}} --- Writing what was read must reproduce the file: the layout is then self-consistent, which no single-direction test can establish.
\begin{itemize}[leftmargin=1.2em,itemsep=0pt,topsep=1pt]\small
\item GROOPS grid: read-write-read is a fixed point: measured 0.0000e+00, bound 0 --- \emph{bitwise; the two files describe the same grid}
\end{itemize}
\item[importGroopsGridSpeed] {\small\emph{speed \textperiodcentered\ \texttt{io.\allowbreak{}importGroopsGrid}}} --- IS THE POINT. One month of a many-month series must not cost what the whole series costs, or a 1.25 GB file is unusable for a one-day arc.
\begin{itemize}[leftmargin=1.2em,itemsep=0pt,topsep=1pt]\small
\item GROOPS grid: one grid of 60 / one grid of 2: measured 1.0304e+00, bound 3 --- \emph{seeking means the cost of one grid does not grow with the file: 30x the payload must not cost 30x. A reader that ignored TimeRange would land near 30}
\end{itemize}
\item[spaceWeatherDstF30Contract] {\small\emph{contract \textperiodcentered\ \texttt{io.\allowbreak{}spaceWeather}}}
\item[spaceWeatherDstF30Reference] {\small\emph{reference \textperiodcentered\ \texttt{io.\allowbreak{}spaceWeather}}}
\item[spaceWeatherF30Precision] {\small\emph{precision \textperiodcentered\ \texttt{io.\allowbreak{}spaceWeather}}} --- The 81-day mean is the only arithmetic of the f30 family: one sum of 81 values of \textasciitilde{}46 sfu and one division. Its rounding is bounded by 81 eps relative; measured against exact rationals.
\begin{itemize}[leftmargin=1.2em,itemsep=0pt,topsep=1pt]\small
\item F30 81-day mean, relative error vs exact rational: measured 1.1102e-16, bound 81 * eps --- \emph{a sum of 81 doubles and one division; the rationals are from an independent Python Fraction sum}
\end{itemize}
\item[spaceWeatherDstF30Vectorization] {\small\emph{vectorization \textperiodcentered\ \texttt{io.\allowbreak{}spaceWeather}}}
\item[spaceWeatherDstF30Robustness] {\small\emph{robustness \textperiodcentered\ \texttt{io.\allowbreak{}spaceWeather}}}
\item[spaceWeatherDstF30Metamorphic] {\small\emph{metamorphic \textperiodcentered\ \texttt{io.\allowbreak{}spaceWeather}}}
\item[readSwarmDensityContract] {\small\emph{contract \textperiodcentered\ \texttt{io.\allowbreak{}readSwarmDensity}}}
\item[readSwarmDensityToleosContract] {\small\emph{contract \textperiodcentered\ \texttt{io.\allowbreak{}readSwarmDensity}}} --- CHAMP, GRACE-A and GRACE-FO 1 accelerometer densities in the same layout, 10 s, six hours each
\item[readSwarmDensityReference] {\small\emph{reference \textperiodcentered\ \texttt{io.\allowbreak{}readSwarmDensity}}} --- Oracle: the raw CSV lines, split on commas with str2double (not textscan), and the calendar.
\item[readSwarmDensityToleosReference] {\small\emph{reference \textperiodcentered\ \texttt{io.\allowbreak{}readSwarmDensity}}} --- Oracle: the raw lines split with str2double; the fill 9.99e+32 of CHAMP's orbit mean is NaN, the rest unchanged.
\item[readSwarmDensityPrecision] {\small\emph{precision \textperiodcentered\ \texttt{io.\allowbreak{}readSwarmDensity}}} --- deg2rad is one multiplication: rad2deg brings every printed degree back to a few ulp.
\begin{itemize}[leftmargin=1.2em,itemsep=0pt,topsep=1pt]\small
\item readSwarmDensity degrees round trip / 180: measured 1.5790e-16, bound 4 * eps --- \emph{deg2rad then rad2deg}
\end{itemize}
\item[readSwarmDensityVectorization] {\small\emph{vectorization \textperiodcentered\ \texttt{io.\allowbreak{}readSwarmDensity}}} --- The three days equal the 1.5-day halves read separately, concatenated (bitwise: nothing is computed across records).
\item[readSwarmDensityRobustness] {\small\emph{robustness \textperiodcentered\ \texttt{io.\allowbreak{}readSwarmDensity}}} --- the product's own units, status and flags are read, not assumed
\item[readSwarmDensityToleosRobustness] {\small\emph{robustness \textperiodcentered\ \texttt{io.\allowbreak{}readSwarmDensity}}} --- CHAMP's orbit mean carries the fill on 359 of these records (validity\_flag\_orbitmean 1): NaN, density kept
\item[readSwarmDensityMetamorphic] {\small\emph{metamorphic \textperiodcentered\ \texttt{io.\allowbreak{}readSwarmDensity}}} --- Physics the product must show: Swarm A and C fly side by side at the same height (same density), Swarm B \textasciitilde{}65 km higher (several times thinner); an epoch present twice is kept once.
\begin{itemize}[leftmargin=1.2em,itemsep=0pt,topsep=1pt]\small
\item Swarm A / C median density ratio, \textbar{}log\textbar{}: measured 6.5472e-03, bound log(1.2) --- \emph{side-by-side pair}
\end{itemize}
\item[readSwarmDensityToleosMetamorphic] {\small\emph{metamorphic \textperiodcentered\ \texttt{io.\allowbreak{}readSwarmDensity}}} --- Physics the products must show on 2008-01-01: CHAMP \textasciitilde{}130 km below GRACE-A flies in several times the density
\item[readSwarmDensitySpeed] {\small\emph{speed \textperiodcentered\ \texttt{io.\allowbreak{}readSwarmDensity}}} --- Against reading the same file's text and splitting it into lines. Bound 20, an estimate on the first run (header decode, one textscan); the measured ratio is in tests/last\_run.txt.
\begin{itemize}[leftmargin=1.2em,itemsep=0pt,topsep=1pt]\small
\item readSwarmDensity / fileread+splitlines time ratio: measured 4.0198e+00, bound 20 --- \emph{first run}
\end{itemize}
\end{description}

\subsection{\texttt{tNgfKernel}}
\noindent Tests of the LP-5 kernel units: gsiCoefficients, accommodation, dragPanelForce, radiationPanelForce, thermalRecoil, assemble, bodyTorques.

\begin{description}[leftmargin=1em,style=nextline,itemsep=2pt,font=\normalfont\ttfamily\small]
\item[partialsContract] {\small\emph{contract \textperiodcentered\ \texttt{kernel.\allowbreak{}partials}}}
\item[gsiCoefficientsContract] {\small\emph{contract \textperiodcentered\ \texttt{kernel.\allowbreak{}gsiCoefficients}}}
\item[accommodationContract] {\small\emph{contract \textperiodcentered\ \texttt{kernel.\allowbreak{}accommodation}}}
\item[dragPanelForceContract] {\small\emph{contract \textperiodcentered\ \texttt{kernel.\allowbreak{}dragPanelForce}}}
\item[radiationPanelForceContract] {\small\emph{contract \textperiodcentered\ \texttt{kernel.\allowbreak{}radiationPanelForce}}}
\item[thermalRecoilContract] {\small\emph{contract \textperiodcentered\ \texttt{kernel.\allowbreak{}thermalRecoil}}}
\item[assembleContract] {\small\emph{contract \textperiodcentered\ \texttt{kernel.\allowbreak{}assemble}}}
\item[bodyTorquesContract] {\small\emph{contract \textperiodcentered\ \texttt{kernel.\allowbreak{}bodyTorques}}}
\item[partialsReference] {\small\emph{reference \textperiodcentered\ \texttt{kernel.\allowbreak{}partials}}} --- Richardson-extrapolated central differences of the force itself: an independent numerical route to the same derivative. A complex step is impossible here and the toolbox says so rather than pretending: MATLAB's erf and erfc raise notFullReal on a complex argument and abs/vecnorm discard the imaginary part (measured 2026-09-08).
\begin{itemize}[leftmargin=1.2em,itemsep=0pt,topsep=1pt]\small
\item d(drag)/d(alpha) vs Richardson central differences: measured 5.9144e-13, bound 1e-9 --- \emph{the extrapolation's own truncation dominates; measured 1e-12}
\item d(drag)/d(speed) vs Richardson central differences: measured 1.9412e-10, bound 1e-9 --- \emph{measured 7e-11}
\item d(drag)/d(densityScale) vs Richardson: measured 2.8201e-11, bound 1e-9 --- \emph{measured 2e-11}
\item d(radiation)/d(pressureScale) vs Richardson: measured 1.0756e-11, bound 1e-9 --- \emph{measured 8e-11}
\item d(thermal)/d(temperature) vs Richardson: measured 2.2264e-11, bound 1e-9 --- \emph{measured 2e-11}
\end{itemize}
\item[gsiCoefficientsReference] {\small\emph{reference \textperiodcentered\ \texttt{kernel.\allowbreak{}gsiCoefficients}}} --- kinetic-theory integration (Simpson, no erf) vs the closed forms; oracle error \textless{} 1e-10
\begin{itemize}[leftmargin=1.2em,itemsep=0pt,topsep=1pt]\small
\item sentman Cp vs kinetic-theory oracle (105 cases): measured 3.9524e-14, bound 1e-12 --- \emph{oracle Simpson error \textasciitilde{}1e-13, coefficients up to 7.8; measured 4e-14}
\item sentman Ctau vs oracle: measured 1.2434e-14, bound 1e-12 --- \emph{measured 1.2e-14}
\item drag projection Cp g + Ctau l vs oracle CD: measured 3.9524e-14, bound 1e-12 --- \emph{measured 4e-14}
\item lift projection Cp l - Ctau g vs oracle CL: measured 9.8810e-15, bound 1e-12 --- \emph{measured 1e-14}
\item specular (maxwell sigma = 1) Cp vs oracle: measured 5.1514e-14, bound 1e-12 --- \emph{measured 5e-14}
\item Cp g + Ctau l vs Doornbos C\_D closed form: measured 8.8818e-16, bound 1e-14 --- \emph{algebraic identity, rounding only}
\end{itemize}
\item[accommodationReference] {\small\emph{reference \textperiodcentered\ \texttt{kernel.\allowbreak{}accommodation}}} --- Oracle: tools/pyxcheck/accommodation\_check.py, a standard-library evaluation of BOTH published Pilinski fits (2010 paper eq. 8, K = 7.50e-17 m\textasciicircum{}3/K; 2011 thesis eq. 7.1, K = 4.98e-17) and of Paul et al. (2022) eqs. 10-13, written from the source texts by an agent that did not open this toolbox (T9), over 12 thermosphere cases from 1e12 to 1e17 m\textasciicircum{}-3 of atomic oxygen and both substrate masses in circulation. Ar and anomalous O are in the inputs and must not enter the mean mass. Bound: the same closed forms in two languages, rounding only; the oracle's mass conversion through N\_A u (0.99999999965e-3 kg/mol) is the same as this toolbox's, so 1e-13 relative covers every column.
\begin{itemize}[leftmargin=1.2em,itemsep=0pt,topsep=1pt]\small
\item accommodation vs the independent formula oracle (12 cases, 8 columns): measured 3.7986e-16, bound 1e-13 --- \emph{closed forms evaluated twice; rounding only}
\item SESAM sticking coefficient: closed form (26b) against quadrature of its definition (26a): measured 1.2212e-14, bound 1e-9 --- \emph{composite Simpson on 20000 panels of a Gaussian times a linear factor; the quadrature error dominates and is O(h\textasciicircum{}4), not the closed form's rounding}
\item SESAM coverage from the quadrature sticking coefficient: measured 2.3870e-15, bound 1e-9 --- \emph{the isotherm has \textbar{}dtheta/d(K P)\textbar{} \textless{}= 1, so it inherits the bound above and cannot amplify it}
\item pilinski2010 / pilinski2011 odds ratio equals the ratio of the two K: measured 2.9488e-15, bound 1e-13 --- \emph{alpha/(1-alpha) = K n\_O T, so the odds ratio is K2010/K2011 = 1.5060 independently of the state}
\end{itemize}
\item[gsiCoefficientsPublishedPlate] {\small\emph{reference \textperiodcentered\ \texttt{kernel.\allowbreak{}gsiCoefficients}}} --- FIRST ORACLE FOR THIS CHAIN THAT THIS PROJECT DID NOT WRITE. Every other one -- the ngfkernel mirror, Sentman's incident closed forms -- is this project's own reading, so a shared misreading would survive all of them. Walker, Mehta \& Koller (2014, JSR 51(5) 1522, eq. 12) print a closed-form drag coefficient for a flat plate under diffuse reflection with incomplete accommodation; tools/pyxcheck/dria\_plate\_check.py evaluates it.  AT alpha = 1 IT IS AN IDENTITY. Their sqrt(T\_kr/T\_inf)/s and this toolbox's Doornbos-form v\_re/v\_in are sqrt((2/3)(1-a) + 2 a k T\_s/(m U\textasciicircum{}2)) and sqrt((1/2)(1-a) + 2 a k T\_w/(m U\textasciicircum{}2)): the same number when a = 1, and a different treatment of the unaccommodated ...
\begin{itemize}[leftmargin=1.2em,itemsep=0pt,topsep=1pt]\small
\item plate Cp vs Walker et al. (2014) eq. 12 at full accommodation (3 speed ratios): measured 1.5966e-15, bound 1e-14 --- \emph{the two formulations are algebraically the same number at alpha = 1; rounding only}
\item \emph{(computed label)}: measured ---, bound measured 2026-09-16")
\end{itemize}
\item[dragPanelForceReference] {\small\emph{reference \textperiodcentered\ \texttt{kernel.\allowbreak{}dragPanelForce}}} --- Woeske (2019) eq. 5.20 with the projected area as reference: the constantCd plate sum reproduces it
\begin{itemize}[leftmargin=1.2em,itemsep=0pt,topsep=1pt]\small
\item constantCd plate sum vs eq. 5.20: measured 6.5562e-16, bound 1e-14 --- \emph{8-term sum, rounding only; measured 1e-16}
\item GRACE-FO Sentman Cd on the projected area (deviation from 3.5): measured 2.1407e-01, bound 0.3 --- \emph{published band 3.2-3.8 at s = 7.9; measured 3.59}
\end{itemize}
\item[radiationPanelForceReference] {\small\emph{reference \textperiodcentered\ \texttt{kernel.\allowbreak{}radiationPanelForce}}} --- sphere of tangent plates: absorbing Cr = 1, Lambert Cr = 1 + 4/9 delta, specular Cr = 1 (standard results)
\begin{itemize}[leftmargin=1.2em,itemsep=0pt,topsep=1pt]\small
\item absorbing sphere Cr - 1: measured 3.7778e-05, bound 1e-4 --- \emph{Fibonacci lattice of 4000 plates, discretisation \textasciitilde{}1/n; measured 3.8e-5}
\item Lambert sphere Cr - (1 + 4/9): measured 3.8450e-05, bound 1e-4 --- \emph{measured 3.8e-5}
\item specular sphere Cr - 1: measured 1.2336e-06, bound 1e-4 --- \emph{measured 1e-6}
\end{itemize}
\item[thermalRecoilReference] {\small\emph{reference \textperiodcentered\ \texttt{kernel.\allowbreak{}thermalRecoil}}} --- Stefan-Boltzmann with CODATA 2018: a black 1 m\textasciicircum{}2 plate at 300 K emits 459.30 W and recoils with (2/3) P / c
\begin{itemize}[leftmargin=1.2em,itemsep=0pt,topsep=1pt]\small
\item emitted flux of a black plate at 300 K vs sigma T\textasciicircum{}4: measured 0.0000e+00, bound 1e-9 W/m\textasciicircum{}2 --- \emph{CODATA 2018 sigma, exact arithmetic}
\item recoil of the +x plate vs (2/3) sigma T\textasciicircum{}4 A / c: measured 9.5163e-13, bound 1e-11 N --- \emph{published 459.3 W/m\textasciicircum{}2 to 4 digits}
\item cube at uniform temperature: net recoil: measured 0.0000e+00, bound 1e-20 N --- \emph{opposite faces cancel exactly}
\end{itemize}
\item[bodyTorquesReference] {\small\emph{reference \textperiodcentered\ \texttt{kernel.\allowbreak{}bodyTorques}}} --- Wertz (1978) eq. 17-30: tau\_x = 3 GM / r\textasciicircum{}3 (I3 - I2) sin(phi) cos(phi) for r in the y-z plane and a diagonal inertia
\begin{itemize}[leftmargin=1.2em,itemsep=0pt,topsep=1pt]\small
\item gravity-gradient torque vs Wertz eq. 17-30: measured 3.9049e-16, bound 1e-14 --- \emph{closed form, rounding only}
\item gravity-gradient torque magnitude at 490 km: measured 4.1272e-06, bound 1e-3 --- \emph{hand calculation: 3 GM/r\textasciicircum{}3 = 3.688e-6 s\textasciicircum{}-2, x 200 kg m\textasciicircum{}2 x sin(0.3) cos(0.3) = 2.0824e-4 N m}
\end{itemize}
\item[partialsPrecision] {\small\emph{precision \textperiodcentered\ \texttt{kernel.\allowbreak{}partials}}} --- The linear parameters have exact derivatives: the force itself, bitwise.
\begin{itemize}[leftmargin=1.2em,itemsep=0pt,topsep=1pt]\small
\item d(drag)/d(densityScale) minus the force: measured 0.0000e+00, bound 0 N --- \emph{the force is linear in the density: the derivative IS the force, bitwise}
\item d(radiation)/d(pressureScale) minus the force: measured 0.0000e+00, bound 0 N --- \emph{as above, bitwise}
\item d(thermal)/d(T) minus 4F/T: measured 0.0000e+00, bound 1e-30 N/K --- \emph{a T\textasciicircum{}4 law differentiated once}
\end{itemize}
\item[gsiCoefficientsPrecision] {\small\emph{precision \textperiodcentered\ \texttt{kernel.\allowbreak{}gsiCoefficients}}} --- hyperthermal limit at s = 1e4: Cp -\textgreater{} 2 g\textasciicircum{}2, Ctau -\textgreater{} 2 g l up to the re-emission sqrt(pi) sqrt(Tw/(Tinf s\textasciicircum{}2)) g and 1/(2 s\textasciicircum{}2)
\begin{itemize}[leftmargin=1.2em,itemsep=0pt,topsep=1pt]\small
\item hyperthermal Cp - 2 g\textasciicircum{}2: measured 9.7091e-05, bound bound --- \emph{re-emission term sqrt(pi) sqrt(Tw/Tinf)/s = 9.71e-5 plus 1/s\textasciicircum{}2, exact to 1e-9 after the alpha bracket rewrite (was 4e-9 off); measured 9.709e-5}
\item hyperthermal Ctau - 2 g l: measured 5.6419e-05, bound 1e-4 --- \emph{P/sqrt(pi) l = 5.6e-5 at s = 1e4; measured 5.64e-5}
\item static-gas pressure limit Cp s\textasciicircum{}2/2 - 1/2: measured 1.0081e-03, bound 2 * s0 --- \emph{first-order term in s (bulk drift), oracle 0.501 at s = 1e-3}
\item head-on s = 8 example value: measured 0.0000e+00, bound 1e-15 --- \emph{closed form in the help text}
\end{itemize}
\item[accommodationPrecision] {\small\emph{precision \textperiodcentered\ \texttt{kernel.\allowbreak{}accommodation}}} --- The published forms at the points where they are exact: K x = 1 gives 1/2, and a gas whose mean molecular mass equals the substrate's gives alpha\_s = 3/4 (mu = 1). Bounds are rounding of one product and one quotient.
\begin{itemize}[leftmargin=1.2em,itemsep=0pt,topsep=1pt]\small
\item \emph{(computed label)}: measured ---, bound 4 * eps --- \emph{x/(1 + x) at x = 1 to rounding}
\item walker2014 coverage at K\_L P\_O = 1: measured 5.5511e-17, bound 4 * eps --- \emph{Langmuir isotherm at unit argument}
\item walker2014 substrate coefficient 3 mu/(1+mu)\textasciicircum{}2: measured 0.0000e+00, bound 4 * eps --- \emph{the same arithmetic in a different association}
\item SESAM sticking coefficient far below the binding energy: measured 1.1102e-16, bound 1e-12 --- \emph{the whole Maxwellian lies below Eb; erfc of a large positive argument underflows to zero and the ratio is exactly 1}
\item SESAM sticking coefficient far above the binding energy: measured 0.0000e+00, bound 1e-12 --- \emph{no part of the distribution reaches below Eb}
\item SESAM K\_L at s\_o = 1 is K\_L,o + K\_L,f: measured 1.9285e-16, bound 1e-12 --- \emph{eq. (22) with the sticking coefficient at its upper limit}
\item SESAM K\_L at s\_o = 0 is K\_L,f: measured 0.0000e+00, bound 1e-12 --- \emph{eq. (22) with the sticking coefficient at its lower limit}
\end{itemize}
\item[dragPanelForceFlowPrecision] {\small\emph{precision \textperiodcentered\ \texttt{kernel.\allowbreak{}dragPanelForce}}} --- (1) a flow shadow of all ones is no shadow, bitwise; (2) ratios equal across the three parts and the species are the plain fraction, to rounding; (3) a two-state surface with coverage 1 is alpha = 1 bitwise and with coverage 0 is alpha = alpha\_s bitwise (1 x + 0 y = x exactly).
\begin{itemize}[leftmargin=1.2em,itemsep=0pt,topsep=1pt]\small
\item uniform flow ratios vs the plain fraction: measured 1.6297e-16, bound 1e-15 --- \emph{the fraction multiplies the coefficient sum per species instead of the dynamic pressure: one reassociation}
\end{itemize}
\item[dragPanelForcePrecision] {\small\emph{precision \textperiodcentered\ \texttt{kernel.\allowbreak{}dragPanelForce}}} --- one plate, one species: the force equals (1/2) rho U\textasciicircum{}2 A [Cp (-n) + Ctau t] with the plate's own theta
\begin{itemize}[leftmargin=1.2em,itemsep=0pt,topsep=1pt]\small
\item +x plate force vs closed form: measured 1.5898e-16, bound 1e-14 --- \emph{rounding of a 3-term product; measured 2e-16}
\item speed ratio of atomic oxygen at 1000 K, 7.6 km/s: measured 0.0000e+00, bound 1e-12 --- \emph{U / sqrt(2 k T / m\_O) = 7.4548 from CODATA constants; rounding}
\item edge-on plate force vs closed form: measured 1.5301e-16, bound 1e-14 --- \emph{gamma = 0 exactly: Cp = 1/(2 s\textasciicircum{}2) + vre P / 2, Ctau = P / sqrt(pi); rounding}
\item edge-on plate: normal over tangential force: measured 1.1102e-16, bound 1e-12 --- \emph{the two closed forms at theta = pi/2}
\end{itemize}
\item[radiationPanelForcePrecision] {\small\emph{precision \textperiodcentered\ \texttt{kernel.\allowbreak{}radiationPanelForce}}} --- single plate closed forms: absorber -P A g d, mirror -2 P A g\textasciicircum{}2 n, Lambert -P A g (d + 2/3 n)
\begin{itemize}[leftmargin=1.2em,itemsep=0pt,topsep=1pt]\small
\item absorbing plate: measured 0.0000e+00, bound 1e-21 N --- \emph{exact expression, rounding}
\item specular plate: measured 0.0000e+00, bound 1e-21 N --- \emph{exact expression, rounding}
\item Lambert plate: measured 0.0000e+00, bound 1e-21 N --- \emph{exact expression, rounding}
\item unlit -z plate: measured 0.0000e+00, bound 0 N --- \emph{gamma \textless{} 0 contributes exactly nothing}
\item absorbed flux a E gamma: measured 0.0000e+00, bound 1e-9 W/m\textasciicircum{}2 --- \emph{E = P c; rounding}
\end{itemize}
\item[thermalRecoilPrecision] {\small\emph{precision \textperiodcentered\ \texttt{kernel.\allowbreak{}thermalRecoil}}} --- steady-state thermalState + thermalRecoil: recoil = (2/3) absorbed power / c on every lit plate
\begin{itemize}[leftmargin=1.2em,itemsep=0pt,topsep=1pt]\small
\item steady state: emitted equals absorbed on lit plates: measured 4.4409e-16, bound 1e-12 --- \emph{T = (q/(eps sigma))\textasciicircum{}(1/4) then eps sigma T\textasciicircum{}4: two roundings}
\item recoil = (2/3) absorbed A / c: measured 4.4409e-16, bound 1e-12 --- \emph{unit normals (O-2 closed 2026-09-07: N renormalised); was 7e-7 with printed normals}
\end{itemize}
\item[assemblePrecision] {\small\emph{precision \textperiodcentered\ \texttt{kernel.\allowbreak{}assemble}}} --- dumbbell: equal and opposite forces on the +x and -x plates of the cube -\textgreater{} F = 0, tau = 2 L f (lever 0.5 m each)
\item[bodyTorquesPrecision] {\small\emph{precision \textperiodcentered\ \texttt{kernel.\allowbreak{}bodyTorques}}} --- magnetic torque m x B in closed form; a principal axis gives zero gravity-gradient torque
\item[partialsVectorization] {\small\emph{vectorization \textperiodcentered\ \texttt{kernel.\allowbreak{}partials}}}
\item[gsiCoefficientsVectorization] {\small\emph{vectorization \textperiodcentered\ \texttt{kernel.\allowbreak{}gsiCoefficients}}}
\item[accommodationVectorization] {\small\emph{vectorization \textperiodcentered\ \texttt{kernel.\allowbreak{}accommodation}}}
\item[dragPanelForceVectorization] {\small\emph{vectorization \textperiodcentered\ \texttt{kernel.\allowbreak{}dragPanelForce}}}
\item[radiationPanelForceVectorization] {\small\emph{vectorization \textperiodcentered\ \texttt{kernel.\allowbreak{}radiationPanelForce}}}
\item[thermalRecoilVectorization] {\small\emph{vectorization \textperiodcentered\ \texttt{kernel.\allowbreak{}thermalRecoil}}}
\item[assembleVectorization] {\small\emph{vectorization \textperiodcentered\ \texttt{kernel.\allowbreak{}assemble}}}
\item[bodyTorquesVectorization] {\small\emph{vectorization \textperiodcentered\ \texttt{kernel.\allowbreak{}bodyTorques}}}
\item[partialsRobustness] {\small\emph{robustness \textperiodcentered\ \texttt{kernel.\allowbreak{}partials}}}
\begin{itemize}[leftmargin=1.2em,itemsep=0pt,topsep=1pt]\small
\item constantCd d/dU minus 2F/U: measured 0.0000e+00, bound 1e-30 N/(m/s) --- \emph{no coefficient depends on the speed there}
\end{itemize}
\item[gsiCoefficientsRobustness] {\small\emph{robustness \textperiodcentered\ \texttt{kernel.\allowbreak{}gsiCoefficients}}} --- leeward plates: exponentially small, finite, non-negative; theta = pi exactly; extreme s finite
\begin{itemize}[leftmargin=1.2em,itemsep=0pt,topsep=1pt]\small
\item theta = pi: Ctau at the sin(pi) rounding level: measured 6.7876e-15, bound 1e-14 --- \emph{sin(pi) = 1.2e-16 times P/sqrt(pi) = 56 at s = 0.01}
\end{itemize}
\item[accommodationRobustness] {\small\emph{robustness \textperiodcentered\ \texttt{kernel.\allowbreak{}accommodation}}}
\item[dragPanelForceRobustness] {\small\emph{robustness \textperiodcentered\ \texttt{kernel.\allowbreak{}dragPanelForce}}} --- a species column of zeros is skipped (no NaN from s = U/0), an exactly edge-on and an exactly head-on plate give finite forces
\begin{itemize}[leftmargin=1.2em,itemsep=0pt,topsep=1pt]\small
\item leeward -x plate force at s = 7.45: measured 9.9742e-33, bound 1e-28 N --- \emph{exp(-s\textasciicircum{}2) = 1e-24 times (1/2) rho U\textasciicircum{}2 A = 8e-5; measured 3.2e-30}
\end{itemize}
\item[radiationPanelForceRobustness] {\small\emph{robustness \textperiodcentered\ \texttt{kernel.\allowbreak{}radiationPanelForce}}} --- grazing incidence (gamma = 0) and zero pressure give exact zeros; a source exactly behind contributes nothing
\item[thermalRecoilRobustness] {\small\emph{robustness \textperiodcentered\ \texttt{kernel.\allowbreak{}thermalRecoil}}}
\item[assembleRobustness] {\small\emph{robustness \textperiodcentered\ \texttt{kernel.\allowbreak{}assemble}}}
\item[bodyTorquesRobustness] {\small\emph{robustness \textperiodcentered\ \texttt{kernel.\allowbreak{}bodyTorques}}}
\item[partialsMetamorphic] {\small\emph{metamorphic \textperiodcentered\ \texttt{kernel.\allowbreak{}partials}}}
\begin{itemize}[leftmargin=1.2em,itemsep=0pt,topsep=1pt]\small
\item \emph{(computed label)}: measured ---, bound 1e-15 --- \emph{an exact power of two through a linear factor}
\item half the exposure halves d/d(alpha): measured 0.0000e+00, bound 1e-15 --- \emph{the shadow multiplies the same scale}
\end{itemize}
\item[gsiCoefficientsMetamorphic] {\small\emph{metamorphic \textperiodcentered\ \texttt{kernel.\allowbreak{}gsiCoefficients}}} --- maxwell with sigma = 0 is sentman bitwise; sigma = 1 has no re-emission (alpha has no effect); the drag/lift projections invert
\begin{itemize}[leftmargin=1.2em,itemsep=0pt,topsep=1pt]\small
\item normal/tangential from drag/lift round trip: measured 1.1102e-16, bound 1e-15 --- \emph{rotation of a 2-vector, rounding}
\end{itemize}
\item[accommodationMetamorphic] {\small\emph{metamorphic \textperiodcentered\ \texttt{kernel.\allowbreak{}accommodation}}}
\begin{itemize}[leftmargin=1.2em,itemsep=0pt,topsep=1pt]\small
\item pilinski2010 depends on n\_O T only: measured 0.0000e+00, bound 4 * eps --- \emph{a power of two moved between two factors of one product}
\item SESAM ram pressure is exactly linear in n\_O: measured 0.0000e+00, bound 4 * eps --- \emph{n\_O enters eq. (15) only through rho\_O; a power of two through one factor}
\end{itemize}
\item[dragPanelForceMetamorphic] {\small\emph{metamorphic \textperiodcentered\ \texttt{kernel.\allowbreak{}dragPanelForce}}} --- density scaling by 2 doubles the force exactly; reversing v on the symmetric cube reverses the force; shadow 0 kills a plate
\begin{itemize}[leftmargin=1.2em,itemsep=0pt,topsep=1pt]\small
\item cube: F(-v) = -F(v): measured 2.4439e-18, bound 1e-14 --- \emph{opposite plates swap roles; rounding}
\end{itemize}
\item[radiationPanelForceMetamorphic] {\small\emph{metamorphic \textperiodcentered\ \texttt{kernel.\allowbreak{}radiationPanelForce}}} --- superposition: two sources together equal the sum of each alone; a duplicated source doubles the force exactly
\begin{itemize}[leftmargin=1.2em,itemsep=0pt,topsep=1pt]\small
\item superposition of two sources: measured 0.0000e+00, bound 1e-15 --- \emph{one addition per source, rounding}
\end{itemize}
\item[thermalRecoilMetamorphic] {\small\emph{metamorphic \textperiodcentered\ \texttt{kernel.\allowbreak{}thermalRecoil}}}
\item[assembleMetamorphic] {\small\emph{metamorphic \textperiodcentered\ \texttt{kernel.\allowbreak{}assemble}}} --- concatenating panel arrays equals adding them; moving the reference point changes the torque by -delta x F
\begin{itemize}[leftmargin=1.2em,itemsep=0pt,topsep=1pt]\small
\item assemble is additive: measured 3.5071e-16, bound 1e-14 --- \emph{8-term sums in two orders, rounding}
\item torque transfer tau' = tau - delta x F: measured 2.7061e-16, bound 1e-14 --- \emph{rounding}
\end{itemize}
\item[bodyTorquesMetamorphic] {\small\emph{metamorphic \textperiodcentered\ \texttt{kernel.\allowbreak{}bodyTorques}}}
\begin{itemize}[leftmargin=1.2em,itemsep=0pt,topsep=1pt]\small
\item gravity gradient scales as r\textasciicircum{}-3: measured 0.0000e+00, bound 1e-14 --- \emph{rounding}
\end{itemize}
\item[partialsSpeed] {\small\emph{speed \textperiodcentered\ \texttt{kernel.\allowbreak{}partials}}} --- One derivative costs about one force evaluation: the coefficients are computed once more.
\begin{itemize}[leftmargin=1.2em,itemsep=0pt,topsep=1pt]\small
\item d(drag)/d(alpha) vs one force evaluation: measured 3.4209e+00, bound 9 --- \emph{budget 3x the expected 3: the derivative evaluates the force and the coefficients again; 5 repeats, median}
\end{itemize}
\item[gsiCoefficientsSpeed] {\small\emph{speed \textperiodcentered\ \texttt{kernel.\allowbreak{}gsiCoefficients}}} --- 1e5 elements at once vs 100 scalar calls: the array form must be at least 20x cheaper per element (measured \textasciitilde{}200x)
\begin{itemize}[leftmargin=1.2em,itemsep=0pt,topsep=1pt]\small
\item gsiCoefficients: 1e5-element array vs 100 scalar calls, time ratio: measured 8.9834e-01, bound 50 --- \emph{budget 3x the expected 1000/60; N = 1e5, 5 repeats, median}
\end{itemize}
\item[dragPanelForceSpeed] {\small\emph{speed \textperiodcentered\ \texttt{kernel.\allowbreak{}dragPanelForce}}}
\begin{itemize}[leftmargin=1.2em,itemsep=0pt,topsep=1pt]\small
\item dragPanelForce: K = 1000 batched vs one epoch, time ratio: measured 7.7380e+00, bound 100 --- \emph{budget 3x the expected \textasciitilde{}30; 8 species x 8 plates; 5 repeats, median}
\end{itemize}
\item[radiationPanelForceSpeed] {\small\emph{speed \textperiodcentered\ \texttt{kernel.\allowbreak{}radiationPanelForce}}}
\begin{itemize}[leftmargin=1.2em,itemsep=0pt,topsep=1pt]\small
\item radiationPanelForce: K = 1000 x 51 sources vs one epoch, time ratio: measured 2.7115e+01, bound 300 --- \emph{budget 3x the expected \textasciitilde{}100 (the source loop is not batched over K); 5 repeats, median}
\end{itemize}
\item[thermalRecoilSpeed] {\small\emph{speed \textperiodcentered\ \texttt{kernel.\allowbreak{}thermalRecoil}}}
\begin{itemize}[leftmargin=1.2em,itemsep=0pt,topsep=1pt]\small
\item thermalRecoil: K = 1e4 vs one epoch, time ratio: measured 5.6605e+00, bound 60 --- \emph{budget 3x the expected \textasciitilde{}20; 5 repeats, median}
\end{itemize}
\item[assembleSpeed] {\small\emph{speed \textperiodcentered\ \texttt{kernel.\allowbreak{}assemble}}}
\begin{itemize}[leftmargin=1.2em,itemsep=0pt,topsep=1pt]\small
\item assemble: K = 1e4 vs one epoch, time ratio: measured 4.5258e+00, bound 60 --- \emph{budget 3x the expected \textasciitilde{}20; 5 repeats, median}
\end{itemize}
\item[bodyTorquesSpeed] {\small\emph{speed \textperiodcentered\ \texttt{kernel.\allowbreak{}bodyTorques}}}
\begin{itemize}[leftmargin=1.2em,itemsep=0pt,topsep=1pt]\small
\item bodyTorques: K = 1e4 vs one epoch, time ratio: measured 1.8522e+00, bound 60 --- \emph{budget 3x the expected \textasciitilde{}20; 5 repeats, median}
\end{itemize}
\end{description}

\subsection{\texttt{tNgfMex}}
\noindent Parity of the compiled single-epoch kernel against the MATLAB units it ports (LP-6b, WP-17b).

\begin{description}[leftmargin=1em,style=nextline,itemsep=2pt,font=\normalfont\ttfamily\small]
\item[singleEpochContract] {\small\emph{contract \textperiodcentered\ \texttt{kernel.\allowbreak{}singleEpoch}}}
\item[singleEpochGatewayContract] {\small\emph{contract \textperiodcentered\ \texttt{kernel.\allowbreak{}singleEpoch}}} --- The gas-surface model reaches the C as a code, and only three codes exist. The gateway must be unreachable except through the unit above: the private folder is what enforces that, and this asserts the folder is where the unit looks for it.
\begin{itemize}[leftmargin=1.2em,itemsep=0pt,topsep=1pt]\small
\item constantCd ignores the gas temperature: measured 0.0000e+00, bound 0 N --- \emph{the model has no speed ratio, so T cannot enter}
\item constantCd accepts the NaN temperature the exponential backend returns: measured 0.0000e+00, bound 0 N --- \emph{an input the model does not read is not validated}
\end{itemize}
\item[singleEpochMirrorReference] {\small\emph{reference \textperiodcentered\ \texttt{kernel.\allowbreak{}singleEpoch}}} --- The frozen Python mirror (WP-33) is the independent authority. Each unit's cases are replayed through the composite by silencing the other two effects: a zero source list and no panel temperature leave only the gas force, and a zero species vector leaves only the radiation and the recoil.
\begin{itemize}[leftmargin=1.2em,itemsep=0pt,topsep=1pt]\small
\item \emph{(computed label)}: measured ---, bound 1e-12 --- \emph{the C, the MATLAB and the mirror are three derivations of the same closed form in double precision}
\item \emph{(computed label)}: measured ---, bound 1e-12 --- \emph{independently derived optical model}
\item singleEpoch absorbed flux vs the mirror: measured 4.7559e-16, bound 1e-12 --- \emph{a E gamma summed over sources}
\item \emph{(computed label)}: measured ---, bound 1e-13 --- \emph{Stefan-Boltzmann and one hemispherical factor: only rounding may differ}
\end{itemize}
\item[singleEpochMatlabParityPrecision] {\small\emph{precision \textperiodcentered\ \texttt{kernel.\allowbreak{}singleEpoch}}} --- CLAIM OF THE LEVEL PROMPT: the compiled path reproduces the MATLAB units to 1e-14. Relative to the largest component of each column, because the three effects differ by two orders of magnitude and an absolute bound would only measure the largest.
\begin{itemize}[leftmargin=1.2em,itemsep=0pt,topsep=1pt]\small
\item singleEpoch force vs the MATLAB kernel chain (5 epochs): measured 4.4897e-16, bound 1e-14 --- \emph{the same closed forms in the same association order; only rounding may differ}
\item singleEpoch torque vs the MATLAB kernel chain: measured 1.9492e-15, bound 1e-14 --- \emph{cross products and sums in the same plate order}
\item singleEpoch per-panel gas force vs the MATLAB kernel: measured 5.9779e-16, bound 1e-14 --- \emph{the reduction cannot hide a per-panel error if the panels are compared too}
\item singleEpoch gamma vs the MATLAB kernel: measured 0.0000e+00, bound 1e-15 --- \emph{one dot product per plate}
\item singleEpoch speed ratio vs the MATLAB kernel: measured 0.0000e+00, bound 1e-15 --- \emph{U / sqrt(2 k T / m) per species}
\item singleEpoch mixture Cp vs the MATLAB kernel: measured 3.1496e-15, bound 1e-13 --- \emph{mass-weighted sum over eight species in the same order}
\end{itemize}
\item[singleEpochEndToEndPrecision] {\small\emph{precision \textperiodcentered\ \texttt{kernel.\allowbreak{}singleEpoch}}} --- TWO PATHS MUST AGREE, and this is where that is checked at the level a user sees. Owner ruling of 2026-09-08: the vectorised MATLAB path and the pointwise C path are BOTH first class; neither replaces the other. Two paths that are both kept are two paths that must be kept in step, so this test composes the environment exactly as Model.evaluate composes it, hands one epoch to the compiled kernel, and requires the acceleration back to be the one Model.accel returns for that epoch.  The fast preset is used deliberately: it sets geom.selfShadow "none" and thermal.model "none", so the shadow arrays are all ones and the recoil is zero, and what is left under test is the composition itself ...
\begin{itemize}[leftmargin=1.2em,itemsep=0pt,topsep=1pt]\small
\item compiled single epoch vs the vectorised path, per effect: measured 1.0294e-15, bound 1e-13 --- \emph{the same environment, the same panel loop, the same reduction: only the association of the mass division and the C rounding may differ}
\item compiled single epoch vs the vectorised path, total acceleration: measured 1.2054e-15, bound 1e-13 --- \emph{Model.accel adds three quotients, this adds three forces and divides once}
\item the relative velocity handed to the C equals the Model's: measured 1.1979e-16, bound 1e-15 --- \emph{same env.wind, same rotation, different product routine: pagemtimes over pages vs one matrix-vector product}
\end{itemize}
\item[singleEpochVectorization] {\small\emph{vectorization \textperiodcentered\ \texttt{kernel.\allowbreak{}singleEpoch}}} --- The compiled path is single-epoch by construction, so the claim is the other way round: K calls to it must reproduce the K-vectorised MATLAB kernel column by column. That is the property an integrator relies on when it swaps one for the other.
\begin{itemize}[leftmargin=1.2em,itemsep=0pt,topsep=1pt]\small
\item K single-epoch calls vs the K-vectorised MATLAB kernel: measured 5.8640e-16, bound 1e-14 --- \emph{epoch k of a batch and the same epoch alone are the same physics}
\end{itemize}
\item[singleEpochRobustness] {\small\emph{robustness \textperiodcentered\ \texttt{kernel.\allowbreak{}singleEpoch}}}
\begin{itemize}[leftmargin=1.2em,itemsep=0pt,topsep=1pt]\small
\item singleEpoch under total shadow: measured 0.0000e+00, bound 0 N --- \emph{a zero fraction multiplies every panel force, both effects}
\item reversed stream: cosine between the gas force and the velocity: measured -9.9997e-01, bound 0 --- \emph{a dissipative force has a non-positive projection on the velocity, whichever plates face the flow}
\end{itemize}
\item[singleEpochMetamorphic] {\small\emph{metamorphic \textperiodcentered\ \texttt{kernel.\allowbreak{}singleEpoch}}}
\begin{itemize}[leftmargin=1.2em,itemsep=0pt,topsep=1pt]\small
\item gas force doubles with the density: measured 0.0000e+00, bound 1e-15 --- \emph{rho enters every species term linearly; the coefficients do not depend on it}
\item radiation force triples with the pressure: measured 1.5678e-16, bound 1e-15 --- \emph{P multiplies the whole optical expression}
\item thermal recoil scales as T\textasciicircum{}4: measured 0.0000e+00, bound 1e-15 --- \emph{Stefan-Boltzmann with an exact power of two}
\item CdScale multiplies the gas column alone: measured 3.8154e-18, bound 1e-15 --- \emph{the scales are applied per effect, where Model.evaluate applies them}
\item CrScale multiplies the radiation column alone: measured 0.0000e+00, bound 1e-15 --- \emph{same, for the radiation column}
\item torque shifts by -d x F when the reference point moves: measured 3.9656e-16, bound 1e-13 --- \emph{the lever arms all change by the same vector, so the sum changes by that vector crossed into the total force}
\end{itemize}
\item[singleEpochSpeed] {\small\emph{speed \textperiodcentered\ \texttt{kernel.\allowbreak{}singleEpoch}}} --- THIS MEASURES AND WHY. HANDOVER D.6 justifies the compiled path with a MATLAB call overhead of 5-15 microseconds, and LP-6's own measurement contradicted it. The figures that comment quoted (88 ms standard, 2 ms fast) are superseded -- WP-C3 took the standard preset to 0.798 ms per epoch at K=100, and most of what the 88 ms recorded turned out to BE per-call overhead rather than model work (O-29, closed 2026-09-18: one Model.accel call costs \textasciitilde{}15 ms of setup in the standard preset before it does any per-epoch work, so D.6's 15 us was wrong by three orders of magnitude in the same direction its author feared). The premise needs a number either way, and this is it: the ratio of the compiled ...
\begin{itemize}[leftmargin=1.2em,itemsep=0pt,topsep=1pt]\small
\item single-epoch kernel: compiled / MATLAB chain (median of 5 x 200): measured 4.5269e-01, bound 1.0
\item single-epoch compiled kernel, absolute cost per epoch: measured 1.1168e+03, bound 10000 us --- \emph{smoke ceiling, not a claim; 901.8 us measured on the CI runner 2026-09-08, MATLAB chain 4149.6 us}
\end{itemize}
\end{description}

\subsection{\texttt{tNgfMirror}}
\noindent Parity of the MATLAB kernel against the independent Python mirror (WP-33).

\begin{description}[leftmargin=1em,style=nextline,itemsep=2pt,font=\normalfont\ttfamily\small]
\item[mirrorFixtureFidelity] {\small\emph{contract \textperiodcentered\ no unit}} --- The exported plate table must still be what geom produces, or every case below is stale.
\begin{itemize}[leftmargin=1.2em,itemsep=0pt,topsep=1pt]\small
\item fixture areas vs live geom.graceFO: measured 0.0000e+00, bound 0 m\textasciicircum{}2 --- \emph{exported from the same object; must be bitwise}
\item fixture normals vs live geom.graceFO: measured 3.3307e-16, bound 1e-15 --- \emph{one CSV round trip through 17 significant digits}
\item fixture centroids vs live geom.graceFO: measured 4.4409e-16, bound 1e-15 m --- \emph{one CSV round trip}
\item fixture optics vs live geom.graceFO: measured 0.0000e+00, bound 1e-15 --- \emph{one CSV round trip}
\item fixture inertia vs live geom.graceFO: measured 0.0000e+00, bound 1e-9 kg m\textasciicircum{}2 --- \emph{printed to 17 digits from Table 4}
\end{itemize}
\item[gsiCoefficientsMirror] {\small\emph{reference \textperiodcentered\ \texttt{kernel.\allowbreak{}gsiCoefficients}}}
\begin{itemize}[leftmargin=1.2em,itemsep=0pt,topsep=1pt]\small
\item \emph{(computed label)}: measured ---, bound 1e-12 --- \emph{two independently derived closed forms in double precision}
\item \emph{(computed label)}: measured ---, bound 1e-11 --- \emph{leeward plates: both sides must be negligible, relative agreement is meaningless there}
\end{itemize}
\item[dragPanelForceMirror] {\small\emph{reference \textperiodcentered\ \texttt{kernel.\allowbreak{}dragPanelForce}}}
\begin{itemize}[leftmargin=1.2em,itemsep=0pt,topsep=1pt]\small
\item \emph{(computed label)}: measured ---, bound 1e-12 --- \emph{independent mixture rule and plate geometry, double precision}
\item dragPanelForce, plates negligible in their case: measured 3.8920e-16, bound 1e-11 --- \emph{leeward plates: negligible on both sides}
\end{itemize}
\item[radiationPanelForceMirror] {\small\emph{reference \textperiodcentered\ \texttt{kernel.\allowbreak{}radiationPanelForce}}}
\begin{itemize}[leftmargin=1.2em,itemsep=0pt,topsep=1pt]\small
\item \emph{(computed label)}: measured ---, bound 1e-12 --- \emph{independently derived optical model, double precision}
\item radiationPanelForce absorbed flux vs the mirror: measured 4.7559e-16, bound 1e-12 --- \emph{a E gamma summed over sources}
\end{itemize}
\item[thermalRecoilMirror] {\small\emph{reference \textperiodcentered\ \texttt{kernel.\allowbreak{}thermalRecoil}}}
\begin{itemize}[leftmargin=1.2em,itemsep=0pt,topsep=1pt]\small
\item \emph{(computed label)}: measured ---, bound 1e-13 --- \emph{Stefan-Boltzmann and one hemispherical factor: only rounding may differ}
\end{itemize}
\item[assembleMirror] {\small\emph{reference \textperiodcentered\ \texttt{kernel.\allowbreak{}assemble}}}
\begin{itemize}[leftmargin=1.2em,itemsep=0pt,topsep=1pt]\small
\item \emph{(computed label)}: measured ---, bound 1e-13 --- \emph{sums of 6 to 8 terms in two orders}
\item assemble torque vs the mirror: measured 0.0000e+00, bound 1e-12 --- \emph{cross products then sums, two orders}
\end{itemize}
\item[bodyTorquesMirror] {\small\emph{reference \textperiodcentered\ \texttt{kernel.\allowbreak{}bodyTorques}}}
\begin{itemize}[leftmargin=1.2em,itemsep=0pt,topsep=1pt]\small
\item \emph{(computed label)}: measured ---, bound 1e-12 --- \emph{independently derived gravity-gradient and dipole forms}
\end{itemize}
\end{description}

\subsection{\texttt{tNgfModel}}
\noindent Tests of NonGravForces.Model: the single point of access (LP-6, WP-14 to WP-17).

\begin{description}[leftmargin=1em,style=nextline,itemsep=2pt,font=\normalfont\ttfamily\small]
\item[modelContract] {\small\emph{contract \textperiodcentered\ \texttt{Model}}}
\item[modelReference] {\small\emph{reference \textperiodcentered\ \texttt{Model}}} --- The measured ACT1B accelerations of the same real day: shape, not level.
\begin{itemize}[leftmargin=1.2em,itemsep=0pt,topsep=1pt]\small
\item ACT1B cross-track correlation shortfall: measured 2.7768e-02, bound 0.10 --- \emph{measured 0.980; the cross-track is radiation-dominated and the model reproduces its shape}
\item ACT1B radial correlation shortfall: measured 4.6458e-03, bound 0.10 --- \emph{measured 0.994}
\item ACT1B along-track correlation (recorded, not a claim): measured 3.1854e-01, bound 1.0 --- \emph{measured 0.190; the excerpt is too short for the drag signal to rise above the instrument bias drift}
\item Sentman lift: drag angle off the stream: measured 8.1709e-01, bound 2.0 deg --- \emph{the flank shear of an 8-plate GRACE-FO model; measured 0.34 deg}
\item constantCd drag direction vs -v\_rel: measured 4.4409e-16, bound 1e-14 --- \emph{constantCd resolves every windward plate along the stream by construction: no lift; measured 1e-16}
\item thermal recoil magnitude vs Siemes 2023 (6.4 nm/s\textasciicircum{}2): measured 9.2528e-10, bound 4e-9 m/s\textasciicircum{}2 --- \emph{published cross-track RMS for GRACE-A 2008; measured median 5.475e-9 over this arc}
\item thermal recoil cross-track RMS vs Hladczuk 2024 case (a) (4.22 nm/s\textasciicircum{}2): measured 1.0433e-09, bound 2.5e-9 m/s\textasciicircum{}2 --- \emph{the residual left by omitting thermal emission bounds its magnitude; measured 5.263e-9, between the GRACE-FO and GRACE-A published figures}
\end{itemize}
\item[modelGroopsReference] {\small\emph{reference \textperiodcentered\ \texttt{Model}}} --- O-19: GROOPS against this toolbox on the real day. TU Graz's operational nonConservativeForces for 2020-01-01 GRACE-C (GRACEFO-1) are MODELLED by GROOPS (owner, 2026-10-26), in the satellite reference frame at 10 s, and published with the attitude (SRF -\textgreater{} CRF, scalar first) and the reduced-dynamic orbit (CRF) they were computed on. This test feeds THEIR state into Model.accel, so the two differ only by force model: density, gas-surface interaction, radiation and thermal. All three files are read with gracefoReading.readGroopsInstrument (\textgreater{}= 1.1, which reads real GROOPS output), the first use of that reader on TU Graz's bytes in this repository (CLAUDE.md section 7).  Every bound is ...
\begin{itemize}[leftmargin=1.2em,itemsep=0pt,topsep=1pt]\small
\item \emph{(computed label)}: measured ---, bound actBound(i) --- \emph{the oracle against the measurement; independent value 2026-10-26: shortfalls 0.010, 0.018, 0.090}
\item GROOPS (TU Graz) vs Model: \textbar{}log\textbar{} of the median magnitude ratio: measured 1.1232e-01, bound log(1.5) --- \emph{measured ratio 1.119 (CI 40411254932); a frame or unit error is a factor of 1e3 or a sign}
\item \emph{(computed label)}: measured ---, bound rhoBound(i) --- \emph{measured 0.017 / 0.0004 / 0.032 (CI 40406425492)}
\item \emph{(computed label)}: measured ---, bound ratioBound(i) --- \emph{measured \textbar{}log\textbar{} 0.404 / 0.166 / 0.180 (CI 40406425492)}
\item GROOPS vs Model RMS of the difference / RMS of GROOPS: measured 3.9100e-01, bound 0.6 --- \emph{measured 0.391 (CI 40411254932); drag dominates the along-track, radiation the cross-track}
\item GROOPS comparison: modelled drag share of \textbar{}a\textbar{} (recorded): measured 4.9719e-01, bound 0.75 --- \emph{measured 0.497 (CI 40411254932): context for the along-track numbers above}
\end{itemize}
\item[modelAct1bReference] {\small\emph{reference \textperiodcentered\ \texttt{Model}}} --- MODEL AGAINST THE MEASUREMENT, the way Huckfeldt et al. (2024, ASR 73, Sec. 3.5, eqs. 7-10) validate theirs: ACT1B with the thruster firings cut, one bias per day and per axis, scale 1, then the daily RMS of model minus calibrated ACT1B. A bias-only fit keeps model errors out of the calibration parameters; a bias-plus-scale fit is recorded beside it, because GRACE-C's ACT1B scale is not 1 (tools/baseline\_reference.json: S\_xx 0.762).  Data: the committed 30 min of each of 2020-01-01..03 (GRACE-C, data/samples/multiday), this toolbox's own state -- GNV1B, SCA1B conjugated once by io.readSCA1B -- and the standard preset. THR1B is committed for 2020-01-01 only; that day's three firings in ...
\begin{itemize}[leftmargin=1.2em,itemsep=0pt,topsep=1pt]\small
\item \emph{(computed label)}: measured ---, bound rbBound(i) m/s\textasciicircum{}2 --- \emph{measured 5.55 / 3.37 / 2.10 nm/s\textasciicircum{}2 (CI 40411254932); Huckfeldt 2024: 3-5 x, \textless{}= 2 y and z; GROOPS 2.7 / 3.1 / 1.8}
\item \emph{(computed label)}: measured ---, bound 1 + 1e-12 --- \emph{least squares with one more column cannot do worse: an identity}
\item \emph{(computed label)}: measured ---, bound log(1.25) --- \emph{measured 1.034 / 1.026 / 0.910 (CI 40411254932); GROOPS against the same ACT1B 0.92 / 0.85 / 0.87}
\item \emph{(computed label)}: measured ---, bound rhoBound(i) --- \emph{measured 0.187 / 0.060 / 0.045 (CI 40411254932); GROOPS against the same ACT1B 0.07 / 0.06 / 0.03}
\item \emph{(computed label)}: measured ---, bound 0.5 --- \emph{what the THR1B mask is worth on the one day it is committed for; measured 0.019 / 0.314 / 0.067 (CI 40411254932)}
\end{itemize}
\item[modelDragDiagnosis] {\small\emph{reference \textperiodcentered\ \texttt{Model}}} --- THE ALONG-TRACK EXCESS COMES FROM. modelGroopsReference and modelAct1bReference found the model's along-track 1.5x GROOPS's RMS (drag mean 1.67e-8 against a GROOPS total of 0.71e-8) and 5.5 nm/s\textasciicircum{}2 from ACT1B against GROOPS's 2.7. Two diagnoses on the same committed data, 10 s grid:  (a) EFFECT SCALES: each reference is written as k\_d drag + k\_r radiation + k\_t thermal of the standard preset, per axis (against ACT1B with one bias per day besides). A k\_d that is constant over the orbit says SCALE (density or Cd); one that varies with the orbit says geometry, wind or the gas-surface model. k = (1, 1, 1) is a feasible point, so the fitted residual can only be smaller: that identity is the ...
\begin{itemize}[leftmargin=1.2em,itemsep=0pt,topsep=1pt]\small
\item \emph{(computed label)}: measured ---, bound 1 + 1e-12 --- \emph{k = (1, 1, 1) is feasible, so the least-squares fit cannot do worse: an identity}
\item drag diagnosis: ACT1B drag-scale fit RMS over bias-only RMS, along-track: measured 2.9517e-01, bound 1 + 1e-12 --- \emph{the bias-only fit is the k\_drag = 1 point of the same design: an identity}
\item drag diagnosis: GROOPS along-track, effect-scale fit RMS over RMS at k = 1: measured 5.0211e-02, bound 0.1 --- \emph{measured 0.050 (CI 40447628923): one drag scale explains 95 \% of the along-track difference}
\item drag diagnosis: \textbar{}log k\_drag\textbar{} along-track vs GROOPS: measured 7.2996e-01, bound log(2.5) --- \emph{measured k 0.482 (CI 40447628923): GROOPS's drag is half of this toolbox's}
\item drag diagnosis: \textbar{}log k\_drag\textbar{} along-track vs ACT1B: measured 1.5657e+00, bound log(10) --- \emph{measured 0.209 (CI 40450067703; 3 x 30 min, ACT1B's own scale and per-day biases in it): the measurement wants less drag too}
\item drag diagnosis: k\_drag window spread (max/min) along-track vs GROOPS: measured 1.1020e+00, bound 1.3 --- \emph{measured 1.10 (0.441-0.486, CI 40447628923): a SCALE, not an orbit-dependent shape}
\item \emph{(computed label)}: measured ---, bound log(2) --- \emph{ratios 0.915 msis21, 0.751 constantCd 2.4, 0.905 alpha 1 (CI 40447628923); the combination recorded}
\item \emph{(computed label)}: measured ---, bound 8e-9 m/s\textasciicircum{}2 --- \emph{measured 5.57 / 4.82 / 4.56 / 4.89 nm/s\textasciicircum{}2 (standard, msis21, constantCd 2.4, alpha 1; CI 40447628923)}
\end{itemize}
\item[modelPrecision] {\small\emph{precision \textperiodcentered\ \texttt{Model}}}
\begin{itemize}[leftmargin=1.2em,itemsep=0pt,topsep=1pt]\small
\item accel minus its three effects: measured 0.0000e+00, bound 0 m/s\textasciicircum{}2 --- \emph{the same three arrays summed in the same order: bitwise}
\item budget total RMS vs accel: measured 0.0000e+00, bound 0 m/s\textasciicircum{}2 --- \emph{bitwise: budget reduces the same array}
\item gcrs vs body magnitude: measured 9.9262e-24, bound 1e-23 m/s\textasciicircum{}2 --- \emph{an orthogonal rotation of a 3e-8 vector; measured 7e-24}
\end{itemize}
\item[modelVectorization] {\small\emph{vectorization \textperiodcentered\ \texttt{Model}}} --- A batch of epochs equals the same epochs one at a time, bitwise, and chunking the series changes nothing: that is what makes a chunked or parallel run trustworthy.
\item[modelRobustness] {\small\emph{robustness \textperiodcentered\ \texttt{Model}}}
\item[modelMetamorphic] {\small\emph{metamorphic \textperiodcentered\ \texttt{Model}}}
\begin{itemize}[leftmargin=1.2em,itemsep=0pt,topsep=1pt]\small
\item thermal inertia vs instantaneous: the step it is worth: measured 9.4977e-10, bound 4e-9 m/s\textasciicircum{}2 --- \emph{[H24] Table 3 (b)-\textgreater{}(c) buys 1.34 nm/s\textasciicircum{}2 of residual on GRACE-FO; measured 0.950e-9 as a model difference, the same order}
\item drag doubles with cdScale: measured 0.0000e+00, bound 1e-15 --- \emph{an exact power of two through a linear factor}
\item radiation doubles with crScale: measured 0.0000e+00, bound 1e-15 --- \emph{as above}
\item drag msis21 / msis00 through the Model, largest departure from 1: measured 9.6465e-02, bound log(2) --- \emph{two climatologies of one thermosphere (the density ratio measured 0.88-0.99 at 490 km that day): a plausibility bound on the wiring}
\item \emph{(computed label)}: measured ---, bound 1e-6 --- \emph{central difference in the option, step 1e-6 relative; its own truncation dominates}
\item accel predicted by the r-partials over a 37 m step: measured 4.5121e-08, bound 5e-3 --- \emph{second-order truncation of a linear extrapolation over 37 m, against a 500 km scale height}
\end{itemize}
\item[modelSpeed] {\small\emph{speed \textperiodcentered\ \texttt{Model}}} --- The fast preset exists to be fast: no self-shadowing, no Earth radiation, no thermal model. Until LP-C it was 45x the cheaper (2 against 89 ms per epoch, both superseded); with the O(1) cube map, flow shadowing and provenance on demand the standard preset cost 1.2 ms per epoch on the host (2026-09-15, 100 epochs; superseded by tools/benchmark\_reference.csv, 0.798 ms on 2026-09-18) and the fast one 0.3 ms, a ratio of 0.25. The fast preset must still be the cheaper one by a clear margin: budget 3x the measured ratio.
\begin{itemize}[leftmargin=1.2em,itemsep=0pt,topsep=1pt]\small
\item fast preset vs standard preset, same arc (40 epochs): measured 5.5297e-01, bound 0.75 --- \emph{measured 0.25 on 2026-09-15 (it was 0.02 before WP-C3; that 88 ms of HANDOVER O-29 was per-call overhead, not physics)}
\end{itemize}
\item[modelDensityScaleContract] {\small\emph{contract \textperiodcentered\ \texttt{Model}}}
\item[modelDensityScaleMetamorphic] {\small\emph{metamorphic \textperiodcentered\ \texttt{Model}}} --- The gas force is linear in the species densities, so the scaled Model's drag is the unscaled Model's drag times the factor at each epoch, and nothing else moves (radiation, thermal, the factor itself is S.at at the Model's own geodetic height). Bound from the rounding of one product per term.
\begin{itemize}[leftmargin=1.2em,itemsep=0pt,topsep=1pt]\small
\item scaled drag / (factor x unscaled drag) - 1: measured 2.2078e-16, bound 1e-12 --- \emph{linear in the densities; a few products' rounding}
\end{itemize}
\end{description}

\subsection{\texttt{tNgfRevolutions}}
\noindent Tests of the per-revolution fit and its attribution (N-03c): app.perRevolution, app.attributeRevolutions.

\begin{description}[leftmargin=1em,style=nextline,itemsep=2pt,font=\normalfont\ttfamily\small]
\item[perRevolutionContract] {\small\emph{contract \textperiodcentered\ \texttt{app.\allowbreak{}perRevolution}}}
\item[attributeRevolutionsContract] {\small\emph{contract \textperiodcentered\ \texttt{app.\allowbreak{}attributeRevolutions}}}
\item[perRevolutionReference] {\small\emph{reference \textperiodcentered\ \texttt{app.\allowbreak{}perRevolution}}} --- The composition: each window is app.fitCalibration on that window's rows, bitwise.
\item[attributeRevolutionsReference] {\small\emph{reference \textperiodcentered\ \texttt{app.\allowbreak{}attributeRevolutions}}} --- The composition: the coefficients are regressDrivers' on the same drivers, bitwise ...
\begin{itemize}[leftmargin=1.2em,itemsep=0pt,topsep=1pt]\small
\item attributeRevolutions beta' of a Sun-plane orbit: measured 2.1416e-04, bound 5e-3 rad --- \emph{the Sun moves \textasciitilde{}0.04 deg/h out of the plane fixed at the first epoch; 3 revolutions \textasciitilde{} 4.7 h}
\item attributeRevolutions eclipse fraction vs asin(Re/R)/pi: measured 5.1845e-04, bound 0.01 --- \emph{conical vs cylindrical shadow and a 10 s sampling: the penumbra is \textasciitilde{}10 s of a 5650 s revolution}
\end{itemize}
\item[perRevolutionPrecision] {\small\emph{precision \textperiodcentered\ \texttt{app.\allowbreak{}perRevolution}}} --- Noise-free with a constant planted scale and bias: every passing window returns the truth; the noise ASD is 1e-30, so what remains is the rounding of the filter and the solve.
\begin{itemize}[leftmargin=1.2em,itemsep=0pt,topsep=1pt]\small
\item perRevolution noise-free scale error: measured 1.1102e-15, bound 1e-9 --- \emph{rounding of one revolution's CRN-filtered solve; fitCalibration's computed bound is \textasciitilde{}1e-12 on longer arcs}
\item perRevolution noise-free bias error, relative: measured 1.0059e-14, bound 1e-6 --- \emph{the bias is the small parameter beside a 2e-7 signal; same rounding, relative to 2e-7}
\end{itemize}
\item[attributeRevolutionsPrecision] {\small\emph{precision \textperiodcentered\ \texttt{app.\allowbreak{}attributeRevolutions}}} --- v3 WP-17's cases over seeded realisations: a scale driven by F10.7 is "density" with the F10.7 coefficient within 20 \%; a pure secular drift is "sensor". A null driver sits beside both, which is what the 3-sigma default is for.
\begin{itemize}[leftmargin=1.2em,itemsep=0pt,topsep=1pt]\small
\item attributeRevolutions wrong verdicts, fraction: measured 0.0000e+00, bound 0 --- \emph{4 seeds x 2 cases, 60 revolutions, 3-sigma significance}
\item attributeRevolutions F10.7 coefficient, relative error: measured 1.3845e-02, bound 0.2 --- \emph{v3 WP-17: within 20 \%}
\end{itemize}
\item[perRevolutionVectorization] {\small\emph{vectorization \textperiodcentered\ \texttt{app.\allowbreak{}perRevolution}}} --- A window does not see the others: the first windows of a longer arc equal those of the arc cut after them, bitwise.
\item[attributeRevolutionsVectorization] {\small\emph{vectorization \textperiodcentered\ \texttt{app.\allowbreak{}attributeRevolutions}}} --- Two axes attributed together equal each attributed alone, bitwise.
\item[perRevolutionRobustness] {\small\emph{robustness \textperiodcentered\ \texttt{app.\allowbreak{}perRevolution}}} --- A window mostly lost to a gap is "window"; with the window gate lowered, the estimator's own minimum-data gate refuses it and the reason is recorded -- in neither case a number.
\item[attributeRevolutionsRobustness] {\small\emph{robustness \textperiodcentered\ \texttt{app.\allowbreak{}attributeRevolutions}}} --- Collinear drivers are a verdict, not an error; a constant driver is dropped and named; gated revolutions (NaN) are left out.
\item[perRevolutionMetamorphic] {\small\emph{metamorphic \textperiodcentered\ \texttt{app.\allowbreak{}perRevolution}}} --- Doubling the observation doubles every estimated parameter of the free axes, bitwise (the equilibration takes a power of two out exactly).
\item[attributeRevolutionsMetamorphic] {\small\emph{metamorphic \textperiodcentered\ \texttt{app.\allowbreak{}attributeRevolutions}}} --- A driver in other units (x 10) has its coefficient divided by 10 and the verdict unchanged.
\begin{itemize}[leftmargin=1.2em,itemsep=0pt,topsep=1pt]\small
\item attributeRevolutions driver x10 -\textgreater{} coefficient / 10: measured 8.8818e-16, bound 1e-12 --- \emph{a rescaled column in a QR solve; rounding only}
\end{itemize}
\item[perRevolutionSpeed] {\small\emph{speed \textperiodcentered\ \texttt{app.\allowbreak{}perRevolution}}} --- One fit per window: twice the revolutions, about twice the cost.
\begin{itemize}[leftmargin=1.2em,itemsep=0pt,topsep=1pt]\small
\item perRevolution, 2x arc / 1x arc: measured 2.0975e+00, bound 4 --- \emph{linear in windows would be 2}
\end{itemize}
\item[attributeRevolutionsSpeed] {\small\emph{speed \textperiodcentered\ \texttt{app.\allowbreak{}attributeRevolutions}}} --- The bootstrap dominates and is linear in the revolutions: twice the rows, about twice the cost.
\begin{itemize}[leftmargin=1.2em,itemsep=0pt,topsep=1pt]\small
\item attributeRevolutions, 2x revolutions / 1x: measured 6.8648e-01, bound 4 --- \emph{linear would be 2}
\end{itemize}
\end{description}

\subsection{\texttt{tNgfSensor}}
\noindent Tests of the sensor layer (LP-7, WP-29; N-04 CRN default): sensor.Electrostatic, sensor.AtomInterferometer, sensor.DragFree.

\begin{description}[leftmargin=1em,style=nextline,itemsep=2pt,font=\normalfont\ttfamily\small]
\item[ElectrostaticContract] {\small\emph{contract \textperiodcentered\ \texttt{sensor.\allowbreak{}Electrostatic}}}
\item[ElectrostaticReference] {\small\emph{reference \textperiodcentered\ \texttt{sensor.\allowbreak{}Electrostatic}}} --- COUPLING IS KINEMATICS, so it is checked against kinematics and not against the same formula written a second time. A point rigidly attached at dr in the body frame sits at R(t) dr in the inertial frame; its inertial acceleration is the second derivative of that, and resolved back into the body frame it must be exactly what the sensor computes. The rotation is ANALYTIC -- a fixed axis whose rate ramps, so both the centrifugal and the Euler term are non-zero -- which leaves the second difference with only its own O(dt\textasciicircum{}2) truncation instead of the attitude product's noise.
\begin{itemize}[leftmargin=1.2em,itemsep=0pt,topsep=1pt]\small
\item proof-mass coupling vs a twice-differentiated analytic rotation: measured 1.9272e-06, bound 1e-5 --- \emph{a centred second difference truncates at (omega dt)\textasciicircum{}2/12 \textasciitilde{} 3e-7 here; the bound is the oracle's error, not the model's}
\item the Euler term is a real part of the reference case: measured 1.6617e+00, bound 1e6 --- \emph{recorded so the reference case cannot silently degenerate into a pure spin, where omegaDot vanishes and half the model is untested}
\item \textbar{}referenceSystem omega\textbar{} vs the sensor's, same magnitude: measured 0.0000e+00, bound 1e-12
\item proof-mass coupling magnitude on the real GRACE-C arc: measured 5.2696e-10, bound 1e-7 m/s\textasciicircum{}2 --- \emph{a few hundred micrometres at the yaw-steering rate: order 1e-9, above the 1e-10 noise floor and below the drag}
\end{itemize}
\item[ElectrostaticPrecision] {\small\emph{precision \textperiodcentered\ \texttt{sensor.\allowbreak{}Electrostatic}}} --- The partials against central differences of observe itself. Every parameter is perturbed in its own natural size: a bias in m/s\textasciicircum{}2, a scale dimensionless, an offset in metres, so one step size cannot be right for all fifteen.
\begin{itemize}[leftmargin=1.2em,itemsep=0pt,topsep=1pt]\small
\item Electrostatic partials vs central differences (15 parameters): measured 1.9630e-12, bound 1e-9 --- \emph{exact linearity plus a step chosen to keep the difference well above eps of the reading: what is left is rounding, not truncation}
\end{itemize}
\item[ElectrostaticVectorization] {\small\emph{vectorization \textperiodcentered\ \texttt{sensor.\allowbreak{}Electrostatic}}} --- K epochs at once and K epochs one at a time are the same numbers, EXCEPT where the filter is on: a filtered sample depends on its neighbours by construction, so a single epoch cannot reproduce it and the mask says so. Both halves are asserted.
\begin{itemize}[leftmargin=1.2em,itemsep=0pt,topsep=1pt]\small
\item Electrostatic, one epoch at a time vs the whole arc: measured 0.0000e+00, bound 0 m/s\textasciicircum{}2 --- \emph{without a filter every sample is independent, so the batch may not change a single one}
\item filtered partials vs the filtered observation, bias column: measured 0.0000e+00, bound 1e-16 --- \emph{the bias enters every sample as 1, so its filtered column is the kernel's step response}
\end{itemize}
\item[ElectrostaticRobustness] {\small\emph{robustness \textperiodcentered\ \texttt{sensor.\allowbreak{}Electrostatic}}}
\begin{itemize}[leftmargin=1.2em,itemsep=0pt,topsep=1pt]\small
\item zero proof-mass offset removes the coupling: measured 0.0000e+00, bound 0 m/s\textasciicircum{}2 --- \emph{with dr = 0 both cross products vanish identically, not approximately}
\item a proof mass on the spin axis feels nothing: measured 0.0000e+00, bound 0 m/s\textasciicircum{}2 --- \emph{omega x dr and omegaDot x dr are both zero when dr is parallel to the axis}
\item the gradient term equals -T dr exactly: measured 2.9985e-23, bound 1e-21 m/s\textasciicircum{}2 --- \emph{one matrix-vector product, subtracted where the observation equation subtracts it; eps(1e-7) is 1.3e-23, so this is the cancellation floor with room}
\end{itemize}
\item[ElectrostaticMetamorphic] {\small\emph{metamorphic \textperiodcentered\ \texttt{sensor.\allowbreak{}Electrostatic}}}
\begin{itemize}[leftmargin=1.2em,itemsep=0pt,topsep=1pt]\small
\item a bias change is a rigid shift of the whole series: measured 3.3881e-21, bound 1e-20 m/s\textasciicircum{}2 --- \emph{bias enters once, outside everything else; eps(1e-5) is 1.7e-21}
\item doubling the scale doubles the non-bias part: measured 1.3949e-15, bound floorRel --- \emph{the scale multiplies the sensed acceleration and only that; the bound is the cancellation floor of removing a bias three times the signal}
\item the coupling is linear in the proof-mass offset: measured 2.6144e-16, bound 1e-14 --- \emph{both cross-product terms are linear operators acting on dr}
\item the centrifugal term is even in omega: measured 0.0000e+00, bound 1e-21 m/s\textasciicircum{}2 --- \emph{omega x (omega x dr) contains omega twice, so its sign cannot depend on omega's}
\item the Euler term is odd in omegaDot: measured 0.0000e+00, bound 1e-21 m/s\textasciicircum{}2 --- \emph{omegaDot x dr contains omegaDot once}
\end{itemize}
\item[ElectrostaticSpeed] {\small\emph{speed \textperiodcentered\ \texttt{sensor.\allowbreak{}Electrostatic}}} --- The reading is vectorised over epochs and the partials are a loop over them, which is the honest shape for a 15-column Jacobian. What must hold is that neither is quadratic: ten times the epochs must not cost a hundred times as much.
\begin{itemize}[leftmargin=1.2em,itemsep=0pt,topsep=1pt]\small
\item Electrostatic observe, 60 epochs / 6 epochs: measured 1.8213e+00, bound 10 --- \emph{linear in K would be 10 and a fixed overhead pulls it below; above 10 would mean the work is not vectorised}
\item Electrostatic partials, 60 epochs / 6 epochs: measured 3.1345e+00, bound 20 --- \emph{a per-epoch loop of fixed work is linear; the allowance is generous because the loop body is tiny and the timer is not}
\end{itemize}
\item[ElectrostaticCrnContract] {\small\emph{contract \textperiodcentered\ \texttt{sensor.\allowbreak{}Electrostatic}}}
\item[ElectrostaticCrnReference] {\small\emph{reference \textperiodcentered\ \texttt{sensor.\allowbreak{}Electrostatic}}} --- 1. The taps ARE timeSeriesPreparation's, unchanged: one authority (R1).
\begin{itemize}[leftmargin=1.2em,itemsep=0pt,topsep=1pt]\small
\item CRN DC gain through the sensor vs the printed 1 + 1.115e-7: measured 1.5984e-11, bound 5e-11 --- \emph{half a unit in the fourth printed digit of 1.115e-7}
\item CRN at f0: interior output vs input, relative to the amplitude: measured 3.1764e-15, bound bound --- \emph{normalised to unit response at f0 and zero phase by symmetry; the bound is 10x an Nf-term dot product's rounding}
\end{itemize}
\item[ElectrostaticCrnPrecision] {\small\emph{precision \textperiodcentered\ \texttt{sensor.\allowbreak{}Electrostatic}}} --- The partials go through the same kernel as the reading. The bias enters every sample as 1, so its filtered column is the kernel's sum in the interior; a scale column is the filtered input it multiplies. Bounds computed from the dot product's rounding.
\begin{itemize}[leftmargin=1.2em,itemsep=0pt,topsep=1pt]\small
\item filtered bias column vs the kernel's sum: measured 6.6613e-16, bound bound --- \emph{the bias column is the step response, flat in the interior}
\item filtered scale column vs the filtered input: measured 0.0000e+00, bound bound --- \emph{the same convolution, so equal to rounding}
\item J*theta vs the filtered reading: measured 3.6993e-15, bound 1e3 * bound --- \emph{fifteen columns summed; cancellation against a 3e-6 bias is the floor, hence the larger factor}
\end{itemize}
\item[ElectrostaticCrnRobustness] {\small\emph{robustness \textperiodcentered\ \texttt{sensor.\allowbreak{}Electrostatic}}}
\item[ElectrostaticCrnMetamorphic] {\small\emph{metamorphic \textperiodcentered\ \texttt{sensor.\allowbreak{}Electrostatic}}} --- SAMPLED, NEVER AVERAGED: 'OutputRate' 1 is the full-rate result at every tenth sample, bit for bit, for the reading, its mask, its epochs and its partials (L1B forms ACT1B so).
\begin{itemize}[leftmargin=1.2em,itemsep=0pt,topsep=1pt]\small
\item the CRN-filtered reading is linear in the acceleration: measured 4.9792e-15, bound bound --- \emph{one convolution, linear by construction; the bound is computed: three Nf-term dot products, ten times their rounding each}
\end{itemize}
\item[AtomInterferometerContract] {\small\emph{contract \textperiodcentered\ \texttt{sensor.\allowbreak{}AtomInterferometer}}}
\item[AtomInterferometerReference] {\small\emph{reference \textperiodcentered\ \texttt{sensor.\allowbreak{}AtomInterferometer}}} --- IDENTITY IS THE ORACLE: for a CONSTANT acceleration the triangular weighting has unit sum, so the phase must be exactly k\_eff a T\textasciicircum{}2 -- the textbook result, and the thing that would break first if the kernel were normalised wrongly.
\begin{itemize}[leftmargin=1.2em,itemsep=0pt,topsep=1pt]\small
\item constant acceleration gives exactly k\_eff a T\textasciicircum{}2: measured 1.7239e-16, bound 1e-15 --- \emph{the kernel is normalised to unit sum, so a constant passes through untouched}
\item the Coriolis phase matches -2 k\_eff T\textasciicircum{}2 \textless{}n, omega x v\textgreater{}: measured 2.0060e-16, bound 1e-15 --- \emph{an independent evaluation of the same closed form on constant inputs}
\item Coriolis phase for a GRACE-like rate and 1 mm/s atom velocity: measured 8.8550e+00, bound 100 rad --- \emph{2 k\_eff T\textasciicircum{}2 \textbar{}omega x v\_at\textbar{} = 8.86 rad = 1.41 fringes; ceiling 100}
\item Coriolis phase as an equivalent acceleration: measured 7.3333e+00, bound 100 --- \emph{in units of a 3e-7 m/s\textasciicircum{}2 drag: 7.3, so the rotation term is not a correction to the measurement but a term larger than it}
\end{itemize}
\item[AtomInterferometerPrecision] {\small\emph{precision \textperiodcentered\ \texttt{sensor.\allowbreak{}AtomInterferometer}}} --- The two parameters enter linearly, so central differences with a LARGE step are exact and the only error is cancellation -- the lesson from the Electrostatic partials.
\begin{itemize}[leftmargin=1.2em,itemsep=0pt,topsep=1pt]\small
\item AtomInterferometer partials vs central differences: measured 7.4584e-14, bound 1e-10 --- \emph{both parameters are exactly linear, so a large step leaves only rounding}
\item wrapped phase plus 2 pi times the fringe reconstructs the phase: measured 4.4409e-16, bound 1e-12 rad --- \emph{the split is exact by construction; a caller resolving the ambiguity must get back what was modelled}
\end{itemize}
\item[AtomInterferometerVectorization] {\small\emph{vectorization \textperiodcentered\ \texttt{sensor.\allowbreak{}AtomInterferometer}}} --- KERNEL HAS FINITE SUPPORT, so an interior sample cannot depend on data outside its own window. That is an exact statement -- perturb a distant sample and the interior phase may not move at all -- and it is stronger than comparing two arcs, which would only measure how the edges are handled.
\begin{itemize}[leftmargin=1.2em,itemsep=0pt,topsep=1pt]\small
\item a sample outside the window cannot move the phase: measured 0.0000e+00, bound 1e-14 --- \emph{the triangular kernel is supported on 2T+1 samples and nothing else enters}
\end{itemize}
\item[AtomInterferometerRobustness] {\small\emph{robustness \textperiodcentered\ \texttt{sensor.\allowbreak{}AtomInterferometer}}}
\begin{itemize}[leftmargin=1.2em,itemsep=0pt,topsep=1pt]\small
\item an unresolved window leaves the acceleration unweighted: measured 0.0000e+00, bound 0 m/s\textasciicircum{}2 --- \emph{T * fs below one sample: the window IS a sample, and pretending otherwise would invent a kernel}
\item velocity along the rotation axis gives no Coriolis phase: measured 0.0000e+00, bound 0 rad --- \emph{omega x v vanishes identically when the two are parallel}
\end{itemize}
\item[AtomInterferometerMetamorphic] {\small\emph{metamorphic \textperiodcentered\ \texttt{sensor.\allowbreak{}AtomInterferometer}}}
\begin{itemize}[leftmargin=1.2em,itemsep=0pt,topsep=1pt]\small
\item the phase doubles with the acceleration: measured 0.0000e+00, bound 1e-15 --- \emph{k\_eff T\textasciicircum{}2 multiplies a linearly and nothing else touches it}
\item doubling the interrogation time quadruples the phase: measured 0.0000e+00, bound 1e-15 --- \emph{T enters squared; for a constant acceleration the weighting is exact at both T}
\item reversing the beam reverses the phase: measured 0.0000e+00, bound 1e-15 --- \emph{the projection \textless{}n, a\textgreater{} is the only place the direction enters}
\item the Coriolis phase is odd in omega: measured 0.0000e+00, bound 1e-15 --- \emph{omega x v contains omega once}
\item a phase bias is a rigid shift: measured 2.2204e-16, bound 1e-12 rad --- \emph{it enters outside the gain and outside the weighting}
\end{itemize}
\item[AtomInterferometerSpeed] {\small\emph{speed \textperiodcentered\ \texttt{sensor.\allowbreak{}AtomInterferometer}}} --- One convolution and a handful of vector operations: linear in K.
\begin{itemize}[leftmargin=1.2em,itemsep=0pt,topsep=1pt]\small
\item AtomInterferometer observe, 2000 epochs / 200 epochs: measured 2.3126e+00, bound 10 --- \emph{linear in K would be 10 and a fixed overhead pulls it below; the convolution is the only term that grows}
\end{itemize}
\item[DragFreeContract] {\small\emph{contract \textperiodcentered\ \texttt{sensor.\allowbreak{}DragFree}}}
\item[DragFreeReference] {\small\emph{reference \textperiodcentered\ \texttt{sensor.\allowbreak{}DragFree}}} --- ORACLE IS THE FREQUENCY RESPONSE, derived on the unit circle from the same difference-equation coefficients but evaluated analytically -- a different route through the same definition, which is what a reference for a filter can be. A steady-state sinusoid in must come out scaled by exactly \textbar{}S(f)\textbar{}.
\begin{itemize}[leftmargin=1.2em,itemsep=0pt,topsep=1pt]\small
\item DragFree rejection vs the analytic \textbar{}S(f)\textbar{}, three decades: measured 6.2793e-16, bound 1e-6 --- \emph{an integer number of periods after 60 time constants; what is left is the residue of the transient, not a modelling difference}
\item DragFree gain at the bandwidth is 1/sqrt(2) for one section: measured 1.1633e-04, bound 1e-3 --- \emph{the bilinear transform warps the corner slightly; 1e-3 is that warping at f\_c = fs/100, not a slack bound}
\item a constant drag is rejected to nothing: measured 0.0000e+00, bound 1e-12 --- \emph{the loop has all the time in the world against a constant; what remains is the decayed transient}
\end{itemize}
\item[DragFreePrecision] {\small\emph{precision \textperiodcentered\ \texttt{sensor.\allowbreak{}DragFree}}} --- Both parameters are linear, so a large step makes the central difference exact.
\begin{itemize}[leftmargin=1.2em,itemsep=0pt,topsep=1pt]\small
\item DragFree partials vs central differences: measured 5.5905e-15, bound 1e-9 --- \emph{bias and gain both enter linearly; the step is large enough that cancellation is not the limit}
\end{itemize}
\item[DragFreeVectorization] {\small\emph{vectorization \textperiodcentered\ \texttt{sensor.\allowbreak{}DragFree}}} --- The three axes go through the same filter but share no state. Perturbing one may not touch another, exactly -- a shared accumulator would show here and nowhere else.
\begin{itemize}[leftmargin=1.2em,itemsep=0pt,topsep=1pt]\small
\item perturbing the y axis leaves x and z untouched: measured 0.0000e+00, bound 0 m/s\textasciicircum{}2 --- \emph{one filter state per axis; a shared accumulator would leak here}
\end{itemize}
\item[DragFreeRobustness] {\small\emph{robustness \textperiodcentered\ \texttt{sensor.\allowbreak{}DragFree}}}
\begin{itemize}[leftmargin=1.2em,itemsep=0pt,topsep=1pt]\small
\item the saturated residual is the drag minus the authority: measured 0.0000e+00, bound 1e-9 --- \emph{the command is clipped in its own direction and the excess is what reaches the payload}
\item an ideal loop rejects the same constant to nothing: measured 0.0000e+00, bound 1e-12 --- \emph{with infinite authority the constant is entirely inside the loop's reach}
\item zero drag gives exactly zero residual: measured 0.0000e+00, bound 0 m/s\textasciicircum{}2 --- \emph{no bias, no input, no state: nothing to produce}
\end{itemize}
\item[DragFreeMetamorphic] {\small\emph{metamorphic \textperiodcentered\ \texttt{sensor.\allowbreak{}DragFree}}}
\begin{itemize}[leftmargin=1.2em,itemsep=0pt,topsep=1pt]\small
\item the residual doubles with the drag below saturation: measured 0.0000e+00, bound 1e-14 --- \emph{an unsaturated loop is a linear filter and nothing else}
\item reversing the drag reverses the residual: measured 0.0000e+00, bound 1e-14 --- \emph{linearity again, in the sign}
\item a bias is a rigid shift of the residual: measured 4.5495e-24, bound biasFloor m/s\textasciicircum{}2 --- \emph{it is added outside the filter and outside the gain; the bound is ten times eps of the residual it is differenced out of}
\item a wider loop rejects a low frequency harder: measured 1.3315e-01, bound 1 --- \emph{eight times the bandwidth against a 0.02 Hz signal: the ratio must be below one, and how far below is the design trade}
\end{itemize}
\item[DragFreeSpeed] {\small\emph{speed \textperiodcentered\ \texttt{sensor.\allowbreak{}DragFree}}} --- The filter is a per-sample recursion, so it is linear in K by construction. What is measured is that nothing quadratic crept into the surrounding bookkeeping.
\begin{itemize}[leftmargin=1.2em,itemsep=0pt,topsep=1pt]\small
\item DragFree observe, 5000 epochs / 500 epochs: measured 7.2289e+00, bound 20 --- \emph{linear in K would be 10; the allowance covers an interpreted loop whose per-sample cost is small next to the timer}
\end{itemize}
\end{description}

\subsection{\texttt{tNgfStubs}}
\noindent Every stub in doc/functionlist.csv raises NonGravForces:\textless{}name\textgreater{}:notImplemented.

\begin{description}[leftmargin=1em,style=nextline,itemsep=2pt,font=\normalfont\ttfamily\small]
\item[stubRaises] {\small\emph{contract \textperiodcentered\ no unit}}
\end{description}

\subsection{\texttt{tNgfTools}}
\noindent The development instruments in tools/ run to completion under the gate.

\begin{description}[leftmargin=1em,style=nextline,itemsep=2pt,font=\normalfont\ttfamily\small]
\item[benchmarkRunsAndKeepsItsShape] {\small\emph{contract \textperiodcentered\ no unit}} --- The smallest run that still exercises every path: one batch size, one timed repetition. The cost is dominated by reading the arc and building the two Models, not by the timing loop, so asking for more would buy no coverage.
\item[benchmarkRefusesAnArcItDoesNotHave] {\small\emph{contract \textperiodcentered\ no unit}} --- The committed GNV1B excerpt is an excerpt, so a large enough batch has no states to time. The count is deliberately not written here -- benchmark.m reads it from the file and puts it in the message; 100000 is past any excerpt this fixture will hold. Until 2026-09-18 this reached ep.day(1:n) and raised MATLAB's "Index exceeds the number of array elements" -- a message naming neither the fixture nor the limit, from no unit the caller can look up. Paired with the call that must SUCCEED, because a refusal test whose sibling does not run proves only that everything refuses.
\end{description}
